% Options for packages loaded elsewhere
% Options for packages loaded elsewhere
\PassOptionsToPackage{unicode}{hyperref}
\PassOptionsToPackage{hyphens}{url}
%
\documentclass[
  letterpaper,
]{scrbook}
\usepackage{xcolor}
\usepackage{amsmath,amssymb}
\setcounter{secnumdepth}{5}
\usepackage{iftex}
\ifPDFTeX
  \usepackage[T1]{fontenc}
  \usepackage[utf8]{inputenc}
  \usepackage{textcomp} % provide euro and other symbols
\else % if luatex or xetex
  \usepackage{unicode-math} % this also loads fontspec
  \defaultfontfeatures{Scale=MatchLowercase}
  \defaultfontfeatures[\rmfamily]{Ligatures=TeX,Scale=1}
\fi
\usepackage{lmodern}
\ifPDFTeX\else
  % xetex/luatex font selection
\fi
% Use upquote if available, for straight quotes in verbatim environments
\IfFileExists{upquote.sty}{\usepackage{upquote}}{}
\IfFileExists{microtype.sty}{% use microtype if available
  \usepackage[]{microtype}
  \UseMicrotypeSet[protrusion]{basicmath} % disable protrusion for tt fonts
}{}
\makeatletter
\@ifundefined{KOMAClassName}{% if non-KOMA class
  \IfFileExists{parskip.sty}{%
    \usepackage{parskip}
  }{% else
    \setlength{\parindent}{0pt}
    \setlength{\parskip}{6pt plus 2pt minus 1pt}}
}{% if KOMA class
  \KOMAoptions{parskip=half}}
\makeatother
% Make \paragraph and \subparagraph free-standing
\makeatletter
\ifx\paragraph\undefined\else
  \let\oldparagraph\paragraph
  \renewcommand{\paragraph}{
    \@ifstar
      \xxxParagraphStar
      \xxxParagraphNoStar
  }
  \newcommand{\xxxParagraphStar}[1]{\oldparagraph*{#1}\mbox{}}
  \newcommand{\xxxParagraphNoStar}[1]{\oldparagraph{#1}\mbox{}}
\fi
\ifx\subparagraph\undefined\else
  \let\oldsubparagraph\subparagraph
  \renewcommand{\subparagraph}{
    \@ifstar
      \xxxSubParagraphStar
      \xxxSubParagraphNoStar
  }
  \newcommand{\xxxSubParagraphStar}[1]{\oldsubparagraph*{#1}\mbox{}}
  \newcommand{\xxxSubParagraphNoStar}[1]{\oldsubparagraph{#1}\mbox{}}
\fi
\makeatother


\usepackage{longtable,booktabs,array}
\usepackage{calc} % for calculating minipage widths
% Correct order of tables after \paragraph or \subparagraph
\usepackage{etoolbox}
\makeatletter
\patchcmd\longtable{\par}{\if@noskipsec\mbox{}\fi\par}{}{}
\makeatother
% Allow footnotes in longtable head/foot
\IfFileExists{footnotehyper.sty}{\usepackage{footnotehyper}}{\usepackage{footnote}}
\makesavenoteenv{longtable}
\usepackage{graphicx}
\makeatletter
\newsavebox\pandoc@box
\newcommand*\pandocbounded[1]{% scales image to fit in text height/width
  \sbox\pandoc@box{#1}%
  \Gscale@div\@tempa{\textheight}{\dimexpr\ht\pandoc@box+\dp\pandoc@box\relax}%
  \Gscale@div\@tempb{\linewidth}{\wd\pandoc@box}%
  \ifdim\@tempb\p@<\@tempa\p@\let\@tempa\@tempb\fi% select the smaller of both
  \ifdim\@tempa\p@<\p@\scalebox{\@tempa}{\usebox\pandoc@box}%
  \else\usebox{\pandoc@box}%
  \fi%
}
% Set default figure placement to htbp
\def\fps@figure{htbp}
\makeatother

\ifLuaTeX
  \usepackage{luacolor}
  \usepackage[soul]{lua-ul}
\else
  \usepackage{soul}
\fi

% definitions for citeproc citations
\NewDocumentCommand\citeproctext{}{}
\NewDocumentCommand\citeproc{mm}{%
  \begingroup\def\citeproctext{#2}\cite{#1}\endgroup}
\makeatletter
 % allow citations to break across lines
 \let\@cite@ofmt\@firstofone
 % avoid brackets around text for \cite:
 \def\@biblabel#1{}
 \def\@cite#1#2{{#1\if@tempswa , #2\fi}}
\makeatother
\newlength{\cslhangindent}
\setlength{\cslhangindent}{1.5em}
\newlength{\csllabelwidth}
\setlength{\csllabelwidth}{3em}
\newenvironment{CSLReferences}[2] % #1 hanging-indent, #2 entry-spacing
 {\begin{list}{}{%
  \setlength{\itemindent}{0pt}
  \setlength{\leftmargin}{0pt}
  \setlength{\parsep}{0pt}
  % turn on hanging indent if param 1 is 1
  \ifodd #1
   \setlength{\leftmargin}{\cslhangindent}
   \setlength{\itemindent}{-1\cslhangindent}
  \fi
  % set entry spacing
  \setlength{\itemsep}{#2\baselineskip}}}
 {\end{list}}
\usepackage{calc}
\newcommand{\CSLBlock}[1]{\hfill\break\parbox[t]{\linewidth}{\strut\ignorespaces#1\strut}}
\newcommand{\CSLLeftMargin}[1]{\parbox[t]{\csllabelwidth}{\strut#1\strut}}
\newcommand{\CSLRightInline}[1]{\parbox[t]{\linewidth - \csllabelwidth}{\strut#1\strut}}
\newcommand{\CSLIndent}[1]{\hspace{\cslhangindent}#1}



\setlength{\emergencystretch}{3em} % prevent overfull lines

\providecommand{\tightlist}{%
  \setlength{\itemsep}{0pt}\setlength{\parskip}{0pt}}



 


\makeatletter
\@ifpackageloaded{bookmark}{}{\usepackage{bookmark}}
\makeatother
\makeatletter
\@ifpackageloaded{caption}{}{\usepackage{caption}}
\AtBeginDocument{%
\ifdefined\contentsname
  \renewcommand*\contentsname{Table of contents}
\else
  \newcommand\contentsname{Table of contents}
\fi
\ifdefined\listfigurename
  \renewcommand*\listfigurename{List of Figures}
\else
  \newcommand\listfigurename{List of Figures}
\fi
\ifdefined\listtablename
  \renewcommand*\listtablename{List of Tables}
\else
  \newcommand\listtablename{List of Tables}
\fi
\ifdefined\figurename
  \renewcommand*\figurename{Figure}
\else
  \newcommand\figurename{Figure}
\fi
\ifdefined\tablename
  \renewcommand*\tablename{Table}
\else
  \newcommand\tablename{Table}
\fi
}
\@ifpackageloaded{float}{}{\usepackage{float}}
\floatstyle{ruled}
\@ifundefined{c@chapter}{\newfloat{codelisting}{h}{lop}}{\newfloat{codelisting}{h}{lop}[chapter]}
\floatname{codelisting}{Listing}
\newcommand*\listoflistings{\listof{codelisting}{List of Listings}}
\makeatother
\makeatletter
\makeatother
\makeatletter
\@ifpackageloaded{caption}{}{\usepackage{caption}}
\@ifpackageloaded{subcaption}{}{\usepackage{subcaption}}
\makeatother
\usepackage{bookmark}
\IfFileExists{xurl.sty}{\usepackage{xurl}}{} % add URL line breaks if available
\urlstyle{same}
\hypersetup{
  pdftitle={CMC GI Surgical Oncology Handbook},
  pdfauthor={Jonathan Salo MD},
  hidelinks,
  pdfcreator={LaTeX via pandoc}}


\title{CMC GI Surgical Oncology Handbook}
\author{Jonathan Salo MD}
\date{}
\begin{document}
\frontmatter
\maketitle

\renewcommand*\contentsname{Table of contents}
{
\setcounter{tocdepth}{2}
\tableofcontents
}

\mainmatter
\bookmarksetup{startatroot}

\chapter*{Introduction}\label{introduction}
\addcontentsline{toc}{chapter}{Introduction}

\markboth{Introduction}{Introduction}

This handbook is designed for education and guidance of surgical
trainees at Carolinas Medical Center.

Please forward corrections or suggestions to
Jonathan.Salo@atriumhealth.org

\part{Inpatient Care}

\chapter{CMC Inpatient}\label{cmc-inpatient}

Colorectal Surgery (Davis/Kasten) and GI Surgical Oncology
(Hill/Salo/Squires) will cover Pineville and CMC. For efficiency, the
services at each hospital will merge for patient care. Each patient will
continue to have an attending surgeon, but rounding and inpatient care
will be provided by the service.

\section{Admissions}\label{admissions}

Provider group: ``GI Surg Onc Attending LCI CMC''

List Attending Surgeon in addition

Patient List is CMC GI Surgical Oncology

\section{Rounds}\label{rounds}

Attending rounds 4-5pm on Tuesday and 4-5pm on Thursday

\section{Resident Epic teams:}\label{resident-epic-teams}

\begin{itemize}
\tightlist
\item
  GI Surgical Oncology Colorectal LCI CMC
\item
  Colorectal Surgery Pineville
\end{itemize}

It is critical that you notify service attendings before the start of
the month to adjust the resident call schedule. Each ``shift'' is 5:50am
to 6pm. At 6pm the resident Epic Teams will be forwarded to the night
team.

Please append a text block to the bottom of each progress note
specifying the Epic Team \emph{GI Surgical Oncology Colorectal LCI CMC}
for that patient to facilitate communication from nursing.

\section{Consults}\label{consults}

Established patients and directed should be discussed with the attending
surgeon.

All new inpatient, PCL, and ER consults will be evaluated by the Red
Team instead of the surgical oncology team. Established patients will be
directed to either the Red Team or surgical oncology team depending on
the time since surgery. Those within the three-month global period after
surgery will continue to be evaluated by the surgical oncology team,
while patients outside of the 90-day global period will be managed by
the Red Team. This change aims to streamline patient care. All
established patients will be listed on the surgical oncology list,
though residents will not be responsible for management if the Red Team
is primary.

\section{Postop Clinic Appt}\label{postop-clinic-appt}

Postoperative patients are generally seen for a Transition of Care visit
within the first week

Discharge appointments are made by sending a message in Canopy the
evening prior (preferred) OR the morning of discharge before 8am to:

\begin{itemize}
\tightlist
\item
  LCI CMC GI, Clerical
\item
  Matthew Carpenter (Salo and Squires)
\item
  Rebecca Wicks (Hill)
\item
  Brandon Galloway
\end{itemize}

Please include the following information in the Epic Message:

\begin{itemize}
\tightlist
\item
  Name of attending
\item
  Ward from which the patient is being discharged
\item
  Desired date for appt @ same time
\item
  Need for bloodwork at first visit
\item
  Other studies to be done after discharge

  \begin{itemize}
  \tightlist
  \item
    Upper GI
  \item
    Chest X-ray
  \item
    Modified Barium Swallow
  \end{itemize}
\end{itemize}

Clinic RNs can be reached at: (Hill) 980-442-6146 or (Salo and Squires)
980-442-6143.

For patients likely to go home over the weekend or holidays, please plan
to send a canopy message before 3pm on Friday or the day prior.

The scheduler will respond with a message to the discharging resident
AND to the ward CNL with the appointment time, which can be included
within the discharge summary. Copies of the message will also be sent to
clinical nurse leaders:

\begin{itemize}
\tightlist
\item
  11Tower: Sharon Hood
\end{itemize}

\section{Conferences}\label{conferences}

\begin{itemize}
\tightlist
\item
  GI Tumor Planning Conf Mon 7-8am (via Teams and LCI I 3rd floor Conf
  Rm)
\item
  Resident Teaching Conf Tues 7-8am 5th floor LCI II. Please review the
  upcoming clinic schedule and choose a case to present.
\item
  Bone and Soft Tissue Conf Fri 7-8am (via Teams)
\end{itemize}

\section{Medicine Consults}\label{medicine-consults}

For medicine consults, please use ``CHG Service Hospitalists CMC'' for
all \emph{new} consult requests

\chapter{Pineville Inpatient}\label{pineville-inpatient}

\section{Rounds}\label{rounds-1}

Starting time for rounds is variable from day to day. Maddie Georgino
will help organize work and timing of rounds, etc.

\section{Resident Epic teams:}\label{resident-epic-teams-1}

\begin{itemize}
\tightlist
\item
  Colorectal Surgery Pineville
\end{itemize}

Residents will be assigned to Epic teams by schedule. It is critical
that you notify service attendings before the start of the month to
adjust the resident Epic schedule. Each ``shift'' is 5:50am to 6pm. At
6pm the resident Epic Teams will be forwarded to the night team.

Please append a text block to the bottom of each progress note
specifying the Epic Team for that patient to facilitate communication
from nursing.

\section{Consults}\label{consults-1}

Established patients and directed should be discussed with the attending
surgeon.

\section{Postop Clinic Appts}\label{postop-clinic-appts}

Postoperative patients are generally seen for a Transition of Care visit
at about two weeks.

Discharge appointments are made by sending a message in Canopy the
evening prior (preferred) OR the morning of discharge before 8am to:

\begin{itemize}
\tightlist
\item
  Hale Mock
\item
  Kamisha Wilson
\item
  Madeline Georgino
\end{itemize}

Please include the following information in the Canopy Message:

\begin{itemize}
\tightlist
\item
  Name of attending
\item
  Ward from which the patient is being discharged
\item
  Desired date for appointment
\item
  Need for Wound Ostomy RN appointment at same time (essential for new
  stomas)
\item
  Need for bloodwork at first visit
\item
  Other studies to be done after discharge

  \begin{itemize}
  \tightlist
  \item
    Upper GI
  \item
    Chest X-ray
  \end{itemize}
\end{itemize}

For patients likely to go home over the weekend or holidays, please plan
to send a canopy message before 3pm on Friday or the day prior.

\section{Conferences}\label{conferences-1}

\begin{itemize}
\tightlist
\item
  GI Tumor Planning Conf Monday 7-8am (Teams)
\item
  Resident Teaching Conf 7-8am in Conference Room. Please review the
  upcoming clinic schedule and choose a case to present.
\item
  Bone and Soft Tissue Conference Friday 7-8am (Teams)
\end{itemize}

\chapter{Rounds}\label{rounds-2}

The following format will help speed communication of data on rounds.

\textbf{ID}: One line description: ``Mr Glenn: PostOp day 3 after low
anterior resection''

\textbf{24 hour events}: Summary of important events in prior 24 hrs

\textbf{Data Communication (organized by system)}

Neuro: Pain control, level of alertness, psychotropic meds, sedatives,
and pain meds.

CardioVascular: Vital signs (normal OR cite the range of systolic blood
pressures and range of heart rate). Heart rhythm. Cardiac meds. Most
recent recommendations of cardiology consult.

Respiratory: Pulmonary exam, oxygen saturation, supplied oxygen,
ventilator setting. Results of CXR.

GI: Diet, bowel function, NG output. Drain outputs can often be
summarized unless they are unusually high or low (and ready to be
removed. New finding of bile in any abdominal drain needs special
emphasis. GI meds (eg protonix, Entereg). Tube feed formula, rate and
duration (continuous or nocturnal). Status of C Diff tests. Results of
JP drain amylase levels (gastroesophageal patients). Results of JP
triglycerides or creatinine, if sent,

Renal: Urine output in 24 hours AND in most recent 8 hour shift.
Presence (or absence) of Foley catheter and plans for removal, if
present. Most recent creatinine. If diuretics administered, dosage and
amount of urine output during the shift when it was administered. Most
recent potassium in any patient receiving (or about to receive)
furosemide (Lasix). Results of Mg and Phos if abnormal.

Heme: Hemoglobin, platelets, DVT prophylaxis. PLEASE CHECK THE MAR
SUMMARY DAILY to be certain that the ordered DVT prophylaxis has been
given.

ID: WBC, Tmax in past 24 hours, culture results.

Endo: Diabetic regimen, blood sugar range, and amount of sliding scale
insulin administered in the prior 24 hours.

\textbf{Assessment/Plan (organized by Problem List):}

Each of the patients problems are addressed with an assessment and plan.
Pre-existing medical problems and postoperative complications need to be
addressed in the plan. An assessment and plan for each organ system is
usually not necessary, except for the most complex patients. Patients
active medical problems should be documented on the Patient List for
rounds. This helps to remind the team about medical problems which the
team is managing:

\begin{itemize}
\tightlist
\item
  Chronic anticoagulation
\item
  Diabetes
\item
  Malnutrition
\end{itemize}

This problem list-oriented approach will also be helpful to writing
problem-oriented notes.

\chapter{Progress Notes}\label{progress-notes}

Progress notes need to reflect the problems which are being managed by
the team. The current medical problems being managed by the team should
be added to the problem list for that hospitalization. This will make it
easier to generate notes which are oriented to the patient's problem
list.

In addition, the progress notes provide a narrative which can later be
used to generate the discharge summary (particularly for complex
patients). Each day, the event of the hospitalization is carried from
note to note. Each day, an additional line is added to the progress note
which summarizes events for that day. This makes it possible to see
within each Progress Note the pertinent events for the hospitalization.
These events would include extubation, re-intubation, positive cultures,
dates lines are inserted or removed, dates of removal of NG tubes and
drains, and transfer to ward or re-admission to ICU. This chronology
assists in treatment decisions (``how old is the IJ line'' or ``when did
we start antibiotics?'' or ``when is the planned antibiotic stop
date''?) but also makes the discharge summary much easier to prepare.

\textbf{Notes should be forwarded to the patient's attending (unless out
of town)}

\textbf{Esophagectomy Events to be Documented:}

\begin{itemize}
\tightlist
\item
  Extubation date/time
\item
  NG Removal date
\item
  Chest tube removal date
\item
  MBS date(s) and results (aspiration \textbar{} penetration)
\item
  ICU DC orders written
\item
  ICU discharge (transfer to ward)
\end{itemize}

\textbf{Esophagectomy Complications to be Documented:}

\begin{longtable}[]{@{}
  >{\raggedright\arraybackslash}p{(\linewidth - 2\tabcolsep) * \real{0.0714}}
  >{\raggedright\arraybackslash}p{(\linewidth - 2\tabcolsep) * \real{0.9286}}@{}}
\toprule\noalign{}
\endhead
\bottomrule\noalign{}
\endlastfoot
N & Delirium \\
& Stroke \\
CV & New arrhythmia req Rx \\
& MI \\
R & Pneumonia (3 of fever \textbar{} WBC \textbar{} infiltrate
\textbar{} abx \textbar{} sputum cx) \\
& Effusion req drainage \\
& Reintubation \\
& Atelectasis req bronchoscopy \\
& ARDS \\
& PE \\
& Ventilation \textgreater48 hours after leaving OR \\
GI & Anastomotic leak (medical rx \textbar{} stent \textbar{}
surgery) \\
& Delayed gastric emptying req botox or NG \textgreater7d \\
& C Diff \\
GU & Urinary Retention \\
& Discharge with foley catheter \\
H & DVT req treatment \\
& Return to OR \\
& Return to ICU \\
\end{longtable}

\textbf{Communication}

Please add an addendum at the BOTTOM of each progress note which
includes a means for contacting the team:

Please message ``GI Surgical Oncology Colorectal LCI CMC via Haiku 24/7.
Messages are automatically forwarded to the General Surgery Resident on
Call evenings and weekends.''

\chapter{Consults}\label{consults-2}

\textbf{Common Haiku Groups for Consults}

\begin{longtable}[]{@{}
  >{\raggedright\arraybackslash}p{(\linewidth - 2\tabcolsep) * \real{0.5000}}
  >{\raggedright\arraybackslash}p{(\linewidth - 2\tabcolsep) * \real{0.5000}}@{}}
\toprule\noalign{}
\begin{minipage}[b]{\linewidth}\raggedright
Service
\end{minipage} & \begin{minipage}[b]{\linewidth}\raggedright
Haiku Group
\end{minipage} \\
\midrule\noalign{}
\endhead
\bottomrule\noalign{}
\endlastfoot
Advanced GI (Dries/Doshi/Lewis/Chauhan) & CMC Advanced GI ERCP
Consult \\
Advanced GI (Narang) & CDHA Advanced GI ERCP \\
Hospitalist & CGH Service Hospitalist CMC - Consults Only \\
Interventional Radiology & Eric Dodson - Haiku or 5-3599 \\
Body Radiology & 5-2270 \\
Ultrasound & 5-3600 \\
Endocrine & CMC Endocrinology Consult Team (Kelli Dunn) \\
Nephrology & Metroline Neph Consult CMC \\
Cardiology & CMC SHVI General Cardiology \\
ID & CMC Infectious Disease Consults \\
General Surgery Red & General Surgery Red CMC floor/consults \\
Pulmonary & CMC Pulmonary Consult PCCC \\
STICU & STICU Charge for TRANSFERS \\
Hospital at Home & GCM HAH Nurse Navigators (referrals) CLT \\
& \\
\end{longtable}

\chapter{Ward Phone list}\label{ward-phone-list}

\begin{longtable}[]{@{}ll@{}}
\toprule\noalign{}
\endhead
\bottomrule\noalign{}
\endlastfoot
3B Obs & 704-355-0142 \\
3K Obs & 704-446-6690 \\
3L & 704-355-8644 \\
3T Onc & 704-355-6312 \\
4B BMT & 704-446-9350 \\
5A Overflow & 704-446-5511 \\
6A & 704-355-2033 \\
6B & 705-355-6275 \\
6T Thoracic & 704-355-6260 \\
7B & 704-355-6262 \\
7T & 704-355-6210 \\
7 Dixon Heart & 704-355-2xxx \\
9A & 704-355-6291 \\
9B & 704-355-6261 \\
9T & 704-355-6290 \\
9 NICU & \\
10A & 704-355-6211 \\
10DE PROG & 704-355-2022 \\
10T & 704-355-6230 \\
10 MICU & \\
11A & 704-446-6311 \\
11B & 704-355-6310 \\
11T & 704-446-XXXX \\
11 STICU & 704-355-2023 \\
KDU Dialysis & 704-355-2078 \\
& \\
& \\
\end{longtable}

\chapter{Department Phone list}\label{department-phone-list}

\begin{longtable}[]{@{}ll@{}}
\toprule\noalign{}
\endhead
\bottomrule\noalign{}
\endlastfoot
Admitting & 704-355-3071 \\
Bed Management & 704-512-6360 \\
Code Blue & 704-355-2345 \\
ER Triage & \\
ER Diag & 704-355-3271 \\
ER Major & 704-355-2157 \\
ER AEC & 704-355-8683 \\
ER Holding & 704-355-3276 \\
Endoscopy & 704-355-2792 \\
Lab & 704-355-9350 \\
Language Line & 704-381-8255 \\
Preop & 704-355-8922 \\
PACU & 704-355-2431 \\
Radiology & 704-355-3404 \\
Rad CT & 704-355-2270 \\
Rad MRI & 704-355-4196 \\
Rad Ultrasound & 704-355-3600 \\
Rad IR & 704-355-3430 \\
Wound Ostomy & 704-355-7605 \\
& \\
& \\
& \\
& \\
& \\
& \\
& \\
& \\
& \\
\end{longtable}

\chapter{Signout}\label{signout}

\textbf{Evening Signout}

The Handoff Tool should be completed for all inpatients, and the
responsible attending designated. This tool is critical for the safe
care of patients by the nigh team. If there are studies which are
pending at the time of signout (CT scan, follow-up Hb), it is critical
that a plan be in place for whom to notify with an abnormal or critical
study. In general, Drs Hill, Squires, and Salo are always available
until 10pm. Attending notification plans (service attending vs covering
attending) for unstable patients should be negotiated before nightfall.

\textbf{Weekend Signout}

The senior resident is responsible for making certain that the weekend
rounding resident is familiar with the patients, their problems, and the
plan of care. A signout email should be prepared Friday afternoon and
forwarded to the service attendings by 6pm for their review. This
signout can then be edited with the attendings' notes and forwarded to
the weekend rounding attending.

\chapter{Discharges}\label{discharges}

\textbf{Discharge Prescriptions}

Prescriptions should be ideally be prepared the day prior to anticipated
discharge and sent to the patient's pharmacy. According to
\href{https://www.ncmedboard.org/images/uploads/license_applications/STOPAct-focusedFAQs-Sept2019.pdf}{North
Carolina STOP guidelines}, opioid prescriptions for postoperative
patients should be for no more than a 7 day supply. At the same time,
patients taking narcotics in the hospital should have the same dosages
for their outpatient prescription, to avoid patients running out of
narcotics between the time of discharge and their first clinic visit.

Note that metoprolol is available in liquid form in the hospital, but is
not available for outpatient prescription

Note that in many cases narcotics for patients living in South Carolina
a prescription from an attending.

\textbf{Additional Appointments}

If followup appointments in additional to surgical followup are needed,
these should be designated on the discharge orders. Particularly:

\begin{itemize}
\tightlist
\item
  Primary Care Physician
\item
  Cardiologist (if new cardiac medicines)
\item
  Co-surgeons (Urology, Thoracic Surgery, GYN)
\end{itemize}

\textbf{Discharge Summary}

The discharge summary documents important events and complications in
the postoperative course and serves to inform the referring physician
and primary physician about these events, but also serves as a blueprint
for post-discharge treatment planning. Please recognize that the first
post-operative visit may be with a resident who may be meeting the
patient for the first time. Key items to include:

\begin{longtable}[]{@{}
  >{\raggedright\arraybackslash}p{(\linewidth - 2\tabcolsep) * \real{0.5000}}
  >{\raggedright\arraybackslash}p{(\linewidth - 2\tabcolsep) * \real{0.5000}}@{}}
\toprule\noalign{}
\endhead
\bottomrule\noalign{}
\endlastfoot
N & Followup plan for chronic pain management \\
& Stroke \\
CV & Postop Arrhythmia? \textbar{} MI? \textbar{} CHF? \\
& If new cardiac meds: Who is managing medications \\
& If afib: CHADS score and anticoagulation plan \\
R & Pneumonia? \textbar{} ARDS? \textbar{} TRACH? \\
& Need for home oxygen? \\
& CXR needed at first postop visit? \\
GI & Delayed gastric emptying? \textbar{} leak? \textbar{} ileus? \\
& Tube feed regimen \\
& Diet at discharge (Low residue \textbar{} Full liquds \textbar{} Meds
with thickened water \textbar NPO) \\
& New stoma (ileostomy \textbar{} colostomy) \\
& Wound care needs (VAc \textbar{} Prevena) \\
GU & Urinary Retention \\
& Discharge with foley catheter \\
H & Complications: DVT \textbar{} PE \\
& Anticoagulation Plan \\
Endo & Insulin regimen at DC (dose will be in med rec) \\
ID & Antibiotics at DC \\
& Return to ICU \\
\end{longtable}

\textbf{Communication}

It is essential that discharge summaries be sent to the patient's
primary MD and referring physician. Please review the initial
consultation note for the names of providers involved in a patient's
care.

\chapter{Education}\label{education}

The service will host Third- and Fourth-year medical students from Wake
Forest University as well as externs.

Student notes should be forwarded to the patient's attending for
attestation and signature.

Third year students will have a `green card' of diagnoses and procedures
which need to be checked off (and signed) during the rotation. Students:
Please remind the chief resident and attendings about items which remain
to be completed..

\section{Medical Student Duty Hours}\label{medical-student-duty-hours}

Hours:

\begin{itemize}
\tightlist
\item
  Students will not work longer hours than residents on the same service
\item
  Students will not work more than 80 hours/week averaged over 4 weeks.
\end{itemize}

Breaks: Students will

\begin{itemize}
\tightlist
\item
  have 4 24-hour periods free from assigned activities over a 4 week
  period
\item
  not work longer than 16 continuous hoours
\item
  have a 8 hour break from clinical/academic hours following a 16-hour
  shift
\item
  can only work a maximum of 5 sequential overnight shifts
\end{itemize}

Exams:

\begin{itemize}
\tightlist
\item
  Must be excused no later than midnight prior to the day of the shift
  or final exam
\end{itemize}

Holidays:

\begin{itemize}
\tightlist
\item
  Must be excused from responsibilities from 5pm on the day prior to the
  holiday on the academic calendar through the holiday\footnote{From
    Faculty as Teacher 2023}
\end{itemize}

\section{Procedures/Diseases}\label{proceduresdiseases}

\begin{itemize}
\tightlist
\item
  Wound Care (VAC/dressing change, identify infection)
\item
  Suture/Staple removal
\item
  Suture Skin
\item
  Foley catheter insertion (adult)
\item
  Insert nasogastric tube (or OG in OR)
\item
  Make an incision, any site
\item
  Participate in intubation, bag mask ventilation in OR
\item
  Xray 3-way of abdomen (interpret)
\item
  Xray chest (interpret)
\item
  Assist with insertion of chest tube or pigtail
\end{itemize}

\section{Ask a Resident (5min
discussion)}\label{ask-a-resident-5min-discussion}

\begin{itemize}
\tightlist
\item
  Abdominal Pain (RUQ)
\item
  Acute Limb Ischemia (vascular disease)
\item
  Diverticulitis
\item
  Neoplastic process (Breast Mass, GI Mass, soft tissue)
\item
  Abdominal wall mass or hernia
\item
  ABC's of trauma, Primary/Secondary survey
\item
  Ileus, small and large bowel obstruction
\item
  Evaluate acute surgical abdomen (participate)
\item
  Post-op fever in surgical inpatient (discuss and participate)
\item
  Post-op pain management (discuss and participate)
\end{itemize}

\section{Medical Student Resources}\label{medical-student-resources}

\href{https://www.youtube.com/watch?v=M3vISruyFkI}{Subcuticular suturing
video}

\section[Recommended Resources:]{\texorpdfstring{Recommended
Resources:\footnote{from Spring 2023 Surgery Syllabus}}{Recommended Resources:}}\label{recommended-resources01-students-2}

\begin{itemize}
\tightlist
\item
  Essentials of general surgery, 6th edition, {[}edited by{]} Peter F.
  Lawrence
\item
  Surgery: A Case Based Clinical Review (Christian de Virgilio)
\item
  Kaplan Surgery Notes
\item
  NMS Surgery CaseBook
\item
  Surgical Recall (Lorne H. Blackbourne)
\item
  UWorld QBank
\item
  Aquifer / Wise-MD
\item
  OnlineMedEd
\end{itemize}

\section{SHELF exam}\label{shelf-exam}

\href{https://www.nbme.org/subject-exams/clinical-science/surgery}{SHELF
exam topic areas}

\section{Entrustable Professional
Activities}\label{entrustable-professional-activities}

\begin{enumerate}
\def\labelenumi{\arabic{enumi}.}
\tightlist
\item
  Gather a history and perform a physical examination
\item
  Prioritize a differential diagnosis following a clinical encounter
\item
  Recommend and interpret common diagnostic and screening tests
\item
  Enter and discuss orders and prescriptions
\item
  Document a clinical encounter in the patient record
\item
  Provide an oral presentation of a clinical encounter
\item
  Form clinical questions and retrieve evidence to advance patient care
\item
  Give or receive a patient handover to transition care responsibility
\item
  Collaborate as a member of an interprofessional team
\item
  Recognize a patient requiring urgent or emergent care and initiate
  evaluation and management
\item
  Obtain informed consent for tests and/or procedures
\item
  Perform general procedures of a physician
\item
  Identify system failures and contribute to a culture of safety and
  improvement
\end{enumerate}

\part{Preoperative Care}

\chapter{Preop Esophageal Ca}\label{preop_esophageal}

Patient who present with esophageal and GE junction diagnoses present in
several manners, which determine their evaluation.

\begin{itemize}
\tightlist
\item
  Superficial tumors
\end{itemize}

These patients typically have Barrett's esophagus and are found to have
small tumors on surveillence endoscopy. These are considered
{[}superficial{]}.

Workup typically involved endosopic ultrasound to determine T stage and
endoscopic mucosal resection {[}EMR{]} which can be both diagnostic and
therapeutic for T1 tumors.

\begin{itemize}
\tightlist
\item
  Tumors without dysphagia
\end{itemize}

Patients with minimal symptoms who may present with upper GI bleeding or
vague symptoms but have no difficulty swallowing. Thesea are considered
{[}early stage{]}. Workup typically involves PET and endoscopic
ultrasound (EUS). Some patients will undergo primary surgery.

\begin{itemize}
\tightlist
\item
  Tumors with dysphagia
\end{itemize}

Patients with esophageal or GE junction tumors who present with
dysphagia are considered {[}locally-advanced{]}. Workup involves PET
scan. Most patients will undergo neoadjuvant therapy. Many will need a
central venous port and some will need a feeding tube. Most need a
referral for medical oncology.

\begin{itemize}
\tightlist
\item
  Tumors with metastatic disease
\end{itemize}

Patients who have signs of metastatic disease on CT or PET are
considered {[}metastatic{]}. These patient typically need a central
venous port of chemotherapy.

\section{Medical Evaluation}\label{medical-evaluation}

\section{Restaging}\label{restaging}

\chapter{Salo Clinic}\label{mie_postop}

\section{Clinic Notes}\label{clinic-notes}

In most cases, a note will already have been started, so first check
under the Notes tab.

If note exists, start a new note and use the template ``.lcir''

At the ``HPI\textgreater\textgreater{}'' prompt, enter history of
present illness:

\subsubsection{Identifier:}\label{identifier}

\begin{itemize}
\tightlist
\item
  67M presentation with iron-deficiency anemia. Usual weight 145\#
\item
  73F presentation with 3 months of dysphagia to solids. Usual weight
  181\#. Weight at presentation 151\#
\item
  44M presentation with 6 weeks of reflux symptoms. Usual weight 165\#
\end{itemize}

Baseline (usual) weight is important for gastroesophageal cancer
patients.

No narrative (this is not a creative writing workshop)

Avoid descriptors such as ``pleasant'' or ``unfortunate''

\subsubsection{Workup}\label{workup}

\begin{quote}
EGD April 14, 2024 by Dr Larry Pennington: Malignant appearing stricture
at 40cm. Pathology (S24-345229) shows poorly-differentiated
adenocarcinoma. EUS May 1, 2024 by Dr Andrew Dries. Staged as uT3 N1 PET
May 15, 2024: Hypermetabolic lesion in the distal esophagus. No regional
lymphadenopathy. No metastatic disease
\end{quote}

Each paragraph contains one study: study -\textgreater{} date
-\textgreater{} operator -\textgreater{} findings

Pathology accession number is helpful for staff to track down outside
path

\subsubsection{Consult Visit}\label{consult-visit}

\begin{quote}
GI Surgical Oncology consultation June 1, 2024. Taking only liquids. No
odynophagia. Hand grip strength R: 31/30/29. ECOG 0. Weight 165\#. Still
smoking.
\end{quote}

\subsubsection{ECOG performance scale}\label{ecog-performance-scale}

\begin{itemize}
\tightlist
\item
  0 Fully active, able to carry on all activities
\item
  1 Ambulatory can do light house work, office work
\item
  2 Ambulatory but can't work up and about more than 50\% of waking
  hours
\item
  3 Capable of only limited selfcare; confined to bed or chair more than
  50\% of waking hours
\item
  4 Completely disabled; cannot carry on any selfcare; totally confined
  to bed or chair
\end{itemize}

\subsection{Update Problem list}\label{problem_list}

Critical to update problem list for co-morbidities which would affect
anesthetic/operative management.

Add details to problem list under ``Overview''

\subsection{Cardiac}\label{cardiac}

\begin{itemize}
\item
  History of heart attack?
\item
  Coronary artery stent? When? Within the last 6mo?
\item
  Is patient on antiplatelet therapy? Which agents?
\item
  Date and location of most recent echocardiogram

  \begin{itemize}
  \item
    If aortic stenosis: Is valve area \textless1cm\^{}2?
  \item
    If pulmonary hypertension: Peak pressures?
  \end{itemize}
\item
  Date and location of stress test. Is there ischemia? Who is
  cardiologist?
\end{itemize}

\subsection{Pulmonary}\label{pulmonary}

\begin{itemize}
\tightlist
\item
  Prior diagnosis of COPD? Dyspnea at rest? Dyspnea walking up one
  flight of stairs?\\
  Has patient had recent pulmonary function tests? Is patient on
  steroids? History of prior lung surgery?
\end{itemize}

\subsection{GI}\label{gi}

Prior abdominal surgery? Hernia repair: Prosthetic mesh? Laparoscopic vs
open inguinal hernia repair Prior history of bariatric surgery?

\subsection{Endocrine}\label{endocrine}

If diabetes: Is insulin required? How much? Steroid use?

\subsection{Renal}\label{renal}

If kidney disease, what is baseline creatinine? Dialysis?

\subsection{Neurologic}\label{neurologic}

History of Stroke? History of TIA? Carotid bruit on exam? History of
dementia?\\
Who manages finances? Who makes medical decisions?

\subsection{Physical Exam}\label{physical-exam}

Cardiopulmonary exam

\begin{itemize}
\item
  Carotid Bruit?
\item
  Cervical Incision?
\item
  Cardiac murmur
\item
  Abdominal incisions - reconcile with past surgical history
\end{itemize}

\subsection{Assessment/Plan}\label{assessmentplan}

Enter assessment/plan after ``A/P'' prompt Do not erase blue-shaded
``Assessment Plan: No problem-specific assessment/plan found for this
encounter \emph{Briefly} summarize assessment/plan Dr Salo will add A/P
and rationale/evidence base

\section{Return Visit Notes}\label{return-visit-notes}

Use notes template ``.lcir'' Select HPI text from prior note and use
Copy Forward button

\subsection{Summarize Recent Surgery}\label{summarize-recent-surgery}

\begin{itemize}
\tightlist
\item
  Minimally-invasive esophagectomy July 2, 2024. Benign postoperative
  course. Discharged on 5 cartons of tube feeds and protein shakes.
  Pathology (CMCS24-004553): T3 N1 M0 1/26 nodes margins negative
\item
  Postop visit July 14, 2024: Doing well. Staples out. Advance diet.
  Metoprolol decreased to 25mg daily x 7days-\textgreater{} stop.
\end{itemize}

\textbf{\emph{Sign notes and forward to Dr Salo}}

\section{}\label{section}

\chapter{Hill Clinic}\label{hill-clinic}

\section{Preop Colorectal Patients}\label{preop-colorectal-patients}

\begin{itemize}
\tightlist
\item
  Opiod Cessation
\item
  Smoking Cessation
\item
  EtOH/Drugs of Abuse -- Social Work
\item
  Nutritonal Evaluation

  \begin{itemize}
  \tightlist
  \item
    All Patients -- Ensure for 3d preop
  \item
    Poor nutrition -- ?Delay surgery
  \end{itemize}
\item
  Preop Anesthesia/ERAS Class
\item
  Expectations of Surgery:

  \begin{itemize}
  \tightlist
  \item
    Length of Stay 1-3 days
  \item
    Diet (self-limiting)
  \item
    Pain control (low-opoid)
  \item
    Activity (OOB at 6am, OOB 3x/day)
  \end{itemize}
\item
  Bowel Prep -- Abx and mechanical
\end{itemize}

\part{Operating Room}

\chapter{Colorectal Cases}\label{colorectal-cases}

\section{Preop}\label{preop}

ADULT SURG Colorectal ERAS MPP Hill This is what has all of the main
ERAS components. Tylenol, gabapentin, Decadron, Alvimopan and heparin
are all given in pre-op holding

ADULT STANDING Antimicrobial colorectal In general, I will give Ancef to
patients with almost all patients with allergies to PCN. They have to
remember the ``severe'' reaction. If it is a severe reaction, please use
the second line antibiotics listed in the power plan.

Type and Screen are not typically needed for colectomy. They will have
an antibody screen from office. If antibodies present, then d/w
attending

Entereg for all colorectal patients (if they are not taking preoperative
opoids)

VTE prophylaxix

Carbohydrate load (2 hours preop)

\section{Intraop}\label{intraop}

Positioning: Attending will typically position patients personally

Right sided=supine; Left sided=lithotomy; All laparoscopic colon cases
will have their arms tucked with a chest tape strap.

NG/OG tubes: usually not needed, unless there is gastric distension.
Patients receive multiple oral medicatinos in preop.

Review anesthesia fluid management during time out. 2L max volume, urine
output not an accurate indicator.

Open procedures: No Epidural

\section{Postop}\label{postop}

Fluids: Goal of 3 liters total in first 24 hours

\textbf{Postop Day 0}

\begin{itemize}
\tightlist
\item
  Goal-directed fluid administration
\item
  OOB (Dangling not compliant)
\item
  Diet
\item
  Low Residue diet
\item
  Ensure Supplements
\item
  Teaching

  \begin{itemize}
  \tightlist
  \item
    Gum, Mag \& Entereg
  \end{itemize}
\item
  Pain Management

  \begin{itemize}
  \tightlist
  \item
    PCA
  \item
    Tylenol 1gm q6 ATC
  \item
    Gabapentin 300mg TID
  \item
    Tramadol PRN
  \end{itemize}
\end{itemize}

Resume all baseline pain meds

Resume all home medications

No therapeutic anti-coagulation

Diabetes medicines - Prefer to resume all oral diabetes meds -
Sliding-scale insulin ordered for all diabetics

\textbf{Postop Day 1}

Labs: K+ and CBC only

Heparin lock IV

Remove Foley

d/c PCA (unless open incision)

Out of bed \textgreater{} 6 hrs

Diet: low residue

Pain management

d/c PCA

Tylenol

Gabapentin 300 tid

Tramadol

Home meds

Oxy if lots of pain

Afternoon rounds

Check patient 2-4PM

Ambulation 2x's by PM

Patient education

\textbf{Postop Day 2}

\begin{itemize}
\tightlist
\item
  No IV fluids unless indicated
\item
  No labs unless indicated
\item
  OOB \textgreater6 hrs
\item
  No PT unless going to rehab or SNF (ask Hill first)
\item
  Pain management: same as POD\#1
\end{itemize}

Discharge planning

Consider early D/C (median LOS=2d)

Otherwise plan

Patient

Nursing

Afternoon rounds

Possible home today!

Check patient 2-3PM

Ambulation 2x's by PM

Check for dehydration

Patient education

Give estimated date of discharge

\textbf{Postop Day 3-5}

\begin{itemize}
\tightlist
\item
  IVF

  \begin{itemize}
  \tightlist
  \item
    Consider bolus of IV fluids
  \item
    Consider maintenance IV if we feel won't resolve soon
  \end{itemize}
\item
  Labs

  \begin{itemize}
  \tightlist
  \item
    No labs unless indicated
  \item
    Consider ordering QOD Chem7 if prolonged ileus
  \end{itemize}
\item
  OOB \textgreater6 hrs

  \begin{itemize}
  \tightlist
  \item
    No PT unless likely to go to rehab or SNIF (ask Hill first)
  \end{itemize}
\item
  Pain management: Same as POD\#1
\item
  Afternoon rounds

  \begin{itemize}
  \tightlist
  \item
    Possible home today!
  \item
    Check patient 2-3PM
  \end{itemize}
\end{itemize}

Ambulation 2x's by PM

Check for dehydration

Patient education

Give estimated date of discharge

\chapter{CV Port (IJ)}\label{cv_port_salo}

\textbf{Room Prep}

\begin{itemize}
\tightlist
\item
  Slider bed (Skytron 3600B) with head section
\item
  C-Arm

  \begin{itemize}
  \tightlist
  \item
    Radiology technician alerted to need for C-Arm
  \item
    Will need lead and thyroid shields for everyone in room
  \end{itemize}
\item
  Ultrasound with hockey-stick probe near patient's RIGHT SHOULDER
\end{itemize}

\textbf{Instruments} - Minor instrument pan

\textbf{Disposables/Meds}

\begin{itemize}
\tightlist
\item
  Confirm choice of port with surgeon. Usual options

  \begin{itemize}
  \tightlist
  \item
    Bard PowerPort VUE with 8Fr attachable catheter (1708062)
  \item
    Bard PowerPort slim Implanted port (for patients with low BMI)
  \end{itemize}
\item
  Heparin 5mL of 1000U/ml labeled as ``1000 U/ml''
\item
  Heparin 5mL of 1000U/ml + 45mL saline labeled as ``100 U/ml''
\item
  Local
\item
  If general anesthesia: Marcaine 0.5\% with epinephrine
\item
  If MAC: Xylocains 1\% with epinephrine
\item
  1000 drape x3 AND blue paper drapes 4 packs of 2 each = 8 total
\item
  Suture

  \begin{itemize}
  \tightlist
  \item
    3-0 Prolene RB-1 double-arm
  \item
    3-0 Vicryl SH
  \item
    4-0 Monocryl PS2
  \end{itemize}
\end{itemize}

\textbf{Position}

\begin{itemize}
\tightlist
\item
  Supine with left arm tucked, right arm on armboard at side.

  \begin{itemize}
  \tightlist
  \item
    Right arm on armboard in case needed by anesthetist
  \item
    NO shoulder roll
  \end{itemize}
\item
  Foley catheter: usually NOT required -- \emph{check with surgeon}
\item
  Lower body Bair Hugger from abdomen to feet with ONE layer of blankets
  on top of Bair Hugger. Velcro strap on thighs.
\end{itemize}

\textbf{Prep}

Chloroprep: RIGHT chest, neck to chin and earlobe, shoulder to include
deltopectoral groove.

\textbf{Drape}

1000 plastic drapes outline the sterile field for the port. The skin is
stretched to avoid a gap between drape and skin. Allow access to the
right sternocleidomastoid, right deltopectoral groove, and sternal
notch. Blue paper drapes on top of 1000 drapes Transverse drape reversed
head-to-foot. Ioban around edges of port field. Skin over SCM is left
without Ioban to facilitate ultrasound

\textbf{Preop evaluation}

\begin{itemize}
\tightlist
\item
  Allergies
\item
  Blood thinners or anti-platelet agents
\item
  Hx of prior central lines or ports
\item
  History of neck surgery
\end{itemize}

\textbf{Operation}

\emph{Reverse Trendelenburg}

Port pocket is constructed 1cm below and parallel to clavicle 3cm long.
It is essential that there is no bleeding in the pocket (to avoid a port
pocket hematoma).

\emph{Trendelenburg}

Ultrasound set up so that the lateral neck appears on the right side of
the screen (as if looking towards the feet)

Right internal jugular vein is identified and its course cephalad-caudad
marked on the skin.

Finder needle (22Ga) \emph{OR} micropuncture kit passed into IJ. The
needle should enter the vein directly beneath the ultrasound probe.

Skin anesthetized and transverse 6mm counter-incision made at needle
entry site

\emph{Respiration held by anesthesia}

16Ga needle passed into IJ under ultrasound. It is essential that the
vein is scanned up and down by rolling the the probe to visualize the
tip of the needle as it passes inferior.

J Wire passed through 16Ga needle and needle withdrawn

\emph{Anesthesia resumes respirations}

Ultrasound used to confirm presence of wire within the vein by scanning
up and down.

\emph{Level bed}

Fluoroscopy used to confirm position of wire. Dilator and sheath
inserted under fluoroscopic visualization (`live'). C-arm backed away
off field.

Tunneler connected to tubing on the end with small numbers. Port
collected to tubing on the opposite end. Tubing must come to rest 1mm
from the body of the port. Clear plastic collar is pushed over the
tubing and the edge must be flush with the body of the port. Tunneler
bent into a curve to avoid injury to the carotid artery.

Catheter tunneled from port pocket to counter-incision over SCM. The
tunneler path describes a gentle arc to avoid kinking the catheter. The
port is moved to the top of the port pocket by traction on the catheter.
Tonsil clamp placed at 25cm and fluoroscopy used to evaluate the length
of the catheter, which is trimmed with a heavy scissors.

Catheter placed through dilator into central circulation. Peel-away
sheath split and removed. Port accessed with straight Huber needle with
100U/ml heparinized saline. Blood is withdrawn into port. Needle is left
in the port and the syringe detached.

C-arm is returned to the field to confirm catheter placement. It may be
necessary to `orbit' the C-arm if the catheter overlies the spine.

Syringe with concentrated flush (1000U/ml) is attached to the needle and
the port flushed (without aspiration of blood). Syringe and needle are
removed.

Port sutured to the underlying pectoralis fascia with 2 sutures of 3-0
Prolene, one forehand and one backhand. Sutures are tied and cut.

The port pocket irrigated and the incision closed with two layers of
subcutaneous 3-0 Vicryl followed by subcuticular 4-0 Monocryl. The skin
is dressed with Dermabond.

\textbf{Postop Orders}

CXR in recovery to confirm central line placement

\chapter{Lap Jejunostomy}\label{lap-jejunostomy}

\textbf{Room Prep}

\begin{itemize}
\tightlist
\item
  EGD cart near patient's LEFT SHOULDER (with ADULT EGD scope)
\item
  If central venous port is placed at the same time:
\item
  Slider bed (Skytron 3600B) with head section
\item
  Radiology technician alerted to need for C-Arm
\item
  BK Ultrasound with hockey-stick probe near patient's RIGHT SHOULDER
\end{itemize}

\textbf{Instruments}

\begin{itemize}
\tightlist
\item
  5mm 30 degree scope AND 5mm 0 degree scope
\item
  SRI laparoscopic Pan
\item
  Salo laparoscopic instruments (or pan with laparoscopic needle
  drivers)
\end{itemize}

\textbf{Disposables/Meds}

\begin{itemize}
\tightlist
\item
  Veress needle (with 10mL syringe and saline)
\item
  5mm Z-thread optical port (3 on table, 2 more in room)
\item
  Transverse drape (for port) AND laparoscopy drape
\item
  Confirm choice of port with surgeon. Usual options

  \begin{itemize}
  \tightlist
  \item
    Bard PowerPort 8Fr xx8062
  \item
    Bard PowerPort 8Fr xx8000 (low profile)
  \end{itemize}
\item
  Heparinized saline: 100U/ml (dilute) and 1000U/ml (concentrated)
\item
  1000 drape x3 AND blue paper drapes 4 packs of 2 each = 8 total
\item
  Micropuncture kit available/not open (from Anesthesia)
\item
  Jejunostomy tube: MIC 0201-14
\item
  Silk 2-0 on RB-1 needle (on Surgical Oncology suture cart)
\end{itemize}

\textbf{Position}

\begin{itemize}
\tightlist
\item
  Supine with left arm tucked, right arm on armboard at side.
\item
  Foley catheter: Usually required -- \emph{check with surgeon}
\item
  Lower body Bair Hugger on thighs. ONE layer of blankets on top of Bair
  Hugger. Velcro strap on thighs. NO PILLOW UNDERNEATH LEGS.
\end{itemize}

\textbf{Prep}

Chloroprep (two sticks) of abdomen (need to keep pubis in field, as well
as right anterior superior iliac spine), both costal margins.

If port: RIGHT chest, neck to chin and earlobe, shoulder to include
deltopectoral groove

\textbf{Drape}

If central venous port: Perimeter of field draped with 1000 (clear
adhesive) drapes. Four 1000 drapes around port site:

\begin{itemize}
\tightlist
\item
  Medial border: From Angle of Louis superiorly along midline to chin.
\item
  Superior border: Inferior to jaw (to allow access to right internal
  jugular vein and SCM)
\item
  Laterally: From inferior to ear down to right shoulder
\item
  Inferior: From lateral shoulder medially to Angle of Louis
\end{itemize}

Abdomen: Two 1000 drapes used inferiorly keeping pubis and right
anterior inferior iliac spine in field. This is critical as the far
inferior/lateral RLQ needs to be in the field for optimal port
placement.

Six Blue Paper Drapes around perimeter of field (on top of 1000 drapes)

If central venous port: Transverse sheet TURNED HEAD-TO-FOOT turned at
an angle to keep deltopectoral groove and SCM within the field

Laparoscopy drape skewed to inferior and right to keep pubis and right
ASIS in the field.

Turn on Bair Hugger only AFTER drapes in place

\textbf{Indications}

Laparoscopic jejunostomy is used for enteral nutrition in patients prior
to planned (or possible) esophagectomy or gastrectomy or those for whom
the stomach is otherwise not available (ie after esophagectomy or
gastrectomy). Patients with metastatic esophageal cancer who need
enteral access are generally treated with a gastrostomy, which does not
require feeding with a pump

Preop (Resident) Preop orderset: search for ``Jejunostomy''

Review Clinical Information (Resident) Review staging scans (especially
PET scan) to identify suspicious areas on imaging which need to be
investigated at the time of laparoscopy Outpatient anticoagulation use
(warfarin, Xaralto, aspirin, Plavix) Review dietitian's recommendations
(how many cans of feeding per day?) If patient is scheduled for central
venous port Confirm that a port has not already been placed Prior
history of central venous lines? Confirm location of port placement with
surgeon (left vs right)

\textbf{Operation}

If a central venous port is placed, the port is performed first. See
\hyperref[cv_port_salo]{IJ Port}

Abdominal access is obtained in one of two ways:

Infraumbilical approach using modified Hasson technique. If the
peritoneum is not easily entered, a Veress needle is used to insufflate,
followed by incision of the fascia with a 15 blade, and a 5mm optical
port (Applied Medical Kii Fios First Entry Z-Thread Trocar) Veress
needle inserted in LEFT upper quadrant just inferior to costal margin.
Abdominal entry with 5mm optical Z-Thread port. 5mm port in right upper
quadrant, 5mm port in RLQ just lateral to rectus, 5mm camera port
between RUQ and RLQ ports

The transverse colon is now elevated (using the umbilical port, if used)
and the ligament of Treitz is identified. The proximal bowel is arranged
in a ``C'' configuration to confirm the proximal and distal ends of
bowel.

A site for placement of the jejunostomy is selected on the skin, left
lateral and just superior to the umbilicus. A site is selected on the
bowel in the most proximal site on the jejunostomy selected which would
allow for placement of the jejunostomy without tension, but at least
20cm from ligament of Treitz.

The proximal jejunum is sutured to the anterior abdominal wall with 2-0
silk. This is usually done with a 9″ suture which is introduced into the
abdomen with a needle driver ``Korean Style'' or ``Paraguayan Style.''
Two cm distal to this suture, a diamond of sutures is placed around the
proposed tube site, and one suture placed distal to avoid torsion. The
final arrangement of sutures is one proximal and one distal and 4-6
sutures around the tube. All sutures were marked with hemostatic clips
to facilitate replacement of the tube via fluoroscopy should the tube
become dislodged.

Using Seldinger technique, a 14Fr Cook catheter introducer kit is placed
within the jejunum, followed by a 18Fr dilator and sheath

A MIC 14Fr jejunustomy tube (0200-14) is inserted and the balloon
inflated with 5mL of sterile WATER. The ballon valve is labeled for
7-10mL of water, but this risks obstruction of the jejunum.

The tube is dressed with a BioPatch and a Tagederm dressing.

The abdomen is desufflated and the port sites closed with 4-0 Monocryl,
followed by Dermabond.

\textbf{Endoscopy}

The scope is set up:

\begin{itemize}
\tightlist
\item
  Suction and aspiration valves inserted and working
\item
  Suction tubing attached Biopsy valve attached and not leaking
\item
  Cart set up for recording by powering on the Stryker SDC digital
  capture box
\end{itemize}

A bite block is used and the scope lubricated. A neonatal scope may be
necessary in patients with a tight stricture. Important findings to
record:

\begin{itemize}
\tightlist
\item
  Level in cm from the incisors, of the most proximal area of Barrett's
  esophagus.
\item
  Level in cm from the incisors of the GE junction Appearance of the GE
  junction on retroflexed view.
\item
  Extent of invasion of the tumor into the cardia or fundus.
\end{itemize}

The scope is withdrawn and the hypopharynx suctioned. The liquid from
the `First Step' disinfectant is suctioned through the scope, followed
by water.

\chapter{Lap Gastrostomy}\label{lap_gastrostomy_salo}

\textbf{Room Prep}

\begin{itemize}
\tightlist
\item
  EGD cart near patient's LEFT SHOULDER (with NEONATAL EGD scope)
\item
  If central venous port is placed at the same time:

  \begin{itemize}
  \tightlist
  \item
    Slider bed (Skytron 3600B) with head section
  \item
    Radiology technician alerted to need for C-Arm
  \item
    BK Ultrasound with hockey-stick probe near patient's RIGHT SHOULDER
  \end{itemize}
\end{itemize}

\textbf{Instruments}

\begin{itemize}
\tightlist
\item
  5mm 30 degree scope AND 5mm 0 degree scope
\item
  SRI laparoscopic Pan (available)
\end{itemize}

\textbf{Disposables/Meds}

\begin{itemize}
\tightlist
\item
  Veress needle (with 10mL syringe and saline)
\item
  5mm Z-thread optical port (3 more in room)
\item
  Transverse drape AND laparoscopy drape
\item
  Confirm choice of port with surgeon. Usual options

  \begin{itemize}
  \tightlist
  \item
    Bard PowerPort 8Fr xx8062
  \item
    Bard PowerPort 8Fr xx8000 (low profile)
  \end{itemize}
\item
  Heparinized saline: 100U/ml (dilute) and 1000U/ml (concentrated)
\item
  1000 drape x4 AND blue plastic adhesive drapes 4 packs of 2 each = 8
  total
\item
  Micropuncture kit available/not open (from Anesthesia)
\item
  20Fr Laparoscopic gastrostomy kit in room/not open
\item
  16Fr MIC gastrostomy tube in room/not open
\item
  Gastrostomy 20Fr Pull PEG (in vending machine)
\item
  GI Anchors (``T-fasteners'') in room/not open
\end{itemize}

\textbf{Endoscope Setup -- Neonatal EGD scope}

\begin{itemize}
\tightlist
\item
  Valves attached and working (suction/aspiration/biopsy cap)
\item
  Suction tubing attached
\item
  Connect water bottle to left-hand port
\item
  Gauze sponges, lubricant, plastic ``tray'' from gauze filled with
  water
\item
  \textbf{Yellow} (small) bite block
\item
  ``First step'' sanitizer
\end{itemize}

\textbf{Anesthesia}

\begin{itemize}
\tightlist
\item
  ET Tube taped to left. Head turned to the left on donut
\item
  No EKG electrodes on anterior right chest
\end{itemize}

\textbf{Position}

\begin{itemize}
\tightlist
\item
  Supine with left arm tucked, right arm on armboard at side.
\item
  Foley catheter: usually NOT required -- \emph{check with surgeon}
\item
  Lower body Bair Hugger on thighs. ONE layer of blankets on top of Bair
  Hugger. Velcro strap on thighs. NO PILLOW UNDERNEATH LEGS.
\end{itemize}

\textbf{Prep}

Chloroprep (two sticks) of abdomen (need to keep pubis in field, as well
as right anterior superior iliac spine), both costal margins.

If port: RIGHT chest, neck to chin and earlobe, shoulder to include
deltopectoral groove

\textbf{Indications}

Laparoscopic gastrostomy is used in patients with esophageal obstruction
. Gastrostomy feedings are much easier than jejunostomy, as they can be
administered via syringe or gravity bag. By contrast, jejunostomy
feedings require administration via pump. Gastrostomy is usually done as
an outpatient unless there are concerns for refeeding.

\textbf{Review Clinical Information (Resident)}

Review staging scans (especially PET scan) to identify suspicious areas
on imaging which need to be investigated at the time of laparoscopy
Outpatient anticoagulation use (warfarin, Xaralto, aspirin, Plavix)
Review dietitian's recommendations (how many cartons per day?)

\textbf{Drape}

If central venous port: Perimeter of field draped with 1000 (clear
adhesive) drapes. Four 1000 drapes around port site:

\begin{itemize}
\tightlist
\item
  Medial border: From Angle of Louis superiorly along midline to chin.
\item
  Superior border: Inferior to jaw (to allow access to right internal
  jugular vein and SCM)
\item
  Laterally: From inferior to ear down to right shoulder
\item
  Inferior: From lateral shoulder medially to Angle of Louis
\end{itemize}

Abdomen: Blue adhesive drapes used inferiorly keeping pubis and right
anterior inferior iliac spine in field. This is critical as the far
inferior/lateral RLQ needs to be in the field for optimal port placement
if the patient needs a jejunostomy.

Six Blue Adhesive Drapes around perimeter of field (on top of 1000
drapes)

If central venous port: Transverse sheet TURNED HEAD-TO-FOOT turned at
an angle to keep deltopectoral groove and SCM within the field

Laparoscopy drape skewed to inferior and right to keep pubis and right
ASIS in the field.

Turn on Bair Hugger only AFTER drapes in place

\textbf{Operation}

If a central venous port is placed, the port is performed first. See
\hyperref[cv_port_salo]{IJ Port}

\textbf{Abdominal access} is obtained in one of two ways:

\begin{itemize}
\tightlist
\item
  Veress needle inserted in left upper quadrant just inferior to costal
  margin. Abdominal entry with 5mm optical port at lateral border of
  rectus just superior to umbilicus
\item
  Infraumbilical approach using modified Hasson technique. If the
  peritoneum is not easily entered, a Veress needle is used to
  insufflate, followed by incision of the fascia with a 15 blade, and a
  5mm optical port
\end{itemize}

The insufflation pressure is decreased to 4mmHg and the abdomen vented
to drop the pressure. A 30 degree scope is passed inferior to the
falciform ligament into the left upper quadrant over the lateral segment
of liver. The post of the scope is positioned to the left, allowing
visualization of the lesser curvature of the stomach with the end of the
scope near the left aspect of the falciform.

A site for placement of the gastrostomy is selected on the skin, using a
22Ga needle as a finder. The site is marked, then infiltrated with local
anesthetic and a 5mm transverse incision made.

Endoscopy is performed and the following noted:

Level in cm from the incisors, of the most proximal area of Barrett's
esophagus. Level in cm from the incisors of the GE junction Appearance
of the GE junction on retroflexed view. Extent of invasion of the tumor
into the cardia or fundus.

A bite block is positioned (unless the patient is edentulous). The
endoscope (usually a neonatal scope) is introduced into the esophagus
and the video capture started. If the tumor will not allow passage of
the scope, do not force the scope.

Once the scope is passed into the stomach, the fundus and duodenal bulb
are suctioned.

Insufflation is then reduced to 4mm and the stomach insufflated to find
the optimal location for tube placement which will minimize tension and
will avoid injury to the right gastroepiploic artery. The reduced
laparoscopic insufflation pressure allows endoscopic insufflation of the
stomach.

Gastrostomy tube placement is done either by Pull or Seldinger
technique.

\textbf{Seldinger Gastrostomy}

Four T-fasteners are then used to affix the stomach to the anterior
abdominal wall. These are arranged at 8:00, 10:00, 2:00, 4:00 relative
to the proposed tube site. T-fasteners are not placed inferior to the
tube site to avoid injury to right gastroepiploic vessels.

The J wire is then passed into the stomach, followed by the dilators, up
to 20Fr. A 16Fr MIC gastrostomy tube is introduced and the balloon
inflated with 7mL of \emph{water}.

\textbf{Pull Gastrostomy}

A 20Fr PULL PEG tube kit is opened. The snare and tube are passed to the
upper operator. The snare is passed through the scope and opened in
anticipation of the passage of the wire.

The angiocath from the kit is placed into the stomach through the
abdominal wall. Once the snare has grasped the angiocath, the needle is
withdrawn and the split wire passed through the Angiocath into the
stomach. The snare is adjusted to grasp the split wire, which is pulled
out through the mouth.

The laparoscopic port site is considered clean. The PEG tube, once it is
pulled through the mouth, is considered dirty. The right abdomen is
covered with a towel to protect the laparoscopic site.

The recording is now stopped. The split wire is joined to the PEG tube,
which is pulled into place by the abdominal operator. The tapered
portion of the tube will be the first source of resistance, which may
require firm traction. The split wire (and PEG tube) is dropped of the
table to the patient's left.

Once the tapered portion of the tube is through the abdominal wall skin,
the next point of resistance will be the bumper of the PEG tube passing
through the tumor. In general, if a 5mm neonatal scope can pass the
tumor, a 20Fr PEG tube can pass as well.

The tube is pulled into position and the measurement at the skin noted.
The stomach is aspirated by the upper operator and insufflation is
resumed at 8mmHg. Tension on the PEG tube is adjusted to allow
apposition of the gastric serosa to the abdominal wall. If the tube is
not easily apposed to the abdominal wall, T fasteners (``GI Anchors'')
must be employed.

The scope is withdrawn and the hypopharynx suctioned. The liquid from
the `First Step' disinfectant is suctioned through the scope, followed
by water.

The abdomen is desufflated and the port sites closed with 4-0 Monocryl.

\chapter{Subtotal Gastrectomy}\label{robo_stgastrectomy_salo}

\textbf{Room Prep}

\begin{itemize}
\tightlist
\item
  EGD cart near patient's LEFT SHOULDER (with ADULT EGD scope)
\end{itemize}

\textbf{Instruments}

\begin{itemize}
\tightlist
\item
  5mm 30 degree scope AND 5mm 0 degree scope
\item
  SRI laparoscopic Pan (available)
\end{itemize}

\textbf{Disposables/Meds}

\begin{itemize}
\tightlist
\item
  Veress needle (with 10mL syringe and saline)
\item
  5mm Z-thread optical port
\item
  Robot prostatectomy drape
\item
  1000 drape x4 AND blue plastic adhesive drapes 2 packs of 2 each = 8
  total
\end{itemize}

\textbf{Endoscope Setup --Adult EGD scope}

\begin{itemize}
\tightlist
\item
  Valves attached and working (suction/aspiration/biopsy cap)
\item
  Suction tubing attached
\item
  Connect water bottle to left-hand port
\item
  Gauze sponges, lubricant, plastic ``tray'' from gauze filled with
  water
\item
  \textbf{Yellow} (small) bite block
\item
  ``First step'' sanitizer
\item
  ICG bottle reconstituted with 10ml Water (ask first)
\item
  25\% albumin (to mix with ICG)

  \begin{itemize}
  \tightlist
  \item
    Will mix 2mL of ICG solution with 5mL of 25\% albumin
  \item
    Save remainder of ICG and keep it sterile
  \end{itemize}
\end{itemize}

\textbf{Position}

\begin{itemize}
\tightlist
\item
  Supine arms abducted on arm boards
\item
  Foley catheter
\item
  Bovie pad
\item
  Lower body Bair Hugger on thighs. ONE layer of blankets on top of Bair
  Hugger. Velcro strap on thighs. NO PILLOW UNDERNEATH LEGS.
\end{itemize}

\textbf{Prep}

Chloroprep of abdomen from midline at the level of the nipples to the
pubis, and table to table

\textbf{Review Clinical Information (Resident)}

Review staging scans (especially PET scan) to identify suspicious areas
on imaging which need to be investigated at the time of laparoscopy
Outpatient anticoagulation use (warfarin, Xaralto, aspirin, Plavix)

\textbf{Drape}

Abdomen: Four Blue Adhesive Drapes around perimeter of field (on top of
1000 drapes)

3/4 sheet over thighs

Prostate Drape

\textbf{Operation}

\textbf{Abdominal access} is obtained in one of two ways:

\begin{itemize}
\tightlist
\item
  Veress needle inserted in left upper quadrant just inferior to costal
  margin. Abdominal entry with 5mm optical port at lateral border of
  rectus just superior to umbilicus
\item
  Infraumbilical approach using modified Hasson technique. If the
  peritoneum is not easily entered, a Veress needle is used to
  insufflate, followed by incision of the fascia with a 15 blade, and a
  5mm optical port
\end{itemize}

\textbf{Ports}

Location of ports (cephalad-caudad) depends upon proximal extent of
tumor.

Abdominal entry in RUQ 8-9cm from midline with OptiView port. This will
be upsized to 8mm robot port (\#1)

Port \#4 as far lateral as possible in LUQ.

Port \#3 in LUQ midway between midline and Port \#4

Port \#2 in midline

Assisting port in LLQ: - 5mm or 12m port inferior to the midway point
between Port \#3 and Port \#4.

Laparoscopy is performed, and biopsies taken as needed. Once it is
determined that there is no peritoneal disease, port \#3A is changed to
a 12mm robotic stapler port.

Retraction port right mid-quadrant for flexible liver retractor held by
Endosopic Bookwalter

Flexible liver retractor can be omitted in very distal lesions. Other
options include suspending the lateral sector of liver with a 10cm
Penrose drain or using a 0 silk on Keith needle to suspend falciform
ligament and distract it to patient's right

\textbf{Porta Hepatis}

\begin{itemize}
\tightlist
\item
  Mark pylorus with cautery by making 2-3 dots 1cm proximal to the
  pylorus
\item
  Dissect pylorus off porta hepatis
\item
  Mobilize duodenum if needed
\end{itemize}

\textbf{Mobilize Greater Curvature}

Can be done with bipolar + monopolar scissors or Extend vessel sealer.
Enter lesser sac and check for peritoneal disease posterior to stomach
on the anterior surface of pancreas

\textbf{Divide Duodenum}

Dissection is carried out to the right of the right gastroepiploic nodal
basin to divide between duodenum/pancreas and transverse mesocolon.
Dissection is brought superiorly to the edge of the duodenum at the
proposed site of transection. Duodenum is divided with a Blue load
stapler

**EGD with ICG injection*

Once the duodenum is dissected, an injection of ICG mixed in albumin is
made in the submucosal around the tumor.

\textbf{Division of greater curvature}

The gastrocolic omentum is dissected at this point (to allow ICG to
transit to lymph nodes). This can be taken up to the inferior edge of
the spleen.

\textbf{Right gastroepiploic pedicle dissection}

Right gastroepiploic artery dissected free and triple-clipped. Right
gastroepiploic vein can be taken with vessel sealer or double-clips.

\textbf{Right gastric artery}

Nodes to the left of the porta hepatis (12a) are dissected and kept with
the specimen. The right gastric artery is divided.

\textbf{Central node dissection}

The antrum is reflected to the patient's left. Nodes are dissected from
the common hepatic artery (Station 8), proximal splen artery (Station
11), and celiac axis (Station 9). The coronary vein is divided with the
vessel sealer. The left gastric artery is triple-clipped. Dissection is
now carried out posteriorly. The lesser curvature is cleared of lymph
nodes (Stations 3) and right paracardial nodes (Station 1) taken
depending upon tumor location.

\textbf{Proximal division}

The greater curvature is dissected proximally, typically taking several
inferior short gastric vessels. The stomach is divided with several Blue
loads.

\textbf{Extraction and frozen section}

An extraction incision is made and the specimen placed in a 12/15mm bag
and removed and sent for frozen section. The extraction site can be
closed or used to create with roux limb. The extraction incision is
closed with two sutures of \#1 Stratafix

\textbf{Roux-en-Y reconstruction}

The division site for jejunum is selected 25cm distal to Treitz. ICG can
be injected intravenously to define the arcades. Up to two arcade
vessels can be divided to yield more length of vascularlized bowel. It
is important that this division of the mesentery be made central to the
vasa recta. The stapler is positioned in port \#3A, if it is not
already. After initial division of the mesentery, the bowel is divided
with a white load stapler.

A roux limb 40cm in length is measured and a jejunojejunostomy
constructed with a white stapler load. Enterotomies made with a hook
cautery. The bowel defect is closed with two running 3-0 Stratafix
sutures on RB-1 needle.

The bare area is fenestrated to the left of the middle colic artery. The
Roux limb is transposed through the mesentery to superior to the
transverse mesocolon.

\textbf{Gastrojejunostomy}

Once the frozen section confirms a clear margin, the Roux limb is
positioned posterior to the proximal gastric pouch. A gastrostomy is
made on the posterior stomach at least 1cm from the staple line. An
enterotomy is made on the antimesenteric border of bowel 7cm from the
end of the limb. The anastomosis is made with a Blue load. The defect is
closed with two sutures of 3-0 Stratafix.

\textbf{NG tube}

After the anastomosis, an NG tube is placed by Anesthesia and bridled.

\textbf{Drains}

Solitary 19Fr drain placed posterior to the anastomosis

\textbf{Closure}

12mm stapler site closed with 0 Vicryl in figure-of-eight fashion.
Transversus abdominus place blocks placed.

\chapter{Esophagectomy 1 Stage}\label{esophagectomy-1-stage}

\section{Indications}\label{indications}

MI Ivor Lewis (``One Stage'') esophagectomy is the most common approach
to esophagectomy. The patient is positioned in `corkscrew' position to
allow simultaneous access to the abdomen and chest by prepping both
abdomen and chest.

\section{Preop Orders}\label{preop-orders}

\begin{itemize}
\tightlist
\item
  Type/Screen
\item
  No subcutaneous heparin if an epidural catheter is planned
\end{itemize}

\section{Room Prep}\label{room-prep}

Double-decker Back Table

EGD Cart (check w/ Salo regarding adult vs neonatal scope)

Doppler box with foot pedal

BK Ultrasound (with laparoscopic probe)

Walter `Long Arm' retractor

Salo Positioning Cart:

\begin{itemize}
\tightlist
\item
  Four black side-rail clamps
\item
  Four rectangular lateral positioners
\item
  Yassargil socket
\item
  Well-leg holder (used as an arm holder)
\end{itemize}

Harmonic Scalpel generator box

Critical supplies (chek prior to pt in room)

\begin{itemize}
\tightlist
\item
  Stapler 25mm DSTXL
\item
  Orvil 25mm
\item
  Stapler 21mm DSTXL
\item
  Orvil 21mm
\item
  Echelon 60mm stapler with 10 gold loads and 2 gray loads
\item
  Gel Port
\end{itemize}

\newpage

\textbf{Anesthesia}

\begin{itemize}
\tightlist
\item
  Dual-lumen ETT tube (\textbf{taped to left})
\item
  Arterial line (left arm will be on arm board)
\item
  Anesthesia will usually place epidural catheter either in preop or in
  OR
\end{itemize}

\section{Position}\label{position}

Supine on blue foam pad. No pillow under knees. Keep siderails clear by
tucking drawsheet into gap between rails.

Mark upper midline with skin marker

Foley catheter (with Criticor temp sensor)

Bovie pad

Lower-body Bair Hugger at level of thighs

Hair clipped from abdomen, right chest, and right axilla.

Pre-existing jejunostomy

\begin{itemize}
\tightlist
\item
  Prep into field
\item
  Remove any eschar at site
\item
  Secure to skin inferiorly with 0 silk
\end{itemize}

Shoulders are shifted to the right in preparation for `corkscrew'
positioning

Lateral positioners positioned with pad extending from greater
trochanter inferiorly.

Velcro strap over thighs and over lateral positioners

Yassargil socket attached to headpiece of bed on left side.

Well-leg holder attached to Yassargil socket and used to support right
forearm (which will cross body). Right arm crosses body and is supported
on well-leg holder.

Lateral positioner placed posterior to spine and scapula to support
right chest, allowing access for thoracoscopy.

Arm holder is dropped towards the floor enough that right arm is brought
forward.

\textbf{Prep}

Chloroprep (two sticks) of abdomen, right chest, right axilla.
Particular attention to prepping as far as possible to the left lateral
side and the right lateral side. Nipple prepped into field.

\textbf{Drape}

Proximal right arm draped with 1000 (clear adhesive) drape.

Two blue plastic ``U'' drapes with center of the ``U'' on either lateral
side with tails forming the perimeter of the field Trauma drape. All of
right chest and axilla is kept within the field, as is the lateral
aspect of the left upper quadrant. The field does not need to extend
inferior to the umbilicus. In general, it is usually possible to keep
all of these area in the field without cutting the drape, except in very
large patients. Ioban strips (4'') aound the periphery of the surgical
field after the trauma drape. Laparoscopic cords (gas, light cord,
camera) to tower at left shoulder. Suction irrigator brought off field.
Laparoscopic LigaSure (2 bars) and Bovie (30/30).

\subsection{Time Out}\label{time-out}

Operation header on consent

Tumor location and likelihood of division of esophagus from abdomen
(two-phase operation) vs division of esophagus from chest (four-phage
operation).

Blood:

\begin{itemize}
\tightlist
\item
  Surgeon's expectation of blood loss
\item
  Availability of blood (type/screen vs type/cross).
\end{itemize}

\emph{Comorbidities:}

Cardiopulmonary disease. If echo, ejection fraction and aortic valve
area (if abnormal)

Beta blockage: Note whether patient on home beta blockade and if so,
whether home medication was taken the morning of surgery.

Anticipated Intraop Problems:

\begin{itemize}
\tightlist
\item
  Possibility of tension left pneumothorax due to carbon dioxide entry
  into left chest during mediastinal dissection
\item
  Expectation of ventilatory difficulties due to carbon dioxide entry
  into the right chest Gastric mobilization -- Greater Curvature
\end{itemize}

\section{Gastric Mobilization}\label{gastric-mobilization}

7cm upper midline incision for handport. Incision is centered over
pylorus. GelPort inserted, and abdomen insufflated to 15mmHg. Two 5mm
ports placed LUQ. Medial LUQ 5mm port placed in angle between left
costal margin and superior edge of GelPort ring. Lateral LUQ 5mm port
placed as far lateral as possible. Depending upon visualization, a third
port may be required between these two and somewhat more inferior.

If feasible, division of gastrocolic omentum starts by delivering
transverse colon into GelPort and dividing ligament with cautery and
LigaSure in the avascular plane just cephalad to transverse colon.
Dissection proceeds as far proximal and distal and feasible. It is
important to avoid damage to the right gastroepiploic artery.

Colon is returned to abdomen and gastricolic ligament divided going
using LigaSure, taking care to avoid the colon and the right
gastroepiploic artery. For patients with a bulky omentum, it may be
helpful to place an additional port in the left mid-quadrant for the
camera to facilitate dissection of omentum off the transverse colon.
Stomach is retracted to the patient's right with the back of the left
hand, placing the gastrocolic ligament (and short gastric arteries) on
stretch. Left gastroepippoic artery divided with LigaSure near its
origin. Short gastric arteries divided with Ligasure close to spleen. As
superior aspect of short gastrics is reached, the dissection plane
shifts medially to create a tunnel towards the base of the left crus.
This places the most superior short gastric vessels on stretch and
facilitates their division. Once all short gastric vessels divided,
peritoneum tethering fundus to diaphragm is incised, and fundus brushed
medially.

\subsection{Distal mobilization}\label{distal-mobilization}

Attention is directed to mobilizing the lateral aspect of the duodenum.
This can either be accomplished with the camera through a LUQ port and a
45 degree camera or by placing an additional 5mm port in the RLQ. Hook
cautery is used to incise the connective tissue lateral and posterior to
the duodenum. The gastrocolic ligament is now dissected distally, taking
care to preserve the integrity of the right gastroepiploic vessels. An
areolar plane generally exists between the fat pad containing the right
gastroepiploic vessels and that containing the transverse mesocolon
vessels.

\subsection{Left gastric artery}\label{left-gastric-artery}

Lymph nodes around the celiac axis and left gastric artery are now
dissected. The extend of dissection depends upon the tumor location and
the presence of nodes her either based upon imaging or palpation.
Dissection begins on the superior edge of the pancreas, and proceeds
superiorly to the right crus. The left gastric (coronary) vein is
usually located to the left of the arery and is divided with the
LigaSure. In a two-phase approach, the left gastric artery is now
divided with a 30mm gray load (2.0mm) Echelon linear stapler. For
patients with mid-esophageal tumors, (for whom a four-phase approach is
used), if there is any question about the resectability of the tumor,
division of the left gastric gastric artery is generally deferred until
the second abdominal phase.

\subsection{Mediastinal dissection}\label{mediastinal-dissection}

The esophagus is now dissected circumferentially at the gastroesophageal
junction. The peritoneum overlying the diaphragm is incised, and
circumferential dissection of esopahgus performed. On the left side, it
is helpful to distract the left crus laterally with a Prestige clamp
placed on the left crus The right pleura is widely entered in otder to
both facilitate the thoracic dissection of the esophagus and placing the
conduit into the right chest in preparation for the final thoracic
phase. Division of Esophagus (Two-Phase only) In a four-phase approach,
the esophagus is divided from the right chest during the first (of two)
thoracic phases. In a two-phase approach, the esophagus is divided from
the abdomen by reaching up into the mediastinum to divide the esophagus
above the tumor. This is only feasible for tumors of the
gastroesophageal junction. In a two-phase approach, the esophagus is
divided with a 60mm Medium-Thich (Gold) Echelon TriStaple stapler.

\subsection{Division of Esophagus
(Four-Phase)}\label{division-of-esophagus-four-phase}

Penrose Drain (Four-Phase only) In a four-phase approach, the esophagus
is divided from the chest. In order to facilitate the thoracic
dissection, a 1/4″ penrose drain is tied around the distal esophagus and
the drain is slid cephalad into the mediastinum. The `tails' of the
drain are directed into the right chest so that they can be grasped from
the right chest during the thoracic phase and can be used to provide
traction on the esophagus

\subsection{Entry into right chest}\label{entry-into-right-chest}

The pneumoperitoneum in the abdomen is vented and the gas pressure
turned down to 8mmHg in preparation for the thoracic phase. The bed is
rotated to the left 20 degrees. A site for entry in to the left chest is
selected just posterior to the tip of the scapula. An incision is made
here and a 5mm optical port with a 5mm 0 degree scope used to enter the
chest. The chest is insufflated at 8mmHg of carbon dioxide, which helps
to both copllapase the lung and depress the diaphragm. Two 12mm ports
are placed. The more superior is placed lateral and superior to the
nipple. The inferior port is placed just lateral to the diaphragmatic
reflection. It is ciritcal to avoid injury to the diaphragm (and liver)
with the inferior port placement. A mini-thoracotomy incision is placed
along the mid-axillary line, frequently in the same interspace as the
inferior/anterior 12mm port. The chest is entered just superior to the
rib and the intercostal muscles divided with the LigaSure device to
allow the ribs to separate. A narrow Deaver retractor is used to guage
the space between the ribs, as the width of the retractor approximates
the diameter of the 25mm stapler. A 5mm `U' port is generally placed as
high as possible midway between the scapular tip port and the
anterior/superior 12mm port t0o the 4 right 5

Division of Esophagus (Two-Phase) In patients with low-lying tumors, it
is possible to divide the esophagus above the tumor from the abdominal
approach. This evaluation is facilitated by review of the properative
endoscopy (and EUS), and the PET scan, particularly the PET obtained
prior to neoadjuvant chemoradiation.

After dissection of the mediastinum, a Echelon stapler with a 60mm
Medium-Thick (Gold) load is inserted through a 12mm port either placed
either through the GelPort or by upsizing the most lateral LUQ

Construction of Conduit The distal esophagus and stomach is exteriorized
through the GelPort. The lesser curvature vessels are divided with the
LigaSure 7-9cm cephalad from the pylorus. The stomach is now placed on
stretch along the greater curvature. An Echelon Medium-Thick (Gold)
stapler is used to construct a 5-6cm wide gastric tube. In constructing
the conduit in a patient with a tumor of the GE jnction invading into
the cardia, it is important to be certain that the staple line to
construct the conduit stays clear of the tumor. Patients with tumors
invading the cardia are at risk of a positive distal margin, meaning
that microscopic tumor may be left in the wall of the gastric conduit.
To make things more complicated, in patients with low-lying esophageal
or GE junction tumors, not all of the length of the conduit are needed
in order to reach to the level of the esophageal transection. After
construction of the anastomosis, the `extra' cephalad portion of the
conduit (near the angle of His) are excised as the `additional gastric
margin.' In order to distinguish which portions of the conduit staple
line which will be used to replace the esophagus and those which will be
included in the `additional gastric margin', both sides of the gastric
conduit staple line are marked with sutures designated `A', `B', `C',
etc proceeding from the Angle of His to the antrum.

The distal esophagus and GE junction are now sent for frozen section.

Feeding jejunostomy The ligament of Trietz is identified and the jejunum
is identified 20cm distal and marked with a directional suture. A site
is selected to the left of the handport incision. A 16Fr Cook Introducer
Kit is used to pass a 16Ga needle, followed by a J wire, through the
left rectus muscle. A skin incision 4mm in length is made adjacent to
the J wire. The 16Fr dilator and sheath are now passed through the
rectus muscle and the dilator and wire removed. A 14Fr Jejunostomy Tube
0301-14 is selected and the `wings' trimmed off with a scalpel. The
jejunostomy tube is passed through the peelaway sheath, which is
removed.

A pursestring suture 1.5cm in diameter is placed on the antimesenteric
border of the jejunum. The pursestring is started and ends on the
lateral aspect. A second 16Fr Cook Introduced Kit is used to introduce
the J wire through the center of the pursestring.

Placement of Drains Two (or three) 19Fr full-fluted Blake drains (72230)
are placed:

JP1: placed into the left pleura through the hiatus. Drain is brought
out through the most lateral 5mm port site on the LUQ JP2: placed into
the right pleura through the hiatus. Drain is brought out through the
next most medial 5mm LUQ port site JP3: (optional) placed in the abdomen
posterior to the left lateral segment of the liver and brought out
through the most medial 5mm LUQ port site Transposition of Conduit The
gastric conduit is now placed into the right pleura through the
mediastinum, with the assistance of a laparoscopic Babcock and gentle
pressure on the greater curvature with the fingers.

Entrance into the R chest (Two Phase) For two-phase operations, ports
are now placed into the right chest, as stated above . In similar
fashion, the inferior pulmonary ligament in dissected and the lung
reflected anterior.

Anastomosis The right chest is entered and the right lung reflected
anterior with the paddle placed in the superior/anterior port. The
gastric conduit is placed into the mediastinum by tucking it medially
from the right pleura into the posterior mediastinum, in order to allow
the conduit to take the most direct path from the hiatus to the proximal
esophagus. Gentle superior tension is now applied to the conduit in
order to eliminate redundancy. The paddle retractor is now moved from
the anterior/inferior port to the anterior/superior 12mm port and the
lung reflected anterior and inferior.

OrVil The OrVil device is now used to place a EEA anvil into the distal
esophagus. In general, a 25mm size is selected, unless the patient has
particularly small frame, in which case a 21mm size is used. Two stay
sutures of 2-0 silk on RB-1 needles are placed in the center of the
esophageal staple line 3mm apart. A Harmonic scalpel is used to divide
the staple line. The OrVil device is passed through the mouth and is
passed through the fenestration in the staple line. The anvil portion of
the OrVil is oriented so that the rounded portion is placed against the
roof of the mouth. The OrVil is guided into the hypopharynx by pulling
on the the tube end of the device. As the anvil approaches the
hypopharynx, the jaw is pulled forward to allow passage of the anvil.

The shaft of the anvil is brought through the esophageal staple line,
and the tube disconnected from the anvil by cutting the blue sutures.

The superior end of the conduit is opened along the staple line and the
DST XL stapler (matching the diameter of the OrVil) introduced through
the Alexis device. The stapler shaft is placed into the open end of the
gastric conduit and the conduit pulled over it (`sock over shoe'). The
stapler spike is brought out through the greater curvature. The anvil is
grasped with a Maryland grasper placed through a superior 5mm port and
the two components of the stapler are mated and the stapler tightened
and fired. The knob of the stapler is rotated two turns
counter-clockwise until a click is felt, at which time the anvil will
flip. The stapler is withdrawn and the donuts examined and sent for
pathologic exam.

The anastomosis is completed by firing a Echelon Medium-Thick (Gold)
linear) across the conduit cephalad to the anastomosis. The excess
conduit is sent for pathologic exam as `additional gastric margin.'

NG Tube A Covidien Salem Sump 18Fr nasogastric tube is passed by the
anesthetist. A laparoscopic BK ultrasound is used to monitor the passage
of the NG tube through the esophagus and into the gastric conduit. The
NG tube is passed to the level that all for dots are outside the nose,
with the 4th dot at the nares. The NG tube is secured with an AMT
bridle.

Chest Tube A 28Fr Blake chest tube is placed through the
anterior/inferior 12mm port and is positioned into the posterior
mediastinum. JP2 is placed near the gastric conduit. The right lung is
re-inflated.

Closure The stapler access port incision is closed with 0 Vicryl to
approximate the serratus muscle. The incisions are closed with 4-0
Monocryl followed by Dermabond.

\chapter{Esophagectomy Robot Ivor
Lewis}\label{esophagectomy-robot-ivor-lewis}

\section{Preoperative evaluation}\label{preoperative-evaluation}

PET scan evaluted for extent of tumor, focusing on proximity to the
carina.

CT scan evaluated for presence of replaced left hepatic artery.

\section{Preop Area}\label{preop-area}

Heparin is held if epidural catheter is planned

Clip hair from abdomen and right chest, including right axilla. Hair is
clipped to base of pubis.

\section{Room Prep}\label{room-prep-1}

Double-decker Back Table

EGD Cart (check w/ Salo regarding adult vs neonatal scope)

BK Ultrasound (with robotic probe AND either hockey stick or small T
probe). Will need sterile ultrasound jelly.

Walter `Long Arm' retractor

Bean Bag

Salo Positioning Cart:

\begin{itemize}
\tightlist
\item
  Four black side-rail clamps
\item
  Four rectangular lateral positioners
\item
  Yassargil socket
\item
  Well-leg holder (used as an arm holder)
\end{itemize}

Critical supplies (check prior to pt in room)

\begin{itemize}
\tightlist
\item
  Stapler 25mm DSTXL
\item
  Orvil 25mm
\item
  Stapler 21mm DSTXL
\item
  Orvil 21mm
\end{itemize}

\newpage

\textbf{Anesthesia}

\begin{itemize}
\tightlist
\item
  Dual-lumen ETT tube (taped to left)
\item
  Arterial line (left arm will be on arm board)
\item
  Anesthesia may place epidural in preop
\end{itemize}

\section{Abdominal Position}\label{abdominal-position}

Supine on blue foam pad.

Foley catheter (with Criticor temp sensor)

Bovie pad

Lower-body Bair Hugger at level of thighs, followed by single blanket
and Velcro strap

Hair clipped from abdomen, right chest, and right axilla.

Pre-existing jejunostomy

\begin{itemize}
\tightlist
\item
  Prep into field
\item
  Remove any eschar at site
\item
  Secure to skin inferiorly with 0 silk
\end{itemize}

Velcro strap over thighs and over lateral positioners

\textbf{Prep}

Patient is positioned frog-leg for prep (and is returned to supine
during draping). Chloroprep (2 sticks) of abdomen and lower chest up to
the level of nipples (in case urgent chest tube needed). Particular
attention to prepping as far as possible to the left lateral side and
the right lateral side. Prep continues inferior to base of pubis. At the
end of prep, bilateral groins are prepped.

\textbf{Drape}

Transverse 1000 drape is placed at the base of the pubis. On either
side, the drape proceeds inferior and lateral to include both groins.
Top gloves are changed.

Blue utility (paper) drape is placed transversely just above level of
pubis. While the abdominal skin incisions are being closed,

1000 Drapes x4 in rectangular fashion with top of field 10cm above
xiphoid and bottom at base of penis. Ioban strips (4'') aound the
periphery of the surgical field after the trauma drape. Laparoscopic
cords (gas, light cord, camera) passed off. Suction irrigator brought
off field.

\subsection{Time Out}\label{time-out-1}

Operation header on consent

\textbf{Time Out}

Review:

\begin{itemize}
\tightlist
\item
  Consent
\item
  Comorbidities, including cardiopulmonary disease
\item
  Anticoagulation plan. If preoperative heparin was held for epidural,
  plans for timing of its administration. Circulator and anesthetist set
  alarms.
\item
  Prior abdominal or thoracic operations
\item
  Allergies
\item
  Blood

  \begin{itemize}
  \tightlist
  \item
    Surgeon's expectation of blood loss
  \item
    Availability of blood (type/screen vs type/cross).
  \item
    Blood product administration limitations, if any
  \end{itemize}
\item
  Introduction by name and role
\item
  Co-surgeons: Confirm availability prior to incision
\item
  Operative plan
\item
  Request ICU bed in STICU
\end{itemize}

\emph{Comorbidities:}

Cardiopulmonary disease. If echo, ejection fraction and aortic valve
area (if abnormal)

Beta blockage: Note whether patient on home beta blockade and if so,
whether home medication was taken the morning of surgery.

Anticipated Intraop Problems:

\begin{itemize}
\tightlist
\item
  Possibility of tension left pneumothorax due to carbon dioxide entry
  into left chest during mediastinal dissection
\item
  Expectation of ventilatory difficulties due to carbon dioxide entry
  into the right chest Gastric mobilization -- Greater Curvature
\end{itemize}

\section{Laparoscopy}\label{laparoscopy}

Pneumoperitoneum established with Veress needle at left subcostal or
inferior to umbilicus \emph{or} modified Hasson approach. Abdominal
entry with 5mm optical port at future location of Port \#2 in
mid-axillary line at level of umbilicus. Laparoscopy to look for
peritoneal implants.

\section{Robotic ports}\label{robotic-ports}

Port \#3 8mm left mid-axillary line 3cm above level of umbilicus
\(\rightarrow\) camera.

hook \textbar{} monopolar scissors \textbar{} Extend vessel sealer.

Port \#4 left anterior axillary line. 8mm port \(\rightarrow\)

Port \#2 right mid-clavicular line. Switch from 5mm to 12mm robotic
stapler port. Place `sleeper' 0 Vicryl figure-of-eight suture and secure
with a hemostat.

Port \#1 left anterior axillary line 8mm port \(\rightarrow\)
fenestrated bipolar

Assistant port 8mm port inferior to umbilicus. Port \#3 will be switched
to this position for jejunostomy later in the case.

Nathanson liver retractor placed in midline using the 5mm optical port
to create a tract. Nathanson held in place with Salo Endoscopic
Bookwalter secured to bedrail just inferior to right armboard.

\section{Abdominal Phase}\label{abdominal-phase}

\subsection{Lesser Sac Dissection}\label{lesser-sac-dissection}

Gastrohpeatic ligament divided. Celiac axis evaluated for resectability.
Right crus dissected to ensure resectability. Dissection carried
anteriorly and left crural dissection begun. Dissection posterior to
esophagus.

Central node dissection begun, including common hepatic, left gastric,
celiac, and proximal splenic aftery nodes. If replaced left hepatic
artery present, may require preservation of left gastric artery. Left
gastric artery divided with clips.

Posterior gastric dissection. Dissection is carried posterior towards
most superior short gastric arteries (as these can be difficult to
approach from the greater curvature side).

\section{Distal Greater Curve}\label{distal-greater-curve}

The right gastroepiploic artery is identified and the lesser sac entere
to the left of midline. Dissection proceeds to the patient's right. This
may require shifting ports:

\begin{itemize}
\tightlist
\item
  Camera to Port 2
\item
  Vessel sealer to Port 3
\item
  Tip Up to Port 4
\end{itemize}

Right gastroepiploic artery carefully preserved and gatrocolic omentum
incised. Dissection proceeds to right to mobilize colon off duodenum.
Kocher maneuver performed and omental adhesions to gallbladder taken
down.

\section{Proximal Greater Curve}\label{proximal-greater-curve}

Instruments are returned to original positions and dissection proceeds
proximally, preserving the left gastroepiploic artery near its orgin.
Posterior dissection from the pancreas and splenic artery is performed,
followed by division of short gastric vessels

\section{Mediastinal dissection}\label{mediastinal-dissection-1}

Penrose drain is placed around esophagus and secured with stapler or
Weck clips. Dissection proceeds to level of inferior pulmonary veins.
Penrose is tucked into the mediastinum.

When Penrose is placed, confirm that heparin has been given.

\section{Pyloric Drainage - Pyloromyotimy
(optional)}\label{pyloric-drainage---pyloromyotimy-optional}

Pyloromyotomy is performed with camera in Port \#2 and hook in Port \#3
and Tip up in Port \#4 (may need to move camera to midline)

\section{Pyloric Drainage - Pyloroplasty
(optional)}\label{pyloric-drainage---pyloroplasty-optional}

Stay sutures of 2-0 silk are placed cephalad and caudad to pyorus.
Pyloric incision is made with hook cautery. Defect is closed with 2-0
silk cut to 6'' (have 2 sutures ready)

\section{Creation of Conduit}\label{creation-of-conduit}

``Ruler'' is created by making 1cm marks with black marker on Tip up.
Tip up grasper place into Port \#4. Robotic stapler is introduced into
Port 2. Staple line starts \textasciitilde5cm from pylorus. First firing
with Green (4.8mm) stapler load followed by Blue loads. Conduit is
created with goal of 3-4cm. Assistant retracts fundus towards left upper
quadrant. Vicryl suture is placed 2cm above the start of the staple line
(approximately 7cm from pylorus).

\section{Preparation of Conduit}\label{preparation-of-conduit}

Camera is moved to Port \#2. Suture cut needle driver is introduced into
Port 3 and Tip Up into Port 4. Conduit is sutured to specimen with
multiple sutures of 2-0 silk (cut to 6''). Arm 4 is undocked and 19 Fr
Blake drain is introduced through the port. The drain is sutured to
specimen with 2-0 silk cut to 4''. JP2 is brought out through the LLQ
5mm assistant port or through the port for Arm \#4.

\section{Prep for Thoracic Phase}\label{prep-for-thoracic-phase}

Once conduit has been constructed, confirm presence in room of bean bag,
two pillows, and an extra OR bed for transition to lateral decubitus
position.

\section{Preparation for ICG}\label{preparation-for-icg}

One vial of ICG is resuspended with 10mL of water. 5mL of the solution
is brought onto the sterile field in a 5mL syringe, which is fitted with
a 25Ga needle. The remainder of the ICG is given to Anesthesia for later
(intravenous) use.

\section{Jejunostomy}\label{jejunostomy}

Camera is moved to midline (or Port 2) and Tip up into Port 3. Camera
re-oriented to 30 degrees up. Ligament of Treitz is identified and site
for jejunostomy marked at 30cm with marking pen on a Prestige with 1 dot
for selected site and 2 transverse marks distally. 3-0 Stratafix on RB-1
is introduced with a non-cutting needle driver into Port 3. Insufflation
pressre reduced to Jejunum is sutured to abdominal wall starting at
10:00. Pursestring is run to 2:00 position and 16Ga needle from a 14Fr
Cook Introducer kit passed through the abdominal wall into the bowel.
Insufflation of air confirms proper position. J wire and 14FR dilator
introduced and left in place. Pursestring continues around proposed site
and then proceeds cephalad to provide a broader area of fixation between
bowel and abdominal wall. Tract is dilated with 18Fr dilator. 14Fr MIC
jejunostomy (trimmed to 10cm) is introduced through 18Fr sheath and
positioned within the bowel. 5mL of water in a slip-tip (non-Luer) 10ml
syringe (from the Cook introducer kit) are instilled into the
jejunostomy balloon.

\section{Closure}\label{closure}

Robot is disengaged and kept sterile. Laparoscopic cords are placed in a
basin on a ring stand to be reused for the thoracic portion.

Abdomen is desufflated and ports removed. Port \#2 site is closed by
tying 0 Vicryl suture. Port sites are closed with 4-0 Monocryl followed
by Dermabond.

\section{ICG injection into groin.}\label{icg-injection-into-groin.}

While port sites are being closed, ``loincloth'' is removed and hockey
stick or small T ultrasound transducer is brought onto the field. Under
ultrasound guidance with sterile ultrasound jelly, ICG is injected into
groin nodes.

\chapter{Thoracic Phase}\label{thoracic-phase}

\section{Transition to Lateral
position}\label{transition-to-lateral-position}

Patient is transferred to a spare OR table and a bean bag placed.
Patient is positioned lateral with tip of scapula at break in the table.
Pillows are used under and between the knees. Arm is supported on a Well
Leg Holder affixed to the left headrail with a Yassargil clamp. Lateral
positioners used on patient's left side to stabilize bean bag and allow
tilt to the left (into modified prone position).

Lower and upper body bair hugger used. Right arm is elevated for prep
and lowered once drapes in place.

\section{Prep and drape}\label{prep-and-drape}

Chloroprep including axilla and proximal right arm. Drape axilla/upper
arm with 1000 drape.

Two blue ``U'' drapes used, folowed by trauma drape (or Universal).
Ioban strips around edges. Once Ioban in place, arm is lowered

Laparoscopic instruments are kept sterile and may be returned to the
field if pleural adhesions are encountered.

\section{Thoracic Ports}\label{thoracic-ports}

\begin{itemize}
\tightlist
\item
  Arm \#4 4th ICS 1cm medial to scapula
\item
  Arm \#3 6th IS more medial than 4th arm and close to posterior
  axillary line. Camera Port.
\item
  Arm \#2 8th ICS 12mm port for Vessel sealer or stapler
\item
  Arm \#1 10th ICS on the scapular line or posterior on 10th ICS A1
  Assitant 9th ICS triangulated beween 1 and 2. Alexis port used for
  extraction incision A2 Assistant 5th ICS.
\end{itemize}

DaVinci docking with ``Anatomy: Renal'' ``Patient Right'' Target on
azygous vein

\section{Thoracic Dissection}\label{thoracic-dissection}

\begin{enumerate}
\def\labelenumi{\arabic{enumi}.}
\item
  Dissect plane between lung and anterior aspect of esophagus. Start
  dissection at the pericardium and proceed cephalad. Once the area of
  the right mainstem bronchus is approached, keep plane more
  superficial.
\item
  Dissect onto pericardium. Confirm that tumor will comes off
  pericardium. Leave pulmonary vein dissection to later. Counter
  traction with suction.
\item
  Dissection over parietal pleura up to azygous arch superficially.
  Dissect lower aspect of the azygous vein.
\item
  Superficial dissection superior to azygous vein to reach superior
  extent of dissection at inferior border of sygous svein, then proceed
  inferiorly along the posterior aspect of the esophagus. Dissection
  proceeds to hiatus.
\item
  Divide azygous vein with either vascular load stapler or Gray load
\item
  Dissection proceeds to hiatus
\item
  Divide thoracic duct with Hem-O-Lock clips at least 5cm superior to
  hiatus.
\item
  Return to dissection plane anterior to the esophagus and deepen plane.
  Leave dissection of Station 7 nodes until esophagus is completely
  freed up. Can switch instruments in Arm 1 and Arm 2.
\item
  Dissection of Station 4R nodes.
\item
  Dissection of subcarinal (Station 7) nodes. Start by dissecting
  esophagus from trachea, then outline inferior border of right and left
  mainstem bronchi. Keep Station 7 lymph nodes on the specimen.
\end{enumerate}

\section{Delivery of Conduit into
chest}\label{delivery-of-conduit-into-chest}

The conduit is delivered into the thorax. This is done using a
laparoscopic graster. Ensure staple line faces you (right side). The
specimen is removed through the assistant port incision. An Alexis port
is used to seal the extraction site. Stomach is pulled into the chest
until the \textbf{Vicryl marker suture} is seen.

Division of esophagus just above level of azygous vein. ICG can be used
to check perfusion. If there is any question about the length of the
conduit, confirm the length of the conduit prior to division of the
esophagus.

Penrose drain is recovered. Specimen is sent for frozen section with
instructions to freeze proximal margin.

Once conduit is delivered into chest, set up BK ultrasound to display on
robotic screen.

\section{Preparation of conduit}\label{preparation-of-conduit-1}

The gastric conduit is placed into the mediastinum by tucking it
medially from the right pleura into the posterior mediastinum, in order
to allow the conduit to take the most direct path from the hiatus to the
proximal esophagus. Gentle superior tension is now applied to the
conduit in order to eliminate redundancy.

\section{Circular stapled anastomosis -
Orvil}\label{circular-stapled-anastomosis---orvil}

OrVil The OrVil device is now used to place a EEA anvil into the distal
esophagus. In general, a 25mm size is selected, unless the patient has
particularly small frame, in which case a 21mm size is used. Two stay
sutures of 2-0 silk on RB-1 needles are placed in the center of the
esophageal staple line 3mm apart. A Harmonic scalpel is used to divide
the staple line. The OrVil device is passed through the mouth and is
passed through the fenestration in the staple line. The anvil portion of
the OrVil is oriented so that the rounded portion is placed against the
roof of the mouth. The OrVil is guided into the hypopharynx by pulling
on the the tube end of the device. As the anvil approaches the
hypopharynx, the jaw is pulled forward to allow passage of the anvil.

The shaft of the anvil is brought through the esophageal staple line,
and the tube disconnected from the anvil by cutting the blue sutures.

The superior end of the conduit is opened along the staple line and the
DST XL stapler (matching the diameter of the OrVil) introduced through
the Alexis device (Port \#2). The stapler shaft is placed into the open
end of the gastric conduit and the conduit pulled over it (`sock over
shoe'). The stapler spike is brought out through the greater curvature.
The anvil is grasped with a Maryland grasper placed through a superior
5mm port and the two components of the stapler are mated and the stapler
tightened and fired. The knob of the stapler is rotated two turns
counter-clockwise until a click is felt, at which time the anvil will
flip. The stapler is withdrawn and the donuts examined and sent for
pathologic exam.

The anastomosis is completed by firing a 60mm Blue load across the
conduit cephalad to the anastomosis. The excess conduit is retrieved
through the Alexis port and sent for pathologic exam as `additional
gastric margin.'

\section{Circular stapled anastomosis -
Luketich}\label{circular-stapled-anastomosis---luketich}

May be helpful to move camera back to Port \#3 to see anastomosis
better. May need to remove Port \#1 for better visibility cephalad.

Esophagus divided with robotic shears.

Pursestring suture of 0 Prolene SH to depth of 3-5mm. Two robotic arms
used to stretch esophagus laterally \(\rightarrow\) anvil introduced
into esophagus with large needle driver \(\rightarrow\) pushed
superiorly \(\rightarrow\) oriented cephalad-caudad. Pursestring tied.
Second pursestring of 0 Prolene placed inside of first.

DST XL 28mm EEA stapler introduced through port \#2 (???) into open end
of conduit \(\rightarrow\) mated and fired. Anastomosis completed with
linear stapler. Distance between anatomosis and final linear staple line
should be 2cm.

\section{Linear Stapled Anastomosis}\label{linear-stapled-anastomosis}

The esophageal stump is opened on the (right side) within the middle
with cold scissors or monopolar cautery hook just proximal to the
transecting stapler line

The site for linear stapler insertion is measured to be at least 6 cm
distal to the transecting staple line on the gastric conduit. The
incision is made using monopolar cautery close to the gastroepiploic
arcade

Four retraction sutures are placed (dorsal, ventral, medial and lateral)
around the opened esophageal lumen, making sure to fixate the mucosa to
the rest of the esophageal wall. These retraction sutures will serve for
applying optimal retraction on the stump for anastomotic suturing as
well as help guide the esophageal stump over the stapler anvil.

The robotic stapler anvil and cartridge jaw (45mm? 60mm? length) are
inserted into the appropriate lumens while the console surgeon guides
the tissue onto the stapler. Once sufficient tissue overlap is achieved
(approximately 60 mm) the stapler is fired (blue reload) and the gastric
conduit and esophageal stump are connected dorsally side-to-side

34Fr gastric tube (?or Bougie) inserted. Anterior wall closed with two
4-0 Stratafix sutures run from each side toward the middle. Second later
with 4-0 PDS

\subsection{Omental Flap}\label{omental-flap}

Omental flap is positioned between the conduit and airway and the
anastomosis is loosely wrapped.

\subsection{NG tube}\label{ng-tube}

NG (Salem Sump silicone 18Fr tube) is passed by anesthetist. Robotic
ultrasound used to confirm location. NG tube is advanced to 4th dot just
about to enter the nose.

\section{Drains}\label{drains}

38Fr Chest tube (Blake or Hollow-bore) placed. 19Fr Blake (JP2)
postioned posterior to anastomosis.

\section{Closure}\label{closure-1}

Extraction site closed with 0 Vicryl. Chest tube secured with 0 silk on
SH needle. Skin incisions closed with 4-0 Monocyrl and Dermabond

Grimminger PP, Hadzijusufovic E, Lang H: Robotic-Assisted Ivor Lewis
Esophagectomy (RAMIE) with a Standardized Intrathoracic Circular
End-to-side Stapled Anastomosis and a Team of Two (Surgeon and Assistant
Only). Thorac Cardiovasc Surg 66:404406, 2018

van der Sluis Ann Surg Oncol22 Supp3:S1350-6 - ICG for RGEA

\section{Postoperative care}\label{postoperative-care}

Admission to ICU. Epidural dosed.

Medications:

\begin{itemize}
\tightlist
\item
  Protonix 40mg IV
\item
  Metoprolol 2.5mg IV q6h (may increase to 5mg)
\item
  Home antihypertensives held unless bp is still elevated after adequate
  beta blockade
\item
  No bowel regimen
\item
  Reglan 10mg IV q6 hours
\end{itemize}

Labs:

\begin{itemize}
\tightlist
\item
  Daily drain amylase from JP2 starting POD\#2
\item
  CBC every other day or daily if elevated above 12
\item
  BMP every other day
\item
  CMP on admission and POD \#4
\end{itemize}

Tube feeds at 30mL/hour start with Osmolite 1.5. Tube feeds advanced
10mL q12 hours once flatus returns.

Intravenous fluid stopped once tube feeds advanced and free water
started 240mL qid via jejunostomy.

POD \#1 Out of bed to chair

Chest tube to water seal on POD2 if no air leak. Chest tube removed once
output below 200mL/day

JP2 removed by POD\#7 if patient is otherwise doing well.

Upper GI POD\#2 \emph{AND} able to stand for 10min \emph{AND} output
less than 400mL/day. If good emptying on upper GI \(\rightarrow\) NG
removed.

Modified Barium Swallow once NG tube removed. Initiation of liquids
dependent upon MBS.

POD \#5 Epidural catheter removed if not done prior. Foley catheter
removed once epidural out. No Flomax to be administered vai jejunostomy

\section{References}\label{references}

\chapter{Esophagectomy Robot/VATS}\label{esophagectomy-robotvats}

\section{Preoperative evaluation}\label{preoperative-evaluation-1}

PET scan evaluated for extent of tumor, focusing on proximity to the
carina.

CT scan evaluated for presence of replaced left hepatic artery.

\section{Preop Area}\label{preop-area-1}

\begin{itemize}
\tightlist
\item
  Heparin is held if epidural catheter is planned
\item
  Type/Screen
\item
  Clip hair from abdomen and right chest, including right axilla. Hair
  is clipped from bilateral groins
\end{itemize}

\section{Room Prep}\label{room-prep-2}

\begin{itemize}
\tightlist
\item
  Double-decker Back Table
\item
  EGD Cart (check w/ Salo regarding adult vs neonatal scope)
\item
  BK Ultrasound with sterile ultrasound jelly

  \begin{itemize}
  \tightlist
  \item
    Laparoscopic probe AND
  \item
    hockey stick probe AND
  \item
    robotic probe
  \end{itemize}
\end{itemize}

Walter `Long Arm' retractor

Salo Positioning Cart:

\begin{itemize}
\tightlist
\item
  Three black side-rail clamps
\item
  Two rectangular lateral positioners
\item
  One round lateral positioner
\item
  Yassargil socket
\item
  Well-leg holder (used as an arm holder)
\end{itemize}

Robotic Instruments:

\begin{itemize}
\tightlist
\item
  Tip Up Grasper
\item
  Fenestrated bipolar Forceps (420205)
\item
  Hot Shears (Monopolar scissors) 470179
\item
  Extend vessel sealer (480322)
\item
  SutureCut Large Needle Driver (470296)
\item
  Medium-Large Clip Applier (420327)
\item
  Large Needle Driver (470006)
\item
  Permanent Cautery Hook (470183) (hold)
\item
  Cadiere Forceps (470049) hold
\end{itemize}

Critical supplies (check prior to pt in room)

\begin{itemize}
\tightlist
\item
  Stapler 25mm DSTXL
\item
  Orvil 25mm
\item
  Stapler 21mm DSTXL
\item
  Orvil 21mm
\end{itemize}

Suture

\begin{itemize}
\tightlist
\item
  Stratafix Spiral PDS 3-0 RB-1 needle 15cm
\end{itemize}

Supplies:

\begin{itemize}
\tightlist
\item
  Micropuncture kit (from Anesthesia Workroom)
\item
  Gel Port available
\item
  14Fr Balloon Jejunostomy tube (8200-14)
\item
  14Fr Cook introducer kit
\item
  18Fr Dilator and peelaway sheath
\end{itemize}

\newpage

\textbf{Anesthesia}

\begin{itemize}
\tightlist
\item
  Epidural placed in preop area
\item
  Dual-lumen ETT tube (taped to left)
\item
  Arterial line (left arm will be on arm board)
\item
  Heparin 5000U SQ 1-2 hours after epidural placed
\end{itemize}

\section{Abdominal Position}\label{abdominal-position-1}

Supine on blue foam pad

Foley catheter

Bovie pad

Lower-body Bair Hugger at level of thighs, followed by single blanket
and Velcro strap

Hair clipped from abdomen, bilateral groins, right chest, and right
axilla.

Pre-existing jejunostomy

\begin{itemize}
\tightlist
\item
  Prep into field
\item
  Remove any eschar at exit site
\item
  Secure to skin inferiorly with 0 silk
\end{itemize}

\textbf{Prep}

Chloroprep (2 sticks) of abdomen and lower chest up to the level of
nipples (in case urgent chest tube needed). Particular attention to
prepping as far as possible to the left lateral side and the right
lateral side. Prep continues inferior to base of pubis. At the end of
prep, bilateral groins are prepped.

\textbf{Drape}

Transverse 1000 drape is placed at the base of the pubis. On either
side, the drape proceeds inferior and lateral to include both groins.
Legs are returned to supine position and Velcro strap placed at thighs.
Top gloves are changed.

Blue utility (paper) drape is placed transversely just above level of
pubis. While the abdominal skin incisions are being closed, ICG will be
injected into bilateral groins.

1000 Drapes x4 in rectangular fashion with top of field 10cm above
xiphoid and bottom at base of penis. Ioban strips (4'') aound the
periphery of the surgical field after the trauma drape. Laparoscopic
cords (CO2, light cord, camera) passed off. Suction irrigator brought
off field.

\subsection{Time Out}\label{time-out-2}

Review:

\begin{itemize}
\tightlist
\item
  Consent
\item
  Comorbidities, including cardiopulmonary disease
\item
  Anticoagulation plan. If preoperative heparin was held for epidural,
  plans for timing of its administration. Circulator and anesthetist set
  alarms.
\item
  Prior abdominal or thoracic operations
\item
  Allergies
\item
  Blood

  \begin{itemize}
  \tightlist
  \item
    Surgeon's expectation of blood loss
  \item
    Availability of blood (type/screen vs type/cross).
  \item
    Blood product administration limitations, if any
  \end{itemize}
\item
  Introduction by name and role
\item
  Co-surgeons: Confirm availability prior to incision
\item
  Operative plan
\item
  Request ICU bed in STICU
\end{itemize}

\emph{Comorbidities:}

Cardiopulmonary disease. If echo, ejection fraction and aortic valve
area (if abnormal)

Beta blockade: Note whether patient on home beta blockade and if so,
whether home medication was taken the morning of surgery.

Anticipated Intraop Problems:

\begin{itemize}
\tightlist
\item
  Possibility of tension left pneumothorax due to carbon dioxide entry
  into left chest during mediastinal dissection
\item
  Expectation of ventilatory difficulties due to carbon dioxide entry
  into the right chest
\end{itemize}

\section{Laparoscopy}\label{laparoscopy-1}

Pneumoperitoneum established with Veress needle at left subcostal or
inferior to umbilicus \emph{or} modified Hasson approach. Abdominal
entry with 5mm optical port at future location of Port \#2 in
mid-clavicular line superior to umbilicus. Laparoscopy to look for
peritoneal implants.

Ports:

\begin{enumerate}
\def\labelenumi{\arabic{enumi})}
\tightlist
\item
  Z-thread optical port at site for Port \#1 in R mid-axillary line 5cm
  superior to umbilicus
\item
  RLQ assisting port (5mm)
\item
  Lower midline assisting port (later will become Port 2A)
\item
  LLQ assisting port (5mm)
\end{enumerate}

Site for jejunostomy is selected 25cm from Treitz. 9'' 2-0 silk suture
is placed on anti-mesenteric aspect of jejunum and tied to itself.
Abdominal wall site for jejunostomy is selected to be superior to ports
3 and 4. Site is marked with skin marker

Falciform suspended with 0 silk suture on Keith needle

Once site for jejunostomy marked, will placed ports 2,3,4. Port 4 is
lateral at level of umbilicus. Port \#3 is somewhat more inferior, to
allow sewing jejunostomy. Port \#2 in midline, inferior enough to allow

\section{Robotic ports}\label{robotic-ports-1}

Port \#1 left anterior axillary line 8mm port 2cm above umbilicus.
\(\rightarrow\) fenestrated bipolar

5mm asstant ports RLQ and LLQ

Assistant port 8mm port inferior to umbilicus. Port \#3 will be switched
to this position for jejunostomy later in the case.

(locate future site for jejunostomy)

Port \#3 8mm left mid-axillary line at level of umbilicus
\(\rightarrow\) camera.

Port \#4 left anterior axillary line. 8mm port \(\rightarrow\) Tip Up

Port \#2 right mid-clavicular line 2cm above umbilicus. Switch from 5mm
to 12mm robotic stapler port. Place `sleeper' 0 Vicryl figure-of-eight
suture and secure with a hemostat.

(suspend falciform with 0 silk on Keith needle)

Nathanson liver retractor placed in midline using the 5mm optical port
to create a tract. Nathanson held in place with Endoscopic Bookwalter
secured to bedrail just inferior to right armboard.

Place loop of 0 silk to suspend superior side of pylorus

Bed positioned to 12 degrees reverse Trendelenberg and tipped to right
10 degrees

\section{Abdominal Phase}\label{abdominal-phase-1}

Robot docked: Anatomy: UPPER ABDOMINAL. Position: PATIENT RIGHT

\begin{longtable}[]{@{}
  >{\raggedright\arraybackslash}p{(\linewidth - 6\tabcolsep) * \real{0.2973}}
  >{\raggedright\arraybackslash}p{(\linewidth - 6\tabcolsep) * \real{0.2162}}
  >{\raggedright\arraybackslash}p{(\linewidth - 6\tabcolsep) * \real{0.2432}}
  >{\raggedright\arraybackslash}p{(\linewidth - 6\tabcolsep) * \real{0.2162}}@{}}
\toprule\noalign{}
\begin{minipage}[b]{\linewidth}\raggedright
Arm 1
\end{minipage} & \begin{minipage}[b]{\linewidth}\raggedright
Arm 2
\end{minipage} & \begin{minipage}[b]{\linewidth}\raggedright
Arm 3
\end{minipage} & \begin{minipage}[b]{\linewidth}\raggedright
Arm 4
\end{minipage} \\
\midrule\noalign{}
\endhead
\bottomrule\noalign{}
\endlastfoot
Fenestrated Bipolar & Camera &
\begin{minipage}[t]{\linewidth}\raggedright
Synchroseal OR\\
Vessel Sealer\strut
\end{minipage} & Tip Up \\
\end{longtable}

\subsection{Distal Greater Curve}\label{distal-greater-curve-1}

The course of the right gastroepiploic artery is identified and the
lesser sac entered to the left of midline. Dissection proceeds to the
patient's right. Right gastroepiploic artery carefully preserved and
gatrocolic omentum incised. Dissection proceeds to right to mobilize
colon off duodenum. Kocher maneuver performed and omental adhesions to
gallbladder taken down.

\section{Posterior Stomach}\label{posterior-stomach}

Posterior dissection elevating stomach. Divide posterior gastric
arteries. Dissection plane up to fundus and to short gastrics.

\subsection{Lesser Sac Dissection}\label{lesser-sac-dissection-1}

Gastrohpeatic ligament divided. Right gastric artery divided at 3rd
arcade (approx 5cm from pylorus). Celiac axis evaluated for
resectability. Right crus dissected to ensure resectability. Dissection
carried anteriorly and left crural dissection begun. Dissection
posterior to esophagus.

Penrose can be looped around body of stomach and used to retract stomach
off central lymph nodes by traction to LLQ port site.

Central node dissection begun, including common hepatic, left gastric,
celiac, and proximal splenic artery nodes. If replaced left hepatic
artery present, may require preservation of left gastric artery. Left
gastric artery divided with clips.

Posterior gastric dissection. Dissection is carried posterior towards
most superior short gastric arteries (as these can be difficult to
approach from the greater curvature side). (Penrose drain can be placed
at this part of the operation)

\section{Proximal Greater Curve}\label{proximal-greater-curve-1}

\begin{longtable}[]{@{}
  >{\raggedright\arraybackslash}p{(\linewidth - 6\tabcolsep) * \real{0.2083}}
  >{\raggedright\arraybackslash}p{(\linewidth - 6\tabcolsep) * \real{0.3056}}
  >{\raggedright\arraybackslash}p{(\linewidth - 6\tabcolsep) * \real{0.2083}}
  >{\raggedright\arraybackslash}p{(\linewidth - 6\tabcolsep) * \real{0.2222}}@{}}
\toprule\noalign{}
\begin{minipage}[b]{\linewidth}\raggedright
Arm 1
\end{minipage} & \begin{minipage}[b]{\linewidth}\raggedright
Arm 2
\end{minipage} & \begin{minipage}[b]{\linewidth}\raggedright
Arm 3
\end{minipage} & \begin{minipage}[b]{\linewidth}\raggedright
Arm 4
\end{minipage} \\
\midrule\noalign{}
\endhead
\bottomrule\noalign{}
\endlastfoot
Tip Up & Fenestrated Bipolar & Camera &
\begin{minipage}[t]{\linewidth}\raggedright
Synchroseal\\
Vessel Sealer\strut
\end{minipage} \\
\end{longtable}

Dissection proceeds proximally, dividing the left gastroepiploic artery
near its orgin. Posterior dissection from the pancreas and splenic
artery is performed, followed by division of short gastric vessels

\section{Mediastinal dissection}\label{mediastinal-dissection-2}

Penrose drain is placed around esophagus and secured with Weck clips.
Dissection proceeds to level of inferior pulmonary veins. Penrose is
tucked into the mediastinum. When Penrose is placed, confirm that
heparin has been given.

\section{Pyloric Drainage -
Pyloromyotomy}\label{pyloric-drainage---pyloromyotomy}

Stay sutures are placed superior and inferior with 2-0 Silk. The
cephalad suture is looped with the prior loop of 0 silk placed at the
Nathanson site

\begin{longtable}[]{@{}
  >{\raggedright\arraybackslash}p{(\linewidth - 6\tabcolsep) * \real{0.2973}}
  >{\raggedright\arraybackslash}p{(\linewidth - 6\tabcolsep) * \real{0.2162}}
  >{\raggedright\arraybackslash}p{(\linewidth - 6\tabcolsep) * \real{0.2432}}
  >{\raggedright\arraybackslash}p{(\linewidth - 6\tabcolsep) * \real{0.2162}}@{}}
\toprule\noalign{}
\begin{minipage}[b]{\linewidth}\raggedright
Arm 1
\end{minipage} & \begin{minipage}[b]{\linewidth}\raggedright
Arm 2
\end{minipage} & \begin{minipage}[b]{\linewidth}\raggedright
Arm 3
\end{minipage} & \begin{minipage}[b]{\linewidth}\raggedright
Arm 4
\end{minipage} \\
\midrule\noalign{}
\endhead
\bottomrule\noalign{}
\endlastfoot
Fenestrated Bipolar & Camera &
\begin{minipage}[t]{\linewidth}\raggedright
Scissors OR\\
Cadiere\strut
\end{minipage} & Tip Up \\
\end{longtable}

Pyloromyotomy is performed with camera in Port \#2 and scissors in Port
\#3 and Tip up in Port \#4 (may need to move camera to midline)

\section{Pyloric Drainage - Pyloroplasty
(optional)}\label{pyloric-drainage---pyloroplasty-optional-1}

Pyloric incision is made with scossors or hook cauterh. Defect is closed
with 2-0 silk cut to 6'' (have 2 sutures ready)

\section{Creation of Conduit}\label{creation-of-conduit-1}

Tip up grasper placed into Port \#4. Robotic stapler is introduced into
Port 1. Staple line starts \textasciitilde5cm from pylorus. First firing
with Green (4.8mm) stapler load followed by Blue loads. Conduit is
created with goal of 4cm. Assistant retracts cut edge of specimen
towards right upper quadrant. Vicryl suture is placed 2cm above the
start of the staple line (approximately 7cm from pylorus).

\section{Preparation for ICG}\label{preparation-for-icg-1}

One vial of ICG is resuspended with 10mL of water. 5mL of the solution
is brought onto the sterile field in a 5mL syringe, which is fitted with
a 25Ga needle. The remainder of the ICG is given to Anesthesia for later
(intravenous) use.

\section{Final dissection of
mediastinum}\label{final-dissection-of-mediastinum}

Once conduit has been disconnected from specimen, additional dissection
of mediastinum cam be performed. A penrose drain is tied around the
esopahgus. Once this is finished the penrose drain is pushed up into the
mediastinum with the tails to the patient's right.

\section{Preparation of Conduit}\label{preparation-of-conduit-2}

\begin{longtable}[]{@{}
  >{\raggedright\arraybackslash}p{(\linewidth - 6\tabcolsep) * \real{0.2500}}
  >{\raggedright\arraybackslash}p{(\linewidth - 6\tabcolsep) * \real{0.1591}}
  >{\raggedright\arraybackslash}p{(\linewidth - 6\tabcolsep) * \real{0.4091}}
  >{\raggedright\arraybackslash}p{(\linewidth - 6\tabcolsep) * \real{0.1591}}@{}}
\toprule\noalign{}
\begin{minipage}[b]{\linewidth}\raggedright
Arm 1
\end{minipage} & \begin{minipage}[b]{\linewidth}\raggedright
Arm 2A
\end{minipage} & \begin{minipage}[b]{\linewidth}\raggedright
Arm 3
\end{minipage} & \begin{minipage}[b]{\linewidth}\raggedright
Arm 4
\end{minipage} \\
\midrule\noalign{}
\endhead
\bottomrule\noalign{}
\endlastfoot
Fenestrated Bipolar & Camera & \textbf{SutureCut Large Needle Driver} &
Tip Up \\
\end{longtable}

If there is entry into the left pleaura, JP1 (19Fr) is brought in
through LLQ 5mm port. This is frequently not necessary. Tip Up retracts
left crus laterally and drain is placed in left pleura. All of the
fluted portion is placed into the chest. The dot is left at the hiatus.

Conduit is sutured to specimen with multiple sutures of 2-0 silk (cut to
6''). After conduit is sutured, arm 4 is removed and JP2 brought in
through that port. JP2 is now sutured to the specimen.

Conduit is now evaluated with ICG and with robotic ultrasound

\section{Jejunostomy}\label{jejunostomy-1}

\begin{longtable}[]{@{}
  >{\raggedright\arraybackslash}p{(\linewidth - 6\tabcolsep) * \real{0.2418}}
  >{\raggedright\arraybackslash}p{(\linewidth - 6\tabcolsep) * \real{0.1868}}
  >{\raggedright\arraybackslash}p{(\linewidth - 6\tabcolsep) * \real{0.3956}}
  >{\raggedright\arraybackslash}p{(\linewidth - 6\tabcolsep) * \real{0.1538}}@{}}
\toprule\noalign{}
\begin{minipage}[b]{\linewidth}\raggedright
Arm 1
\end{minipage} & \begin{minipage}[b]{\linewidth}\raggedright
Arm 2
\end{minipage} & \begin{minipage}[b]{\linewidth}\raggedright
Arm 2A
\end{minipage} & \begin{minipage}[b]{\linewidth}\raggedright
Arm 3
\end{minipage} \\
\midrule\noalign{}
\endhead
\bottomrule\noalign{}
\endlastfoot
Fenestrated Bipolar & Camera (30 up) & \textbf{SutureCut Large Needle
Driver} & Tip Up \\
\end{longtable}

Prior suture on jejunum is identified and grasped with suture passer at
the pre-marked site. Two passes are made with the suture passer similar
to transfascial sutures on a Stoppa hernia repair. The 22Ga needle from
a \textbf{micropuncture kit} is introduced through the pre-marked site
to extend 1cm into the abdominal cavity. 3-0 Stratafix on RB-1 is
introduced with a non-cutting needle driver into Port 3. Insufflation
pressure reduced to 8m Hg. Jejunum is sutured to abdominal wall starting
at 10:00. Pursestring is run to 2:00 position and the J wire from the
micropuncture kit threaded into the small bowel. Correct placement
within the bowel lumen is confirmed by rapid insufflation of air through
the needle. The micropuncture dilator is next used, followed by the J
wire from a 14Fr Cook introducer kit. The tract is dilate with a 14FR
dilator followed by 18Fr dilator introduced and left in place.
Pursestring continues around proposed site and then proceeds cephalad to
provide a broader area of fixation between bowel and abdominal wall.
Tract is dilated with 18Fr dilator. 14Fr MIC jejunostomy 8200-14
(trimmed to 10cm) is introduced through 18Fr sheath and positioned
within the bowel. 5mL of water in a slip-tip (non-Luer) 10ml syringe
(from the Cook introducer kit) are instilled into the jejunostomy
balloon.

\section{Drains}\label{drains-1}

Once jejunostomy in place, JP1 and JP2 are routed lateral to jejunostomy
and brought out through Port \#4 and the LLQ assisting port. Enough
slack is left on JP2 so that its outside length approximates that of
JP1.

\section{Closure}\label{closure-2}

Robot is disengaged. Laparoscopic cords are placed in a basin on a ring
stand to be reused for the thoracic portion.

Abdomen is desufflated and ports removed. Port \#2 site is closed by
tying 0 Vicryl suture. Port sites are closed with staples. A Red Rubber
Robinson catheter (or 8Fr feeding tube) is placed through the Nathanson
subxiphoid access site and covered with a Tegaderm. The RRR is used in
case it is necessary to observe the transposition of the conduit from
the abdomen.

\section{ICG injection into groin.}\label{icg-injection-into-groin.-1}

While port sites are being closed, ``loincloth'' is removed and hockey
stick or small T ultrasound transducer is brought onto the field. Under
ultrasound guidance with sterile ultrasound jelly, ICG is injected into
groin nodes.

\section{Thoracic Phase}\label{thoracic-phase-1}

\section{Transition to Corkscrew
Position}\label{transition-to-corkscrew-position}

Patient's chest is rotated to the left. Rectangular positioner supports
scapula. Circular positioner on Left side supports chest. Rectangular
positioner along left greater trochanter. Right arm rests on pink foam
pad on Well Leg Holder. Lower and upper body bair hugger used. Right arm
is elevated for prep and lowered once drapes in place.

\section{Prep and drape}\label{prep-and-drape-1}

Chloroprep including axilla and proximal right arm. Prep should include
all of abdomen. Drape axilla/upper arm with 1000 drape.

Two blue ``U'' drapes used, folowed by trauma drape (or Universal). JP
drains are draped out of the abdominal field but ports 1,2,2A, AND 3 are
prepped into teh field. Ioban strips around edges. Once Ioban in place,
right arm is lowered

Laparoscopic instruments are kept sterile and may be returned to the
field if pleural adhesions are encountered.

\subsection{Thoracoscopy}\label{thoracoscopy}

Table is tipped 10-15 degrees to the left. Gas pressure turned down to
8mmHg in preparation for the thoracic phase. A site for entry in to the
left chest is selected just posterior to the tip of the scapula OR
lateral to nipple. An incision is made here and a 5mm optical port with
a 5mm 0 degree scope used to enter the chest. The chest is insufflated
at 8mmHg of carbon dioxide, which helps to both copllapase the lung and
depress the diaphragm. One 12mm port is placed just lateral to the
diaphragmatic reflection. It is ciritcal to avoid injury to the
diaphragm (and liver) with the inferior port placement. A
mini-thoracotomy incision is placed along the mid-axillary line,
frequently in the same interspace as the inferior/anterior 12mm port.
The chest is entered just superior to the rib and the intercostal
muscles divided with the LigaSure device to allow the ribs to separate.
A narrow Deaver retractor is used to guage the space between the ribs,
as the width of the retractor approximates the diameter of the 25mm
stapler. A 5mm `U' port is generally placed as high as possible midway
between the scapular tip port and the anterior/superior 12mm port t0o
the 4 right 5

\section{Thoracic Dissection}\label{thoracic-dissection-1}

The 1688 ``Green'' scope is used to identify the course of the thoracic
duct. If a disruption is found, the duct is mass ligated as low as
possible with 3-0 Stratafix on RB-1 needle.

\section{Delivery of Conduit into
chest}\label{delivery-of-conduit-into-chest-1}

The conduit is delivered into the thorax. This is done using a
laparoscopic grasper. Ensure staple line faces you (right side). The
specimen is removed through the assistant port incision. Stomach is
pulled into the chest until the \textbf{Vicryl marker suture} is seen.

Division of esophagus just above level of azygous vein. ICG can be
administered intravenously to check perfusion. If there is any question
about the length of the conduit, confirm the length of the conduit prior
to division of the esophagus.

Penrose drain is recovered. Specimen is sent for frozen section with
instructions to freeze proximal margin.

\section{Abdominal Exposure
(optional)}\label{abdominal-exposure-optional}

If needed to monitor passage of conduit into chest, the abdomen is
insufflated using the RRR placed in the sub-xiphoid port. Ports \#1,
\#2, \#3 are re-opened and Nathanson retractor replaced. Robotic camera
used to monitor passage of conduit into chest.

\section{Preparation of conduit}\label{preparation-of-conduit-3}

The gastric conduit is placed into the mediastinum by tucking it
medially from the right pleura into the posterior mediastinum, in order
to allow the conduit to take the most direct path from the hiatus to the
proximal esophagus. Gentle superior tension is now applied to the
conduit in order to eliminate redundancy.

\section{Circular stapled anastomosis -
Orvil}\label{circular-stapled-anastomosis---orvil-1}

OrVil The OrVil device is now used to place a EEA anvil into the distal
esophagus. In general, a 25mm size is selected, unless the patient has
particularly small frame, in which case a 21mm size is used. Two stay
sutures of 2-0 silk on RB-1 needles are placed in the center of the
esophageal staple line 3mm apart. A Harmonic scalpel is used to divide
the staple line. The OrVil device is passed through the mouth and is
passed through the fenestration in the staple line. The anvil portion of
the OrVil is oriented so that the rounded portion is placed against the
roof of the mouth. The OrVil is guided into the hypopharynx by pulling
on the the tube end of the device. As the anvil approaches the
hypopharynx, the jaw is pulled forward to allow passage of the anvil.

The shaft of the anvil is brought through the esophageal staple line,
and the tube disconnected from the anvil by cutting the blue sutures.

The superior end of the conduit is opened along the staple line and the
DST XL stapler (matching the diameter of the OrVil) introduced through
the Alexis device (Port \#2). The stapler shaft is placed into the open
end of the gastric conduit and the conduit pulled over it (`sock over
shoe'). The stapler spike is brought out through the greater curvature.
The anvil is grasped with a Maryland grasper placed through a superior
5mm port and the two components of the stapler are mated and the stapler
tightened and fired. The knob of the stapler is rotated two turns
counter-clockwise until a click is felt, at which time the anvil will
flip. The stapler is withdrawn and the donuts examined and sent for
pathologic exam.

The anastomosis is completed by firing a 60mm Gold load across the
conduit cephalad to the anastomosis. The excess conduit is retrieved
through the Alexis port and sent for pathologic exam as `additional
gastric margin.'

\section{NG Tube}\label{ng-tube-1}

A Covidien Salem Sump 18Fr nasogastric tube is passed by the
anesthetist. A laparoscopic BK ultrasound is used to monitor the passage
of the NG tube through the esophagus and into the gastric conduit. The
NG tube is passed to the level that all for dots are outside the nose,
with the 4th dot at the nares. The NG tube is secured with an AMT
bridle.

\section{Chest Tube}\label{chest-tube}

A 28Fr Blake chest tube is placed through the anterior/inferior 12mm
port and is positioned into the posterior mediastinum. JP2 is placed
near the gastric conduit. The right lung is re-inflated.

\section{Closure}\label{closure-3}

The stapler access port incision is closed with 0 Vicryl to approximate
the serratus muscle. The incisions are closed with 4-0 Monocryl followed
by Dermabond.

\section{Postoperative care}\label{postoperative-care-1}

Admission to ICU. Epidural dosed.

Medications:

\begin{itemize}
\tightlist
\item
  Protonix 40mg IV daily
\item
  Metoprolol 2.5mg IV q6h (may increase to 5mg)
\item
  Home antihypertensives held unless bp is still elevated after adequate
  beta blockade
\item
  No bowel regimen
\item
  Reglan 10mg IV q6 hours
\end{itemize}

Labs:

\begin{itemize}
\tightlist
\item
  Daily drain amylase from JP2 starting POD\#2
\item
  CBC every other day or daily if elevated above 12
\item
  BMP every other day
\item
  CMP on admission and POD \#4
\end{itemize}

Tube feeds at 30mL/hour start with Osmolite 1.5. Tube feeds advanced
10mL q12 hours once flatus returns.

Intravenous fluid stopped once tube feeds advanced and free water
started 240mL qid via jejunostomy.

POD \#1 Out of bed to chair

Chest tube to water seal on POD2 if no air leak. Chest tube removed once
output below 200mL/day

JP2 removed by POD\#7 if patient is otherwise doing well.

Upper GI POD\#2 \emph{AND} able to stand for 10min \emph{AND} output
less than 400mL/day. If good emptying on upper GI \(\rightarrow\) NG
removed.

Modified Barium Swallow once NG tube removed. Initiation of liquids
dependent upon MBS.

POD \#5 Epidural catheter removed if not done prior. Foley catheter
removed once epidural out. No Flomax to be administered via jejunostomy

\part{Postop Care}

\chapter{Colectomy Postop}\label{colectomy-postop}

\textbf{Clinic}

\begin{itemize}
\tightlist
\item
  Opiod Cessation
\item
  Smoking Cessation
\item
  EtOH/Drugs of Abuse -- Social Work
\item
  Nutritonal Evaluation

  \begin{itemize}
  \tightlist
  \item
    All Patients -- Ensure for 3d preop
  \item
    Poor nutrition -- ?Delay surgery
  \end{itemize}
\item
  Preop Anesthesia/ERAS Class
\item
  Expectations of Surgery:

  \begin{itemize}
  \tightlist
  \item
    Length of Stay 1-3 days
  \item
    Diet (self-limiting)
  \item
    Pain control (low-opoid)
  \item
    Activity (OOB at 6am, OOB 3x/day)
  \end{itemize}
\item
  Bowel Prep -- Abx and mechanical
\end{itemize}

\textbf{Preop Holding}

\begin{itemize}
\tightlist
\item
  Colon PowerPlan (Hill)
\item
  Antibiotic PowerPlan
\item
  No PCN only for severe allergy
\item
  Entereg (if no preop opiods)
\item
  VTE prophylaxis
\item
  Carbohydrate load (2hr preop)
\end{itemize}

\textbf{OR}

\begin{itemize}
\tightlist
\item
  Goal-directed fluid administration
\item
  2L total in OR
\item
  3L total/first 24hrs
\item
  Open procedures: No epidural
\end{itemize}

\textbf{Postop Day 0}

\begin{itemize}
\tightlist
\item
  Goal-directed fluid administration
\item
  OOB (Dangling not compliant)
\item
  Diet

  \begin{itemize}
  \tightlist
  \item
    Low Residue diet
  \item
    Ensure Supplements
  \end{itemize}
\item
  Teaching
\item
  Gum, Mag \& Entereg
\item
  Pain Management

  \begin{itemize}
  \tightlist
  \item
    PCA
  \item
    Tylenol 1gm q6 ATC
  \item
    Gabapentin 300mg TID
  \item
    Tramadol PRN
  \item
    Resume all baseline pain meds
  \end{itemize}
\item
  Home Medications

  \begin{itemize}
  \tightlist
  \item
    Resume all home medications (write on postop 0 so it can be
    administered morning of POD1)
  \item
    No therapeutic anti-coagiation
  \item
    Diabetes medicines

    \begin{itemize}
    \tightlist
    \item
      Prefer to resume all oral diabetes meds
    \item
      Sliding-scale insulin ordered for all diabetics
    \end{itemize}
  \end{itemize}
\end{itemize}

\textbf{Postop Day 1}

\begin{itemize}
\tightlist
\item
  Labs: K+ and CBC only
\item
  Heparin lock IV
\item
  Remove Foley
\item
  d/c PCA (unless open incision)
\item
  Out of bed \textgreater{} 6 hrs
\item
  Diet: low residue
\item
  Pain management

  \begin{enumerate}
  \def\labelenumi{\arabic{enumi})}
  \tightlist
  \item
    d/c PCA
  \item
    Tylenol
  \item
    Gabapentin 300 tid
  \item
    Tramadol
  \item
    Home meds
  \item
    Oxy if lots of pain
  \end{enumerate}
\item
  Afternoon rounds
\end{itemize}

\begin{enumerate}
\def\labelenumi{\arabic{enumi})}
\tightlist
\item
  Check patient 2-4PM
\item
  Ambulation 2x's by PM
\item
  Patient education
\end{enumerate}

\textbf{Postop Day 2}

\begin{itemize}
\tightlist
\item
  No IV fluids unless indicated
\item
  No labs unless indicated
\item
  OOB \textgreater6 hrs
\item
  No PT unless going to rehab or SNF (ask Hill first)
\item
  Pain management

  \begin{itemize}
  \tightlist
  \item
    Same as POD\#1
  \end{itemize}
\item
  Discharge planning

  \begin{itemize}
  \tightlist
  \item
    Consider early D/C (median LOS=2d)
  \end{itemize}
\item
  Otherwise plan

  \begin{itemize}
  \tightlist
  \item
    Patient
  \item
    Nursing
  \end{itemize}
\item
  Afternoon rounds
\end{itemize}

\begin{enumerate}
\def\labelenumi{\arabic{enumi})}
\tightlist
\item
  Possible home today!
\item
  Check patient 2-3PM
\item
  Ambulation 2x's by PM
\item
  Check for dehydration
\item
  Patient education
\item
  Give estimated date of discharge
\end{enumerate}

\textbf{Postop Day 3-5}

\begin{itemize}
\item
  IVF

  \begin{itemize}
  \tightlist
  \item
    Consider bolus of IV fluids
  \item
    Consider maintenance IV if we feel won't resolve soon
  \end{itemize}
\item
  Labs

  \begin{itemize}
  \tightlist
  \item
    No labs unless indicated
  \item
    Consider ordering QOD Chem7 if prolonged ileus
  \item
    OOB \textgreater6 hrs
  \item
    No PT unless likely to go to rehab or SNIF (ask Hill first)
  \end{itemize}
\item
  Pain management

  \begin{itemize}
  \tightlist
  \item
    Same as POD\#1
  \end{itemize}
\item
  Afternoon rounds

  \begin{enumerate}
  \def\labelenumi{\arabic{enumi})}
  \tightlist
  \item
    Possible home today!
  \item
    Check patient 2-3PM
  \item
    Ambulation 2x's by PM
  \item
    Check for dehydration
  \item
    Patient education
  \item
    Give estimated date of discharge
  \end{enumerate}
\end{itemize}

\chapter{LAR + Ileostomy}\label{lar-ileostomy}

\textbf{Clinic}

\begin{itemize}
\tightlist
\item
  Opiod Cessation
\item
  Smoking Cessation
\item
  EtOH/Drugs of Abuse -- Social Work
\item
  Nutritional Evaluation

  \begin{itemize}
  \tightlist
  \item
    All Patients -- Ensure for 3d preop
  \item
    Poor nutrition -- ?Delay surgery
  \end{itemize}
\item
  \emph{Arrange Wound Ostomy Nursing}
\item
  \emph{Pre-approval for Home Health}
\item
  Preop Anesthesia/ERAS Class
\item
  Expectations of Surgery:

  \begin{itemize}
  \tightlist
  \item
    Length of Stay 1-3 days
  \item
    Diet (self-limiting)
  \item
    Pain control (low-opoid)
  \item
    Activity (OOB at 6am, OOB 3x/day)
  \end{itemize}
\item
  Bowel Prep -- Abx and mechanical
\end{itemize}

\textbf{Preop Holding}

\begin{itemize}
\tightlist
\item
  Colon PowerPlan (Hill)
\item
  \emph{CCM order}
\item
  \emph{WOCN order}
\item
  Antibiotic PowerPlan
\item
  No PCN only for severe allergy
\item
  Entereg (if no preop opiods)
\item
  VTE prophylaxis
\item
  Carbohydrate load (2hr preop)
\item
  Confirm stoma marking
\item
  Confirm CCM aware of ostomy
\end{itemize}

\textbf{OR}

\begin{itemize}
\tightlist
\item
  Goal-directed fluid administration

  \begin{itemize}
  \tightlist
  \item
    2L total in OR
  \item
    3L total/first 24hrs
  \end{itemize}
\item
  Open procedures: No epidural
\end{itemize}

\textbf{Postop Day 0}

\begin{itemize}
\tightlist
\item
  Goal-directed fluid administration
\item
  OOB (Dangling not compliant)
\item
  Diet

  \begin{itemize}
  \tightlist
  \item
    Low Residue diet
  \item
    Ensure Supplements
  \end{itemize}
\item
  Teaching
\item
  Gum, Mag \& Entereg
\item
  Pain Management

  \begin{itemize}
  \tightlist
  \item
    PCA
  \item
    Tylenol 1gm q6 ATC
  \item
    Gabapentin 300mg TID
  \item
    Oxycodone once at night
  \item
    Tramadol PRN
  \item
    Resume all baseline pain meds
  \end{itemize}
\item
  Home Medications

  \begin{itemize}
  \tightlist
  \item
    Resume all home medications
  \item
    No therapeutic anti-coagiation
  \item
    Diabetes medicines

    \begin{itemize}
    \tightlist
    \item
      Prefer to resume all
    \item
      Sliding-scale insulin if needed
    \end{itemize}
  \end{itemize}
\item
  \emph{Wound Ostomy teaching}
\item
  \emph{CCM for Home Health for stoma care}
\end{itemize}

\textbf{Postop Day 1}

\begin{itemize}
\tightlist
\item
  Labs: K+ and CBC only
\item
  d/c ``pre'' plan \& colon visit
\item
  Heparin lock IV
\item
  Remove Foley
\item
  d/c PCA (unless laparotomy)
\item
  Out of bed \textgreater{} 6 hrs
\item
  Diet: low residue
\item
  Pain management

  \begin{enumerate}
  \def\labelenumi{\arabic{enumi})}
  \tightlist
  \item
    d/c PCA
  \item
    Tylenol
  \item
    Gabapentin 300 tid
  \item
    Tramadol
  \item
    Home meds
  \item
    Oxy if lots of pain
  \end{enumerate}
\item
  Discharge planning
\end{itemize}

\begin{enumerate}
\def\labelenumi{\arabic{enumi})}
\tightlist
\item
  Possible home today!(25\% will go home POD1)
\item
  Check patient 2-3PM
\end{enumerate}

\begin{itemize}
\tightlist
\item
  Afternoon rounds
\end{itemize}

\begin{enumerate}
\def\labelenumi{\arabic{enumi})}
\tightlist
\item
  Possible home today!
\item
  Check patient 2-4PM
\item
  Ambulation 2x's by PM
\item
  Patient education
\end{enumerate}

\textbf{Postop Day 2}

\begin{itemize}
\tightlist
\item
  No IV fluids unless indicated
\item
  No labs unless indicated
\item
  OOB \textgreater6 hrs
\item
  No PT unless going to rehab or SNF (ask Hill first)
\item
  Pain management

  \begin{itemize}
  \tightlist
  \item
    Same as POD\#1
  \end{itemize}
\item
  \emph{Wound Ostomy Teaching}
\item
  Discharge planning

  \begin{itemize}
  \tightlist
  \item
    Consider early D/C (median LOS=2d)
  \end{itemize}
\item
  Otherwise plan

  \begin{itemize}
  \tightlist
  \item
    Patient
  \item
    Nursing
  \item
    \emph{CCM for home health for stoma care}
  \end{itemize}
\item
  Afternoon rounds
\end{itemize}

\begin{enumerate}
\def\labelenumi{\arabic{enumi})}
\tightlist
\item
  Possible home today!
\item
  Check patient 2-3PM
\item
  Ambulation 2x's by PM
\item
  Check for dehydration
\item
  Patient education
\item
  Give estimated date of discharge
\end{enumerate}

\textbf{Postop Day 3-5} - IVF - Consider bolus of IV fluids - Consider
maintenance IV if we feel won't resolve soon - Labs - No labs unless
indicated - Consider ordering QOD Chem7 if prolonged ileus - OOB
\textgreater6 hrs - No PT unless likely to go to rehab or SNIF (ask Hill
first) - Pain management - Same as POD\#1 - Afternoon rounds 1) Possible
home today! 2) Check patient 2-3PM 3) Ambulation 2x's by PM 4) Check for
dehydration 5) Patient education 6) Give estimated date of discharge

\chapter{Abdominoperineal Resection
Postop}\label{abdominoperineal-resection-postop}

\textbf{Clinic}

\begin{itemize}
\tightlist
\item
  Opiod Cessation
\item
  Smoking Cessation
\item
  EtOH/Drugs of Abuse -- Social Work
\item
  Nutritional Evaluation

  \begin{itemize}
  \tightlist
  \item
    All Patients -- Ensure for 3d preop
  \item
    Poor nutrition -- ?Delay surgery
  \end{itemize}
\item
  \emph{Arrange Wound Ostomy Nursing}
\item
  \emph{Pre-approval for Home Health}
\item
  Preop Anesthesia/ERAS Class
\item
  Expectations of Surgery:

  \begin{itemize}
  \tightlist
  \item
    Length of Stay 1-3 days
  \item
    Diet (self-limiting)
  \item
    Pain control (low-opoid)
  \item
    Activity (OOB at 6am, OOB 3x/day)
  \end{itemize}
\item
  Bowel Prep -- Abx and mechanical
\end{itemize}

\textbf{Preop Holding}

\begin{itemize}
\tightlist
\item
  Colon PowerPlan (Hill)
\item
  \emph{CCM order}
\item
  \emph{WOCN order}
\item
  Antibiotic PowerPlan
\item
  No PCN only for severe allergy
\item
  Entereg (if no preop opiods)
\item
  VTE prophylaxis
\item
  Carbohydrate load (2hr preop)
\item
  Confirm stoma marking
\item
  Confirm CCM aware of ostomy
\end{itemize}

\textbf{OR}

\begin{itemize}
\tightlist
\item
  Goal-directed fluid administration

  \begin{itemize}
  \tightlist
  \item
    2L total in OR
  \item
    3L total/first 24hrs
  \end{itemize}
\item
  Open procedures: No epidural
\end{itemize}

\textbf{Postop Day 0}

\begin{itemize}
\tightlist
\item
  No sitting

  \begin{itemize}
  \tightlist
  \item
    Order Sign over bed ``No sitting''
  \end{itemize}
\item
  Goal-directed fluid administration
\item
  OOB (Dangling not compliant)
\item
  Diet

  \begin{itemize}
  \tightlist
  \item
    Low Residue diet
  \item
    Ensure Supplements
  \end{itemize}
\item
  Teaching
\item
  Gum, Mag \& Entereg
\item
  Pain Management

  \begin{itemize}
  \tightlist
  \item
    PCA
  \item
    Tylenol 1gm q6 ATC
  \item
    Gabapentin 300mg TID
  \item
    Oxycodone once at night
  \item
    Tramadol PRN
  \item
    Resume all baseline pain meds
  \end{itemize}
\item
  Home Medications

  \begin{itemize}
  \tightlist
  \item
    Resume all home medications
  \item
    No therapeutic anti-coagiation
  \item
    Diabetes medicines

    \begin{itemize}
    \tightlist
    \item
      Prefer to resume all
    \item
      Sliding-scale insulin if needed
    \end{itemize}
  \end{itemize}
\item
  \emph{Wound Ostomy teaching}
\item
  \emph{CCM for Home Health for stoma care}
\end{itemize}

\textbf{Postop Day 1}

\begin{itemize}
\tightlist
\item
  Labs: K+ and CBC only
\item
  d/c ``pre'' plan \& colon visit
\item
  Heparin lock IV
\item
  Out of bed \textgreater{} 6 hrs
\item
  Diet: low residue
\item
  Pain management

  \begin{enumerate}
  \def\labelenumi{\arabic{enumi})}
  \tightlist
  \item
    d/c PCA
  \item
    Tylenol
  \item
    Gabapentin 300 tid
  \item
    Tramadol
  \item
    Home meds
  \item
    Oxy if lots of pain
  \end{enumerate}
\end{itemize}

\textbf{Postop Day 2}

\begin{itemize}
\tightlist
\item
  No IV fluids unless indicated
\item
  No labs unless indicated
\item
  OOB \textgreater6 hrs
\item
  No PT unless going to rehab or SNF (ask Hill first)
\item
  Pain management

  \begin{itemize}
  \tightlist
  \item
    d/c PCA (unless laparotomy)
  \end{itemize}
\item
  \emph{Wound Ostomy Teaching}
\item
  Foley voiding challenge

  \begin{itemize}
  \tightlist
  \item
    250mL saline into foley -\textgreater{} pull
  \item
    Must void within 2 hours
  \end{itemize}
\item
  Discharge planning

  \begin{itemize}
  \tightlist
  \item
    Consider early D/C
  \end{itemize}
\item
  Otherwise plan

  \begin{itemize}
  \tightlist
  \item
    Patient
  \item
    Nursing
  \item
    \emph{CCM for home health for stoma care}
  \end{itemize}
\item
  Afternoon rounds
\end{itemize}

\begin{enumerate}
\def\labelenumi{\arabic{enumi})}
\tightlist
\item
  Possible home today!
\item
  Check patient 2-3PM
\item
  Ambulation 2x's by PM
\item
  Check for dehydration
\item
  Patient education
\item
  Give estimated date of discharge
\end{enumerate}

\textbf{Postop Day 3-5} - IVF - Consider bolus of IV fluids - Consider
maintenance IV if we feel won't resolve soon - Labs - No labs unless
indicated - Consider ordering QOD Chem7 if prolonged ileus - OOB
\textgreater6 hrs - No PT unless likely to go to rehab or SNIF (ask Hill
first) - Pain management - Same as POD\#1 - Afternoon rounds 1) Possible
home today! 2) Check patient 2-3PM 3) Ambulation 2x's by PM 4) Check for
dehydration 5) Patient education 6) Give estimated date of discharge

\chapter{Gastrectomy Postop}\label{gastrectomy-postop}

\section{Postoperative PPI}\label{postoperative-ppi}

Proton-pump inhibitors have been shown to reduce risk of marginal
ulceration but also are associated with increased risk of delayed
gastric emptying after pancreaticoduodenectomy and gastric bypass.

There is a limited evidence base for duration of PPIs after gastrectomy.
The American Society for Metabolic and Bariatric surgery (Edwards et al.
2025) recommends at least 90 days of PPIs. A retrospective study found
lower rate of marginal ulcer with 90 days of PPI after gastric bypass
compared with 30 days.

\section{Nasogastric tubes}\label{nasogastric-tubes}

The duration of nasogastric tubes will be determined by the attending
surgeon, operation performed, and NG output. Experience with advanced
recovery protocols would suggest that prolonged nasogastric tube
drainage is not only unnecessary but may slow recovery.

For Dr Salo's patients, an upper GI is generally performed on
postoperative day 2 using Isovue through the nasogastric tube. The
patient stands upright and gastric emptying is evaluated

\section{Drains}\label{drains-2}

Drain amylase is monitored daily starting on the second day after
surgery.

\chapter{Gastrostomy Feeds}\label{gastrostomy-feeds}

Gastrostomy feeds are generally given via bolus feeds 3-4 times per day.
Tube feed goal is usually 1 to 1.5 cartons at a time (4-6 cartons
total). For patients at risk of refeeding, tube feeds are started at one
half carton qid. The amount per feed can be increased over the next few
days to goal.

Patients with severe reflux may require gastrostomy feeds via infusion
pump, but this is much less common. Patients with gastrostomy feeds DO
NOT routinely need an infusion pump initially.

Osmolite 1.5 is generally used as an initial formula. Glucerna is
formulated with carbohydrates with low glycemic index, which is
particularly helpful for bolus feeding (via gastrostomy). However, the
high fiber context of Glucerna may make this formula more prone to
causing clogging of gastrostomy tubes.

\section{Free Water Flushes}\label{free-water-flushes}

Most patients require free water in addition to tube feeds. Most
patients require 240mL 4 times per day via gastrostomy

\section{Med Administration via
G-Tubes}\label{med-administration-via-g-tubes}

Patients who have difficulty with dysphagia or complete esophageal
obstruction will need to have their medicines administered via
jejunostomy tube. The process of designing a medicine regimen which can
safely be administered via enteral tube can be a challenge and may
require a consultation with the hospital pharmacist. One the other hand,
the financial consequences of a clogged feeding tube are substantial.
\textbf{Flomax is never given via enteral feeding tubes due to risk of
clogging}.

Several common medicines are available in liquid form:

\begin{itemize}
\tightlist
\item
  Acetaminophen (pediatric formulation)
\item
  Gabapentin
\item
  Oxycodone
\item
  Hydrocodone + acetaminophen
\item
  Reglan
\end{itemize}

Apart from delayed-release medicines and Flomax (tamsulosin), \emph{most
medicines can be crushed and resuspended and administered via
gastrostomy tube if needed}.

Nexium can be administered by opening the capsule and resuspending the
beads in 50mL of water. The beads are resuspended and administered. In
this case, the beads don't dissolve in the syringe, but can be
administered without risk of clogging.

\chapter{Jejunostomy Feedings}\label{jejunostomy-feedings}

Due to the osmotic load, jejunostomy feedings are given via enteral
(Kangaroo) pump rather than bolus feeding. Feedings are generally begun
as continuous (around-the-clock) and are then transitioned to nocturnal
(generally 6pm to 10am) prior to discharge.

Jejunostomy feedings carry a small but significant risk
(\textasciitilde1\%) of small bowel necrosis, as evidenced by the
findings of pneumatosis on CT scan and in some cases small bowel
necrosis and perforation. The existing literature would suggest that the
early symptoms associated with small bowel necrosis are abdominal
distension. As a result, patients on jejunostomy tube feeding need to be
carefully monitored for distension, and tube feeds held if distension
develops.(Taylor et al. 2014) Tube feeds are generally started at
30mL/hour in the immediate postoperative patients. In patients who are
awake and in whom it is possible to determine whether or not there are
issues of tube feed intolerance, the rate of tube feedings is increased
10mL/hour every 12 hours to a goal of 60mL/hour for women and 75mL/hour
for men. In patients who are intubated/sedated, the advancement of tube
feeds is individualized, and decisions are made on a daily basis on
rounds.

For patients who receive tube feedings preoperatively, the same formula
is generally used after surgery.

For patients with BMI less than 35, Osmolite 1.5 is used as a starting
formula. If Osmolite is not tolerated due to diarrhea, Vital 1.5 can be
used. For patients with BMI greater than 35, Promote is used as a
starting formula, as it contains a lower amount of carbohydrates then
Osmolite 1.5. Patients who are intolerant to Promote can be switched to
Vital High Protein.

Patients less than 150 pounds receive 4 cans of tube feedings,
administered at 60mL/hour x 16 hours (6pm to 10am). Patients greater
than 150 pounds receive 5 cans of tube feedings, administered at
75mL/hour x 16 hours (6pm to 10am). Patients with BMI \textgreater35
receive 5 cans of Promote (or Vital High Protein).

Glucerna is formulated with carbohydrates with low glycemic index, which
is particularly helpful for bolus feeding (via gastrostomy) or patients
who are eating. Glycemic index is less important in patients on
continuous tube feeds (eg via jejunostomy). In addition, the high fiber
context of Glucerna may make this formula more prone to causing clogging
of jejunostomy tubes. Free Water Flushes

Once patients are transferred to the ward, free water via jejunostomy
should be ordered as 240mL via jejunostomy qid. This is entered under
Flush Enteral Tube for Free Water Requirements. Diarrhea with
Jejunostomy Feedings

Patients who experience diarrhea during jejunostomy administration
(especially diarrhea at night) will need this addressed. Several steps

\begin{itemize}
\tightlist
\item
  Send stool for C Diff
\item
  Consider changing tube feedings to a more easily digestible formula.
  Patients on Osmolite 1.5 can be changed to Vital 1.5. Patients on
  Vital 1.5 can be changed to Vivonex. Patients on Promote can be
  changed to Vital High Protein
\item
  Stop tube feedings for 2-4 hours to allow diarrhea to resolve
\item
  restart tube feeds at a lower rate (eg 20mL/hour lower than the prior
  rate).
\item
  Bannatrol can be given to patients taking an oral diet. Bannatrol is
  NOT given via jejunostomy tube to avoid clogging.
\item
  Lomotil is generally used as a last resort.
\end{itemize}

\section{Med Administration via
J-Tubes}\label{med-administration-via-j-tubes}

Patients who have difficulty with dysphagia or complete esophageal
obstruction will need to have their medicines administered via
jejunostomy tube. The process of designing a medicine regimen which can
safely be administered via enteral tube can be a challenge and may
require a consultation with the hospital pharmacist. One the other hand,
the financial consequences of a clogged feeding tube are substantial.
\textbf{Flomax is never given via jejunostomy tubes due to risk of
clogging}.

Several common medicines are available in liquid form:

\begin{itemize}
\tightlist
\item
  Acetaminophen (pediatric formulation)
\item
  Gabapentin
\item
  Oxycodone
\item
  Hydrocodone + acetaminophen
\item
  Reglan
\end{itemize}

Apart from delayed-release medicines and Flomax (tamsulosin), most
medicines can be crushed and resuspended and administered via
gastrostomy tube if needed.

Nexium can be administered by opening the capsule and resuspending the
beads in 50mL of water. The beads are resuspended and administered. In
this case, the beads don't dissolve in the syringe, but can be
administered without risk of clogging.

\section{Occluded Jejunostomy Tubes}\label{occluded-jejunostomy-tubes}

Jejunostomy tubes which are refractory to the usual non-invasive means
(warm water, Coca-Cola) will need to be changed over a wire in
Interventional Radiology.

\section{Jejunostomy feedings on preop
Patients}\label{jejunostomy-feedings-on-preop-patients}

Scope Anesthesia policy states the jejunostomy feeds are considered
equivalent to oral solid food in terms of NPO interval. In general,
jejunostomy tube feeds should be held at midnight the night before
surgery.

\chapter{Jejunostomy + Diabetes}\label{jejunostomy_diabetes}

\textbf{Inpt -- Jejunostomy \textasciitilde{} Diabetes}

Patients on tube feeds are typically started on continuous
(around-the-clock) tube feedings, and subsequently changed to a
nocturnal regimen (typically 6pm to 10am). The diabetic management for
these patients differs depending upon their tube feeding regimen:

\textbf{Diabetics on Continuous Tube Feeds.}

Non-insulin diabetic patients are generally initially given Osmolite 1.5
as it provides higher caloric density and does not contain fiber, which
tends to clog the feeding tubes. Diabetic patients requiring insulin can
be initially trialed on Osmolite 1.5 but are changed to Promote or
Glucerna 1.5 if their blood sugars prove difficult to control (as
evidence by either the need for EndoTool or requiring a q6 hour regimen
of insulin N + insulin R).

Diabetic patients who need insulin while receiving tube feedings are
typically treated initially with continuous tube feedings and
around-the-clock insulin. Patients with large insulin requirements may
need hourly intravenous insulin (with dosages calculated via EndoTool).
Once their insulin requirements are stabilized, they are transitioned to
a q6 hour regimen consisting of N and R insulin, typically twice as many
units of N insulin as R (for instance, 6Units of N + 3 Units of R
insulin every six hours). Patients who require EndoTool or a q6hr
regimen need an endocrinology consultation with Dr Kelli Dunn to assist
in diabetic management.

\textbf{Diabetics on Nocturnal Tube Feeds}

Most patients are transitioned from continuous tube feeds to nocturnal
prior to discharge. Diabetic patients on nocturnal tube feedings
typically receive tube feeding from 6pm to 10am and receive insulin at
initiation of tube feeds (6pm) and again at 6 hours later (midnight). A
typical regimen might be 18U of 70/30 at 1800 and 18U of 70/30 at MN. An
alternative might be 12U NPH + 6U Regular at 1800 and 16U NPH and 8U
Regular insulin at MN. Because it can take several days to determine the
correct insulin regimen, diabetic patients receiving jejunostomy
feedings are cycled as early in their hospital course as possible to
avoid delaying discharge for blood sugar management.

Diabetic patients receiving insulin will need careful coordination of
tube feeding and insulin administration when they are being transitioned
from continuous to nocturnal tube feeds. Patients on continuous tube
feeds may receive insulin on a q6 hour schedule, while those receiving
nocturnal tube feeds receive insulin at 1800 and MN. When patients on
continuous tube feeds are transitioned, the tube feeds are stopped at
10am, to be restarted at 6pm that evening. It is critical that as soon
as the tube feeds are stopped at 10am, that the q6 hour insulin as
stopped as well.

In either case it is critical that if tube feedings are stopped,
standing insulin administration (either q6 hour OR 1800 and MN) be
stopped as well. In these cases, sliding scale insulin is generally
continued.

\begin{longtable}[]{@{}llll@{}}
\toprule\noalign{}
Day & Time & Tube Feeds & Insulin \\
\midrule\noalign{}
\endhead
\bottomrule\noalign{}
\endlastfoot
SUN & MN & 60mL/hr & 8N+4R \\
Mon & 6am & 60mL/hr & 8N+R \\
Mon & Noon & Stop & None \\
Mon & 6pm & 75mL/hr & 16N + 8R \\
Mon & MN & 75mL/hr & 16N+ 8R \\
Tues & 10am & Stop & None \\
Tues & 6pm & 75mL/hr & 16N+ 8R \\
Tues & MN & 75mL/hr & 16N+ 8R \\
Weds & 10am & Stop & None \\
\end{longtable}

In this example, the patient was receiving 8N + 4R every 6 hours, so
total insulin units per day is (8+4) x 4 = 48units. Because the
carbohydrate load of 60mL/hour x 24 is roughly equivalent to 75mL/hour x
16 hours, the total insulin administered is roughly the same. When
converted to nocturnal dosing, the patient now received 16+8 = 24 units
twice (6pm and MN) = 48U. In practice, it may be wiser to begin by
adjusting the dose down a little, to perhaps 14U N and 7U R for the
first night.

\chapter{Esophagectomy ICU}\label{esophagectomy-icu}

See
\href{https://gisurgonc.github.io/talks/eso_postop_care.htm}{Esophagectomy
Postop Care}

\textbf{ICU Care}

All patients admitted to STICU with Surgical Critical Care Consultation.
Contact STICU High Team or STICU Low Team (depending upon bed
assignment)

\textbf{Ward Care}

Once stable, patients are transferred to 6T. If a 6T bed is not
available, please notify Dr Salo. Historically, over 95\% of patients
are transferred to 6T after leaving the ICU.

\textbf{Neuro}

Multimodal pain control:

\begin{itemize}
\tightlist
\item
  gabapentin (300mg tid liquid via Jejunostomy)
\item
  Tylenol (1000mg q6hrs as pediatric liquid via Jejunostomy).
\item
  Epidural in most patients. Epidurals will usually be removed by the
  fifth postoperative day.
\item
  {[}ASRA Guidelines for anticoagulation{]}
  (\url{https://rapm.bmj.com/content/rapm/43/3/263.full.pdf})
\item
  Once epidural removed: PCA with subsequent conversion to oxycodone
  elixir via Jejunostomy.
\item
  No ketorolac (Toradol) given risk of anastomotic
  failure\hspace{0pt}1\hspace{0pt}
\item
  Home anxiolytics are generally administered at half the home dose
\end{itemize}

\textbf{Cardiovascular}

Postoperative atrial fibrillation is a common occurrence (20\%) after
esophagectomy, with the risk increasing in older patients. For patients
over age 70, half of patients will develop atrial fibrillation in the
postoperative period.

For prevention of atrial fibrillation, beta blockade is used. For
patients receiving beta blockers prior to surgery, continuation of beta
blockade is recommended. For others, patients are given metoprolol 2.5
mg IV q6hrs which can be titrated up to 10mg IV q6hrs as needed. See STS
Guidelines.\hspace{0pt}2\hspace{0pt}

Home anti-hypertensives are usually held in order to allow
beta-blockade. Patients who have elevated blood pressures once they are
adequately beta-blocked (eg HR 60-70) are usually restarted on their
home anti-hypertensives. \emph{Home anti-hypertensives are not routinely
restarted postoperatively}

\textbf{Respiratory}

Chest X-ray on admission to ICU. Chest tubes generally consist of a 28Fr
Blake drain placed into the right chest. This is placed to water seal
when output is less than 200mL/day and there is no leak visible in the
Pleurevac container. Chest tubes are usually removed once output is less
than 150mL/day and drainage is clear without evidence of chyle (milky
appearance).

\textbf{Gastrointestinal} All patients receive pantoprazole 40mg IV
daily and metoclopramide 5-10mg IV q 6hrs

\textbf{Nutrition}

All esophagectomy patients receive a feeding jejunostomy at the time of
operation. In the immediate post-operative period, patients receive
Osmolite 1.5 starting the day of surgery once they are off pressors.
Tube feeds are started at 20mL/hour until flatus and then advanced at
10mL per hour every 8 hours, up to a goal of 60mL per hour (x24 hours).
Patients who are on tube feeds prior to surgery are generally restarted
on their home tube feed formula.

Patients who do not tolerate Osmolite are switched to Vital 1.5, which
is pre-hydrolyzed. In order to allow enough time for switching of tube
feedings, patients are generally switched from Vital to their home tube
feeding formula 3-4 days prior to discharge. This is typically done when
they are transferred out of the ICU.

Obese patients (BMI\textgreater30) are started on Promote at 20mL/hour
and increased to a goal of 60mL/hour. Promote contains a more protein
relative to carbohydrdates. An alternative is Vital High Protein, which
is similar to Promote but using hydrolyzed proteins.

Diarrhea in patients on tube feeds (especially nocturnal diarrhea) needs
to be addressed. Despite STICU guidelines for nutrition in trauma
patients, diarrhea (especially night-time diarrhea) is justification for
alteration in tube feeds. See Diarrhea and Jejunstomy Feeding

\emph{Diabetic Patients}

Patients who require EndoTool in the ICU will need an endocrinology
consult. See also \hyperref[jejunostomy_diabetes]{Jejunostomy Feedings
with Diabetes}

\emph{`Free Water'}

In addition to tube feeds, most patients will receive `free' water
flushes through the jejunostomy. This is typically done as 240mL four
times per day (for a total of 32oz)

\emph{Nasogastric Tubes}

A silicone nasogastric tube (Covidien Salem Sump) is placed in all
patients during surgery and the position confirmed by ultrasound
intraoperatively. Tubes are positioned so that all 4 dots are outside.
Gastric emptying is evaluated with upper GI prior to NG tube removal.
Once extubated, upper GI is typically performed on the 2nd through 4th
postoperative day. Conversely, patients who are intubated will keep
their nasogastric tube until extubated. Radiology ordered as ``upper GI
Series''. In the comments section please add ``IsoVue through NG tube.
Contact Dr Salo for study''. See Evaluation of GI Function

\emph{Drains}

\begin{itemize}
\tightlist
\item
  JP1: 19Fr Blake drain in left pleura. The exit site the most lateral
  drain and is secured with a blue suture
\item
  JP2: 19Fr Blake drain in right pleura. The exit site is medial to JP1
  and is secured with a black suture
\item
  JP3: 19 Fr Blake drain in abdomen. If used, the exit site is most
  medial and is secured with a blue suture
\item
  JP4: 15Fr Blake drain in neck (for cervial incision)
\item
  JP5: 15Fr Blake drain in subcutaneous tissue of incision
\end{itemize}

\emph{Evaluation for Anastomotic leak}

\hyperref[drain_amylase]{Drain amylase} is an inexpensive, specific, and
relatively sensitive test for anastomotic leak. Fluid from JP2 (right
chest) is sent for ``Body Fluid Amylase'' starting on postoperative day
\#4 and continued until postoperative day 9. Drain amylase over 400IU/ml
is considered positive and prompts a \hyperref[ct_esophagram]{CT
esophagram} for confirmation. See also
\hyperref[leak_evaluation]{Evaluation for Anastomotic Leak}

\textbf{Renal} Total fluids (IV + tube feeds) are generally run at
75mL/hour. Foley catheter is removed on the first or second day after
surgery. Some patients will need diuresis on the 3rd or 4th
postoperative day

\textbf{Heme} All patients require VTE prophylaxis with Lovenox or
heparin SQ.

Preoperative anti-platelet agents are started on the first (aspirin) or
second (Plavix) postoperative day if there is no excessive bleeding from
the chest tube or JP drains. Patients on preoperative anticoagulation
are transitioned to therapeutic Lovenox on the second postoperative day
if no signs of bleeding.

\textbf{ID} Prophylactic Cefazolin and Flagyl are administered for 24
hours and stopped

\chapter{Esophagectomy Ward}\label{esophagectomy-ward}

See
\href{https://gisurgonc.github.io/talks/eso_postop_care.htm}{Esophagectomy
Postop Care}

\textbf{Ward Care}

Once stable, patients are transferred to 6T. If a 6T bed is not
available, please notify Dr Salo. Historically, \textgreater98\% of
patients are transferred to 6T after leaving the ICU.

\section{Anti-Hypertensives}\label{anti-hypertensives}

Patients are transitioned to enteral metoprolol once they are
transferred to the ward. Patients on IV metoprolol at 2.5mg IV q6 hours
are started on 25mg enteral bid, while patients receiving 5mg IV q6
hours are started on 50mg bid. For patients who are not taking medicines
by mouth, liquid metoprolol can be ordered while an inpatient. (Liquid
metoprolol is \textbf{\emph{not}} available as a home medicine.)

Home anti-hypertensives are usually held in order to allow
beta-blockade. Patients who have elevated blood pressures once they are
adequately beta-blocked (eg HR 60-70) are usually restarted on their
home anti-hypertensives. \emph{Home anti-hypertensives are not routinely
restarted postoperatively}. Home antihypertensives are restarted on a
selective basis.

\section{Epidural Catheter}\label{epidural-catheter}

Epidural catheters are routinely placed. Foley catheters are usually
left in place until the epidural is removed. Preoperative heparin is
\textbf{\emph{not}} administered if an epidural is to be placed, but is
administered by anesthesia at least one hour after epidural placements.
Postoperatively, prophylactic Lovenox 40mg sq daily is used. Therapeutic
Lovenox and anti-platelet agents are not used until the epidural
catheter is removed.

\section{Chest tubes}\label{chest-tubes}

Chest tubes generally consist of a 28Fr Blake drain placed into the
right chest. This is placed to water seal when output is less than
200mL/day and there is no leak visible in the Pleurevac container. Chest
tubes are usually removed once output is less than 150mL/day and
drainage is clear without evidence of chyle (milky appearance).

\subsection{Chest tube removal}\label{chest-tube-removal}

\begin{itemize}
\tightlist
\item
  Supplies needed: Suture removal kit, 4x4 gauze, xeroform or Vaseline
  gauze, 2'' silk tape, red biohazard bag
\item
  Notify the bedside RN
\item
  Explain procedure to patient including the need for prolonged
  expiration during removal
\item
  Prepare a bandage with 4x4 gauze and Xeroform
\item
  Cut and remove suture
\item
  Cover the exit site with the Xeroform bandage
\item
  Rehearse process with patient: prolonged expiration
\item
  Pull tube and apply bandage
\item
  Place tube and Atrium/Pleurevac in red biohazard bag and dispose of in
  soiled utility room
\item
  Return patient room to prior condition (including lowering bed)
\end{itemize}

Chest x-ray 4 hours after removal to look for a pneumothorax.

\section{GI Medicines}\label{gi-medicines}

All patients receive pantoprazole 40mg IV daily and metoclopramide
5-10mg IV q 6hrs. This is later switched to enteric PPI and reglan. For
patients \textless age 75, remeron is added as 15mg enteral qhs.

\section{Evaluation for leak}\label{leak_evaluation}

Experience suggests that the median time to the diagnosis of leak is 7
day after surgery. Currently the risk of anastomotic leak at CMC after
transthoracic esophagectomy is 2\%. The most sensitive test for the
diagnosis of leak is CT esophagram (see below). Based upon institutional
experience, patients are divided into low risk of leak vs high risk for
leak based upon drain amylase level and white blood cell count. Patients
with a normal drain amylase (400IU/ml) between postoperative day 4 and
day 7 AND white blood cell count less than 12 are considered low risk.
Patients with either an elevated drain amylase or WBC greater than 12
are evaluated with CT esophagram. In a group of 100 patients, several
were found to have elevated drain amylase in the first four days (up to
2000Iu/ml) which subsequently declined, and no evidence of leak was
found

\subsection{Drain Amylase}\label{drain_amylase}

JP2 is placed into the right pleura, passes through the hiatus, and is
brought out through a trocar site in the medial left upper quadrant.
Drain amylase from JP2 is tested beginning on postoperative day \#4
until discharge or postoperative day \#9. JP2 is generally removed prior
to patient discharge. JP1 is placed in the left pleura and is generally
not tested for amylase.

\subsection{CT esophagram}\label{ct_esophagram}

This is the most sensitive study for the detection of anastomotic leak.
In order to obtain sufficient sensitivity, it requires a pre-contrast
scan and the administration of contrast into the esophagus. The need for
a pre-contrast scan means that the presence of remnant barium in the
esophagus (from a Modified Barium Swallow) makes it more difficult to
interpret the scan and should be avoided. Awake patients can drink the
contrast. Patients with an NG tube present at the time of the study
generally will have contrast administered through their NG tube. Almost
all esophagectomy patients will have a Covidien silicone Salem Sump 18Fr
tube in place. This has four marks on the tube at 45, 55, 65, and 75cm.
The tube is typically positioned with the 4th mark at the nares, which
means that the tip is 45cm from the nares and is usually within the
gastric conduit AND below the anastomosis. The NG tube also contains
side holes which extend 8.5cm above the tip of the tube. The optimal
study is done with the NG tube withdrawn so that there are two
side-holes above the anastomosis, which means that the tip is
approximately 5cm below the anastomosis. A scout CT will be performed,
and the radiologist will determine how far back the NG tube needs to be
withdrawn. This is communicated to the CT technician, who asks the nurse
(or physician) to withdraw the tube the calculated amount (making a note
of the starting position of the tube relative to the four marks). This
should keep the tip below the level of the anastomosis, so that after
the CT scan, the NG tube can be (blindly) advanced back to its original
position (at approximately 45cm from the nares).

Patients who clinically deteriorate prior to post-operative day 7 in
whom there is high suspicion for a leak (fevers, pleural effusion on
chest X-ray, elevated JP amylase, respiratory failure) will undergo a CT
esophagram earlier than post-operative day 7.

\section{Anastomotic Leak Treatment}\label{anastomotic-leak-treatment}

If a patient is demonstrated to have a leak, they are made NPO AND will
have any pleural effusion treated with a pigtail catheter. Conservative
management will generally be successful for most patients with leaks
provided 1) there is no evidence of conduit necrosis 2) nutrition is
optimized and 3) empyema is treated. Patients with leaks may even need a
decortication, which emphasizes the importance of CT scan in the sick
post-operative esophagectomy patient, so that an empyema can be
diagnosed and treated. Patients with leaks who show signs of systemic
illness may need to be considered for an intraluminal stent. Patients
who are profoundly ill need to be evaluated for gastric necrosis with
upper endoscopy.

\textbf{Nutrition}

All esophagectomy patients receive a feeding jejunostomy at the time of
operation. In the immediate post-operative period, patients receive
Osmolite 1.5 starting the day of surgery once they are off pressors.
Tube feeds are started at 20mL/hour until flatus and then advanced at
10mL per hour every 8 hours, up to a goal of 60mL per hour (x24 hours).
Patients who are on tube feeds prior to surgery are generally restarted
on their home tube feed formula.

Patients who do not tolerate Osmolite are switched to Vital 1.5, which
is pre-hydrolyzed. In order to allow enough time for switching of tube
feedings, patients are generally switched from Vital to their home tube
feeding formula 3-4 days prior to discharge. This is typically done when
they are transferred out of the ICU.

Obese patients (BMI\textgreater30) are started on Vital High Protein at
20mL/hour and increased to a goal of 60mL/hour. Promote contains a more
protein relative to carbohydrdates.

Diarrhea in patients on tube feeds (especially nocturnal diarrhea) needs
to be addressed. Despite STICU guidelines for nutrition in trauma
patients, diarrhea (especially night-time diarrhea) is justification for
alteration in tube feeds. See Diarrhea and Jejunostomy Feeding

\emph{Diabetic Patients}

Patients on tube feeds are typically started on continuous
(around-the-clock) tube feedings, and subsequently changed to a
nocturnal regimen (typically 6pm to 10am). See
\hyperref[jejunostomy_diabetes]{Jejunostomy Feeds in Diabetic Patients}

\emph{`Free Water'}

In addition to tube feeds, most patients will receive `free' water
flushes through the jejunostomy. This is typically done as 240mL four
times per day (for a total of 32oz)

\emph{Nasogastric Tubes}

A silicone nasogastric tube (Covidien Salem Sump) is placed in all
patients during surgery and the position confirmed by ultrasound
intraoperatively. Tubes are positioned so that all 4 dots are outside.
Gastric emptying is evaluated with upper GI prior to NG tube removal.
Once extubated, upper GI is typically performed on the 2nd through 4th
postoperative day. Conversely, patients who are intubated will keep
their nasogastric tube until extubated.

Upper GI is ordered as \emph{FL Upper GI Track Single Contrast}

\emph{Drains}

\begin{itemize}
\tightlist
\item
  JP1: 19Fr Blake drain in left pleura. The end of the tube is cut at an
  angle. The exit site the \emph{usually} the most lateral/inferior
  drain.
\item
  JP2: 19Fr Blake drain in right pleura. The exit site is \emph{usually}
  medial/superior to JP1
\item
  JP3: 19 Fr Blake drain in abdomen. If used, the exit site is
  \emph{usually} most medial
\item
  JP4: 15Fr Blake drain in neck (for cervical incision)
\item
  JP5: 15Fr Blake drain in subcutaneous tissue of incision
\end{itemize}

\emph{Evaluation for Anastomotic leak}

Drain amylase is an inexpensive, specific, and relatively sensitive test
for anastomotic leak. Fluid from JP2 (right chest) is sent for ``Body
Fluid Amylase'' starting on postoperative day \#4 and continued until
postoperative day 9. Drain amylase over 400IU/ml is considered positive
and prompts a CT esophagram for confirmation. See also Evaluation for
Anastomotic Leak

\textbf{Renal} Total fluids (IV + tube feeds) are generally run at
75mL/hour. Foley catheter is removed on the first or second day after
surgery. Some patients will need diuresis on the 3rd or 4th
postoperative day

\textbf{Heme} All patients require VTE prophylaxis with Lovenox.

Preoperative anti-platelet agents are started on the first (aspirin) or
second (Plavix) postoperative day if there is no excessive bleeding from
the chest tube or JP drains. Patients on preoperative anticoagulation
are transitioned to therapeutic Lovenox on the second postoperative day
if no signs of bleeding.

\textbf{ID} Prophylactic Cefazolin and Flagyl are administered for 24
hours and stopped

\section{Labs}\label{labs}

Once on the ward, BMP and CBC are checked every other day. Patients with
leukocytosis (\textgreater12,000) are monitored with daily CBC until
resolved.

\section{Discharge Medicines}\label{discharge-medicines}

Patients after esophagectomy typically go home with the following
medicines:

Proton pump inhibitors (will continue for 2 years) Oxycodone elixir via
feeding tube. Current STOP guidelines dictate that patients receive no
more than a 7 day supply of opiods at discharge Reglan 10mg po qid (will
stop at 6 weeks post-op) Remeron 15mg qhs (will continue for 3 months
post-op) Tylenol 1000mg q6 hours as elixir (pediatric form) Gabapentin
300mg as liquid tid x 14 days Metoprolol if started postoperatively
(most patients). If not on beta blockers preoperatively, this will be
cut in half at the first visit and stopped at the second postoperative
visit. It is important to limit the use of medicines via jejunostomy in
order to lower the risk of clogging the feeding tube. While liquid
metoprolol is available for inpatients, it is not available from
outpatient pharmacies.

\section{Diet at discharge}\label{diet-at-discharge}

\begin{itemize}
\tightlist
\item
  70\% will pass their MBS for thin liquids are and discharged on
  protein shakes
\item
  15\% will pass for nectar thick but fail for this liquids. They are
  discharged taking medicines with a sip of thickened water but
  otherwise NPO
\item
  15\% will fail for nectar thick and are discharged NPO with medicines
  via jejunostomy
\end{itemize}

\part{Outpatient}

\chapter{Outpatient Jejunostomy Care}\label{outpt_jejunostomy}

\section{Jejunostomy Tubes}\label{jejunostomy-tubes}

Two types of tubes are used

\textbf{Dacron cuffed MIC tube}
\href{https://avanos.com/product-catalog/product/mic-jejunostomy-tubes/0301-14__mic-jejunostomy-feeding-tube-14-fr/}{0301-14}.
These tubes are usually only placed after esophagectomy. The dacron cuff
is designed to sit in the subcutaneous tissue. If the cuff becomes
incorporated, the tube can stay in placed for extended periods of time.
Unfortunately, the cuff is prone to infection. Patients with infection
of a cuffed jejunostomy will need antibiotics (eg Augmentin) and orders
for stat \emph{IR J tube Manip Exchange}

Removal of cuffed tubes requires a small surgical procedure. Supplies
required:

\begin{itemize}
\tightlist
\item
  ER Laceration tray
\item
  30mL 1\% xylocaine (prefer with epinephrine)
\item
  15 blade scalpel
\item
  Non-sterile gloves
\end{itemize}

Patient will need a surgical consent and a procedure note.

\textbf{Balloon Tube MIC}
\href{https://avanos.com/product-catalog/product/mic-jejunal-j-feeding-tubes/8200-14__mic-jejunal-feeding-tube-with-enfit-connector-14-fr/}{8200-14}

These tubes are held in place by a balloon simlar to a foley catheter.
The balloon is typically inflated with 5mL of water. Saline is not used
in the balloon as the salt may crystallize in the balloon valve. Balloon
tubes can be removed by deflating the balloon and traction. A gauze is
placed over the site. Patients are instructed that the site can drain
for a few days. A jejunal-cutaneous fistula develops in approximately
1\% of patients, who will will placement of an ostomy bag over the site
and total parenteral nutrition (TPN) in order for the fistulous tract to
heal.

\chapter{Outpatient Esophagectomy}\label{outpt_mie}

Patient are usually seen 7-10 days after hospital discharge.

\section{Postop Evaluation}\label{postop-evaluation}

By system:

\textbf{Neurologic - Pain control}

Patient are typically discharged on:

\begin{itemize}
\tightlist
\item
  acetaminophen 1000mg qid
\item
  gabapentin 300mg tid (liquid)
\item
  oxycodone elixir 5mg q6 hrs prn
\end{itemize}

If patients need additional oxycodone, this is renewed 2 weeks at a
time, typically to correspond to a clinic day. Their current daily usage
at the time of their first postop visit is estimated, and the first
prescription plans for the same daily amount for a two week period. If
they are taking 5mg of oxycodone 5 times per day, a prescription is
prepared for 5 x 14 = 70 doses. The prescription is annotated ``Two week
supply at 5 doses per day''. Patients are told the prescription will be
renewed 2 weeks later at a reduced dose (typically one fewer dose per
day). If the oxycodone bottle is lost, patients are prescribed tramadol
(50mg q6 hours). Patient are told that if they need long-term pain
management, they will need the assistance of their primary care
provider.

Patients are instructed that they are not ``cleared'' for driving as
long as they are taking oxycodone

\textbf{Cardiovascular}

Beta blockers are used around the time of surgery to prevent atrial
fibrillation. Management depends upon patients's preoperative
medications:

\ul{Patients not on beta-blockers or antihypertensives prior to surgery}
are started on metoprolol 25 to 50mg enteral bid in the hospital. At
their first postoperative visit, metoprolol is cut to once per day for
one week, then stopped.

\ul{Patients taking antihypertensives but not on beta-blockers prior to
surgery} are started on metoprolol 25 to 50mg enteral bid in the
hospital. In some cases, their anti-hypertensives are held to prevent
hypotension. At their first postoperative visit, metoprolol is
discontinued and their home antihypertensives resumed.

\ul{Patients taking beta-blockers} are kept on their beta blockers in
the hospital and after surgery

If patients have been experienced atrial fibrillation in the hospital or
have new cardiac medicines, they will need cardiology follow up for
management. Typically this is amiodarone.

\textbf{Respiratory}

Pleural effusion (typically on the right) is a risk after esophagectomy.
If the lung exam suggests a pleural effusion, a chest x-ray PA/Lateral
can be obtained same-day in LCI I.

Chest incisions should be inspected for signs of wound infection and
retained silk (black) suture at the chest tube site

Patients who are hypoxemic should be further investigated with chest
x-ray. Delayed pulmonary embolism would be unusual, but would need to be
considered and evaluated with a CT angiogram.

\textbf{GI}

Patients who pass their modified barium swallow are discharged on
protein shakes and water by mouth. They are advanced to a Phase I diet
at their first post-operative visit.

Patients who do not pass their modified barium swallow will need an
outpatient MBS ordered at the first postop visit. They will continue
with either 1) meds via J tube or 2) meds with thickened water,
depending upon their MBS results in the hospital.

Skin staples are removed, and the jejunostomy site inspected for
infection. Infections are treated with Augmentin for 7-10 days. In some
cases, infected cuffed jejunostomy tubes will need to be replaced in
interventional radiology. See \href{31-outpt_jejunostomy.htm}{J tubes}

Tube feeds are typically run for 16 hours nocturnal - 5 cartons for men
(75mL/hour) and 4 cartons for women (60mL/hour). Free water is usually
240mL 4x/day.

\hyperref[outpt_tubefeeds_diabetes]{Diabetics on Tube Feeds receiving
insulin} require special management.

70\% of patient are taking oral liquids when discharged, including
protein shakes. For patients that can get one Premier protein, this
represents 30gm of protein, which is the equivalent of 2 cartons of tube
feeds. These patients can often have their tube feeds decreased to 3
cartons/night, which allows them to run for only 12 hours.

\textbf{Diarrhea} Loose stools during the day are to be expected, but
diarrhea at night is problematic both for sleep and quality of life.
This may require a change in formula from intact protein (Osmolite 1.5)
to a peptide formula (Vital 1.5). If patients are taking enough protein
orally through protein shakes, it may be possible to decrease the rate
of tube feeds.

\textbf{Constipation} can be easily treated by increasing free water to
240mL 6x/day (48oz) or even 240mL 8x/day (64oz). Miralax is generally
not required with a jejunostomy.

Reglan is used for 6 weeks after surgery, then stopped. If patients have
reflux symptoms, this is renewed for another 6 weeks.

Remeron (if used) is given for 12 weeks after surgery, then stopped.

Proton pump inhibitors are used for 12 months after surgery, to prevent
anastomotic stricture.

\textbf{Renal}

Patient with renal dysfunciton in the hospital should have BMP checked
at their first postoperative visit.

\textbf{Hematologic}

Patients on Plavix or anticoagulation should have had theses medicines
restarted before discharge.

\textbf{ID}

New fevers at home require investigation with CXR and CBC or CT
chest/abdomen/pelvis as indicated

\section{Medication Management}\label{medication-management}

Beta blockade and anti-hypertensives

Atrial fibrillation occurs in the postoperative period in 20\% of
patients. For this reason, beta blockers are used to atrial fibrillation
prevention. In the ICU, metoprolol is administered intravenously and is
subsequently changed to enteral administration. Because metoprolol also
lowers blood pressure, patients on home anti-hypertensives frequently
have them held during their hospital stay (and after discharge). At the
first postoperative visit, patients who were not treated with beta
blockade preoperatively are generally switched from metoprolol back to
their home anti-hypertensives.

\section{Special Cases}\label{special-cases}

\subsection{Weaning Tube Feeds in
Diabetics}\label{outpt_tubefeeds_diabetes}

As outpatients begin eating more orally, their tube feeds are reduced.

Weaning from 5 cans to 4 cans: Easiest method is to maintain the same
schedule (16 hours) and reduce insulin dosage by 20\%. For instance, the
above patient who is on 5 cans at 75mL/hour x 16 hours is receiving 16N
+ 8R at start of tube feeds and 6 hours later. This patient could be
weaned by reducing rate form 75mL/hour x 16 hours to 60mL/hour x 16
hours and reducing insulin to 12N + 6R at the start of tube feeds AND
another dose of 12N + 6R after 6 hours.

Weaning from 4 cans to 3 cans: One option is to decrease the duration of
tube feeds from 16 hours to 12 hours, while maintaining rate of
60mL/hour. In this case, the insulin dosage could be kept the same at
the start of tube feeds, BUT the dose 6 hours after the start of tube
feeds could be omitted.

Once patients are on 3 cans per night, further weaning can be
accomplished by skipping tube feeds (and insulin) every other night in a
``tube feed holiday''. This allows an evening of interrupted sleep and
can tend to increase the appetite the morning after tube feeds are held.

\part{Esophageal Cancer}

\chapter{Esophageal Overview}\label{EsoIntro}

Esophageal cancers can be grouped into 4 treatment categories:

\begin{itemize}
\tightlist
\item
  \hyperref[superficial]{Superficial} \(\rightarrow\) Endoscopic therapy
\item
  \hyperref[localized]{Localized} \(\rightarrow\) Primary surgery
\item
  \hyperref[locally_advanced]{Locally Advanced} \(\rightarrow\)
  Trimodality therapy
\item
  \hyperref[metastatic]{Metastatic} \(\rightarrow\) Systemic therapy
\end{itemize}

Patients with minimal dysphagia, no weight loss, and small (\textless3cm
length) tumors are evaluated with endoscopic ultrasound:

\begin{itemize}
\tightlist
\item
  If uT1 on EUS and \textless2cm in size, \hyperref[emr]{endoscopic
  mucosal resection} yields more information and may be therapeutic for
  tumors with negative margins and without high-risk features.
\item
  If uT2N0 on EUS, and PET scan shows a small tumor (MTV
  \textless10cm\textsuperscript{3}), \hyperref[primary_surgery]{primary
  surgery} is preferred in patients who are good surgical risks
\item
  If T3 or N+ on EUS, if PET shows no metastatic disease,
  \hyperref[trimodality]{neoadjuvant therapy} is optimal)
\end{itemize}

Patients with dysphagia to solids or weight loss or tumor length
\textgreater3cm are unlikely to have T1-2 tumors and can be evaluated
with \hyperref[pet]{PET scan}.

\begin{itemize}
\tightlist
\item
  If PET shows disease confined to the esophagus and regional nodes,
  \hyperref[trimodality]{trimodality therapy} (chemoradiation followed
  by surgery) is optimal.
\item
  If PET shows metastatic disease, patients are eligible for palliative
  chemotherapy with radiation for treatment of symptoms of dysphagia.
\item
  If PET shows extra-regional lymph node disease, patient is at high
  risk for distant disease and can be treated with induction
  chemotherapy followed by chemoradiation and surgical evaluation.
\end{itemize}

\chapter{EsoCa SCORE - JR}\label{EsoObjJB}

\href{https://www.surgicalcore.org/modulecontent.aspx?id=1000534}{Junior
Resident SCORE - Esophageal Neoplasms}

\href{https://www.surgicalcore.org/modulecontent.aspx?id=164703}{Junior
Resident SCORE - Esophagectomy}

\textbf{Anatomy}

\textbf{Epidemiology and Prevention}

\begin{itemize}
\tightlist
\item
  Worldwide distribution
\item
  Risk factors
\end{itemize}

\textbf{Presentation}

\begin{itemize}
\tightlist
\item
  Symptoms
\item
  Nutritional consequences
\end{itemize}

\textbf{Diagnosis and Staging}

\begin{itemize}
\tightlist
\item
  Physical findings
\item
  Role of PET scan
\item
  Role of EUS
\item
  TNM staging system
\item
  Treatment Categories

  \begin{itemize}
  \tightlist
  \item
    Superficial
  \item
    Localized
  \item
    Locally-advanced
  \item
    Metastatic
  \end{itemize}
\end{itemize}

\textbf{Operative Treatment}

\begin{itemize}
\tightlist
\item
  Indications for surgery
\item
  Operative anatomy
\item
  Gastric tube construction
\item
  Ivor Lewis
\item
  Transhiatal
\item
  McKeown
\end{itemize}

\textbf{Complications}

\begin{itemize}
\tightlist
\item
  Postoperative hypovolemia
\item
  Chylothorax
\item
  Anastomotic leak - chest
\item
  Anastomotic leak - neck
\item
  Atrial fibrillation
\item
  Aspiration
\item
  Vocal cord paralysis
\item
  Stomach ischemia
\end{itemize}

\textbf{Intraop Decision-making}

\begin{itemize}
\tightlist
\item
  Liver metastasis
\item
  Peritoneal lesion
\end{itemize}

\textbf{Nonoperative management}

\begin{itemize}
\tightlist
\item
  Palliative Radiation
\item
  Esophageal stents
\end{itemize}

Resources

\href{https://gisurgonc.github.io/talks/eso_ca_rx_talk.htm}{Esophageal
Cancer Talk}

\href{https://gisurgonc.github.io/talks/eso_ca_cases.htm}{Esophageal
Cancer Cases}

\chapter{EsoCa Objectives - Chief}\label{EsoObjCR}

\href{https://www.surgicalcore.org/modulecontent.aspx?id=1000690}{Chief
Resident SCORE}

\textbf{Epidemiology and Prevention}

\begin{itemize}
\tightlist
\item
  Worldwide distribution
\item
  Risk factors for squamous cell carcinoma
\item
  Risk factors for adenocarcinoma
\end{itemize}

\textbf{Presentation}

\begin{itemize}
\tightlist
\item
  Symptoms
\item
  Nutritional consequences
\end{itemize}

\textbf{Diagnosis and Staging}

\begin{itemize}
\tightlist
\item
  Physical findings
\item
  Role of PET scan
\item
  Role of EUS
\item
  TNM staging system
\item
  Treatment Categories

  \begin{itemize}
  \tightlist
  \item
    Superficial
  \item
    Localized
  \item
    Locally-advanced
  \item
    Metastatic
  \end{itemize}
\end{itemize}

\textbf{Multidisciplinary Management}

\begin{itemize}
\tightlist
\item
  High-grade dysplasia
\item
  Endoscopic therapy
\item
  Primary surgical therapy
\item
  Trimodality therapy
\end{itemize}

\textbf{Operative Mangement}

\begin{itemize}
\tightlist
\item
  Ivor Lewis
\item
  Transhiatal
\item
  Left thoracoabdominal
\item
  McKeown
\item
  Alternative conduits
\item
  Complications and their management
\end{itemize}

\textbf{References}

\begin{itemize}
\tightlist
\item
  CROSS Trial \href{https://www.youtube.com/watch?v=4XWNr4rJXQ}{Behind
  the Knife Podcast}
\item
  MAGIC Trial
\item
  FREGAT
\item
  Dutch TIME trial
\end{itemize}

\chapter{Staging}\label{staging}

The staging workup begins once a diagnosis is made on endoscopy.

The first step is to make a preliminary determination whether the tumor
is early stage (and can be treated with endoscopy or primary surgery) or
later stage (and treated with chemoradiation followed by surgery or
with)

Patients with minimal dysphagia, no weight loss, and tumors with less
than 3cm cranio-caudal extent have a reasonable change of being T1 or T2
tumors. Tumors \textless3cm in length are much more likely to represent
T1-2 lesions than those \(\geq\) 3cm(Hollis et al. 2017). For these
patients, determining the precise T stage is important in their workup,
so \textbf{endoscopic ultrasound} is the most frequent staging study
after diagnosis.

Patients who present with dysphagia are likely to have T3 or T4 disease,
which is generally treated with neoadjuvant chemoradiation followed by
surgery. Data from Memorial Sloan Kettering (Ripley et al. 2016) among
61 patients with esophageal cancer who presented with dysphagia, 54
(89\%) were found on EUS to have uT3-4 tumors. On the other hand, among
53 patients without dysphagia, 25 (47\%) were uT1-2, and were
potentially candidates for primary surgery. Their conclusion was that
EUS could be omitted from the workup of patients with dysphagia, but is
useful in patients without dysphagia.

PET can be helpful in evaluating patients who may have T1-2 disease, and
might be candidates for primary surgical therapy. A comparison of PET
and EUS (Malik et al. 2017) showed that uT1-2 tumors had median
metabolic tumor volume (MTV) of 6.7cm\textsuperscript{3}, compared with
uT3-4 tumors, with a median MTV of 35.7cm\textsuperscript{3}.

\chapter{Barrett's Esophagus}\label{barretts-esophagus}

AGA recommendations for endoscopic ablation of Barretts:

Strong recommendation for ablation of high grade dysplasia. Recommends
surveillance at 3,6,12 months, then annually.

Shaheen randomized trial of radiofrequency ablation (RFA) for dysplasia
(Shaheen et al. 2009). Mean age 66 years. Mean length of Barrett's
esophagus 5.3cm. RFA resulted in 77\% rate of eradication of Barrett's
at one year.

Meta-abalysis suggests ablation can reduce rate of development of cancer
from 6.6 per 100 person years to 1.9 (Rastogi et al. 2008)

RFA for Barrett's national registry (Ganz et al. 2008)

RFA for low-grade dysplasia (Phoa et al. 2014) resulted in 25\% risk
reduction in progression to HGD.

\section{Anti-reflux surgery for Barrett's
esophagus}\label{anti-reflux-surgery-for-barretts-esophagus}

\chapter{Superficial EsoCa (T1)}\label{superficial}

Superficial esophageal cancer is usually asymptomatic, which means that
the diagnosis is generally made in the context of surveillance for
Barrett's esophagus.

Nodular Barrett's esophagus can be best evaluated with endoscopic
mucosal resection, which can provide further staging information if an
adenocarcinoma is found, such as depth of invasion, differentiation, and
lymphovascular invasion.

Larger lesions arising within Barrett's esophagus should first be
evaluated with endoscopic ultrasound (EUS)

EUS is less sensitive for T1 lesions (Bergeron et al. 2014)
-\textgreater{} use EMR for diagnosis (Maish and DeMeester 2004)

T1a tumors have a low risk of nodal metastasis (Dunbar and Spechler
2012)

\section{Endscopic Mucosal Resection (EMR)}\label{emr}

For patients with nodular Barrett's esophagus or small tumors judged to
be T1 by endoscopic ultrasound, endoscopic mucosal resection (EMR) can
be diagnostic and potentially curative.

EMR is likely sufficient for small tumors with favorable pathologic
factors(Nurkin et al. 2014):

\begin{itemize}
\tightlist
\item
  Size less than 2cm
\item
  Lateral and deep margins clear
\item
  Absence of lymphovascular invasion
\item
  Well- or moderately- differentiated
\end{itemize}

Need for RFA of Barrett's after EMR: (Haidry et al. 2013)

Combination therapy with EMR and RFA results in lower rate of recurrence
than EMR alone.(O. Pech et al. 2008) 1000 patients with T1a
adenocarcinoma werer treated with EMR followed by RFA with only 2
(0.2\%) of patient dying of esophageal cancer with median followup of
almost 5 years.(Oliver Pech et al. 2014)

Review of EMR: (Soetikno et al. 2005)

\section{Endoscopic submucosal dissection (ESD).}\label{esd}

For T1a tumors, EMR and ESD appear to have similar outcomes (Terheggen
et al. 2017) Terheggen, G. et al.~A randomised trial of endoscopic
submucosal dissection versus endoscopic mucosal resection for early
Barrett's neoplasia. Gut 66, 783--793 (2017).

Endoscopic submucosal dissection is a technique for deeper endoscopic
removal of esophageal lesions using endoscopic cautery, which dissects
through the submucosa. ESD has a higher rate of curative resection (Cao
et al. 2009) albeit at the cost of prolonged operative times and
increased risk of complications such a bleeding. (Repici et al. 2010)

ESD takes more time and has higher R0 resection rate but similar
recurrence rate at 2 years (Terheggen et al. 2017)

\chapter{Localized EsoCa}\label{localized}

\section{T1b Tumors}\label{t1b-tumors}

\section{T2N0 Tumors}\label{t2n0-tumors}

Multiple studies have failed to show the additional benefit of
chemotherapy or chemoradiation for pT2N0M0 esophageal cancer patients
treated with radiation.

Neoadjuvant chemo not likey to be helpful for early stage disease - FFCD
9901 (C. Mariette et al. 2014) enrolled patients with T1-2 or T3N0
tumors to chemoradiation followed by surgery versus surgery alone. The
majority of the tumors (72\%) were squamous cell carcinoma.Postoperative
mortality was significaly increased in the chemoradiation arm (11.1\% vs
3.4\%).

Meta-analysis of 5265 patients in 10 studies showed that while
neoadjuvant therapy was associated with a reduction in positive margin
rate, there was no difference in terms of recurrence or
survival.(\textbf{mota176?})

French trial FREGAT(Markar et al. 2016)

Retrospective review of the National Cancer DataBase failed to
demonstrate a difference in survival of cT2N0M0 esophageal cancer with
or without preoperative chemoradiation.(Speicher et al. 2014)

A retrospective report from Johns Hopkins examined outcomes of T2N0
squamous cell carcinoma patients and showed equivalent outcomes for
primary surgery vs neoadjuvant chemoradiation followed by surgery (Zhang
et al. 2012)

Retrospective reveiw from Utrecht found longer survival for cT2N0
patients treated with trimodality therapy compared with primary surgery
(\textbf{wirsik799?})

Indiana University: 731 patient undergoing esophagectomy, of whom 68
were cT2N0. Among 35 patients who had surgery alone, 17 (49\%) were
clinically understaged. Understaging was more common among
poorly-differentiated tumors (Hardacker et al. 2014)

\section{Staging of T2N0 Tumors}\label{staging-of-t2n0-tumors}

The challenge for treatment decision-making is the limited sensitivity
of endoscopic ultrasound in ruling out pT3 or pN+ disease. In other
words, if a patient who is thought to have cT2N0 disease undergoes
resection, and is found on pathology to have pT3 or
pN\textsuperscript{+} disease, this would dictate the need for
postoperative chemoradiation. In general, chemoradiation after
esophagectomy is difficult for patients to tolerate, with a 40 \% chance
of failure to complete therapy.

Data from the Cleveland Clinic looked at 53 patients judged to be T2N0
by endoscopic ultrasound (uT2N0) were treated with primary surgery.
Pathologic examination showed that 17 (37\%) were understaged by
endoscopic ultrasound, and were pathologic (pT3) in 4 or node positive
(pN\textsuperscript{+}) in 13 cases. These patients were treated with
postoperative adjuvant chemoradiation.(Rice et al. 2007)

It is critical, therefore, in patients for whom primary surgery is
contemplated, to attempt to identify those with occult T3 or N+ disease.

\subsection{PET for staging of early-stage esophageal
cancer}\label{pet-for-staging-of-early-stage-esophageal-cancer}

86 patients with esophageal adenocarcinoma and squamous cell carcinoma
(Mantziari et al. 2020) evaluated with PET: clinical T3/T4 was
significantly predicted by SUVmax 8.25, Total Lesion Glycolysis
\textgreater42 and Metabolic Tumor Volume \textgreater{}
10cm\textsuperscript{3}

Cohort of patients treated with primary surgery: risk factors for
clinical understaging are dsyphagia, poor differentiation, and tumor
length \textgreater3cm (Worrell et al. 2018)

\section{Primary Surgery}\label{primary_surgery}

NCCN recommends PET scan

Most common sites of metastasis are liver, lung, bones, adrenal.

PET detects occult metastasis in 10-20\% of cases Kim et al. (2009).
Among 129 patients with esophageal cancer, PET detected additional sites
of disease in 41\% and changed management in 38\% (Chatterton et al.
2009)

PET for restaging detects interval development of metastatic disease in
8-17\% of cases (Vliet et al. 2008)

\chapter{Locally Advanced EsoCa}\label{locally_advanced}

Tumors that are T2N\textsuperscript{+}M0 or T3NxM0 are considered
locally-advanced. The high rate of failure with surgery alone has led to
development of adjunctive therapies.

\section{Trimodality Therapy}\label{trimodality}

Trimodality therapy consists of chemoradiation (chemoRT) followed by
surgery. Between 2010 and early 2024, this approach was the standard of
care for locally-advanced esophageal cancer. In 2024, the ESOPEC trial
showed superior survival for \emph{adenocarcinoma} with perioperative
``sandwich'' FLOT chemotherapy compared to chemoRT. For squamous cell
carcinoma, trimodality therapy continues to be the standard of care.

CROSS trial randomized 364 patients with resectable esophageal and
gastroesophageal junction tumors (75\% adenocarcinoma) to neoadjuvant
chemoradiation consisting of 4,140 cGy of radiation with concurrent
carboplatin and paclitaxel or surgery alone.(Hagen et al. 2012) Clinical
node-positive disease was present in 16\%. Pathologic complete response
was seen in 23\% of adenocarcinoma and 49\% of squamous cell carcinomas.
Median overall survival was 49 months after trimodality vs 24 months
after surgery alone (p=0.003). Survival at 5 years was 47\% with
trimodality vs 34\% with surgery alone. Squamous cell carcinomas
appeared to have particular benefit, with a hazard ratio of 0.42 for
squamous cell vs 0.74 for adenocarcinoma. Median survival was improved
for adenocarcinoma from 27.1 months to 43.2 months, but the median
survival for squamous cell increased from 27.1months to 81.6 months for
squamous cell. Rate of R0 resection was higher with chemoradiation (92\%
vs 69\% p\textless0.001) and local recurrence rates lower (14\% vs 34\%
P\textless0.001), and peritoneal recurrence lower (4\% vs 14\%
P\textless0.001). Despite the relatively low dose of radiation, in-field
recurrences were less than 5\%. The primary cause of failure was distant
disease (31\%) and local/regional failure (14\%). Therapy was
well-tolerated with 17\% grade 3 toxicity.(Oppedijk et al. 2014)

Ten-year followup of the CROSS trial (Eyck et al. 2021) showed that the
primary benefit of CROSS regimen was in reducing local and loco-regional
recurrences. There was no difference in isolated distant recurrence
betwee the two arms.

NeoResII trial i(Nilsson et al. 2023) investigated interval between end
of CROSS neoadjuvant therapy and surgery. Original hypothesis was that
prolonged interval between chemoRT and surgery might increase pCR rate.
Patients randomized to short interval (4-6 weeks) vs long interval
(10-12 weeks). No difference in pCR rate or overall survival. However,
bottom quartile for survival was shorter in the prolonged interval
group. There does not appear to be an advantage in prolonged interval
between chemoradiation and for some patients (perhaps those with more
aggressive tumors), prolonged interval may be detrimental.

Alternative to carboTaxol for chemoradiation is FOLFOX (SOG trial
(Leichman et al. 2011))

Ongoing PROTECT trial compares FOLFOX to paclitaxel and carboplatin
(Messager et al. 2016)

See also \hyperref[eso_dcrt]{Definitive ChemoRT}

\section{Neoadjuvant Chemo
(adenocarcinoma)}\label{neoadjuvant-chemo-adenocarcinoma}

ESOPEC Trial (\textbf{heoppner323?}) compared FLOT chemotherapy with
CROSS chemoradiation for \emph{adenocarcinoma} of the esophagus and
found better survival with FLOT. FLOT resulted in superior 3-year
survival (57\% vs 51\%). Pathologic CR rate 16.7\% vs 10\%. Grade 3
toxicity 58\% vs 50\%. 90-day mortality 3.1\% with FLOT vs 5.6\% CROSS.
pCR rate was much lower than original CROSS trial (23\%). Notably,
patients with an incomplete pathologic response after CROSS (90\%) did
not receive adjvant nivolumab, which has since become the standard of
care due to the Checkmate 577 trial(Kelly et al. 2021).

POET Trial (Pre-Operative therapy in Esophageal adenocarcinoma Trial)
treated 119 patients with adenocarcinoma of the gastroesophageal
junction with either neoadjuvant chemotherapy (5-FU, leucovorin,
cisplatin) followed by surgery or induction chemotherapy with the same
agents, followed by chemoradiation (4000cGy with concurrent cisplatin
and etoposide). The study suffered from slow accrual, but there was a
suggestion of improved 3-year survival with preoperative chemoradiation
(47.4\% vs 27.7\% \emph{p}=0.07) as well as better local control (76.5\%
vs 59\%). In addition, chemoradiation was associated with a higher
pathologic complete response rate (15.6\% vs 2\%)(Stahl et al. 2009).

Neo-AEGIS Trial (Reynolds et al. 2017) Randomized patients with GE
junction adenocarcinomas to CROSS vs modified MAGIC (prior to 2018) or
FLOT (after 2018). Trial included Siewert I, II, and III tumors. Trial
was underpowered and did not show a difference between the regimens.

OEO2 clinical trial (Allum et al. 2009) Randomized trial of preoperative
chemotherapy with ECF for esophageal and GE junction cancers. Improved
survival with preoperative chemotherapy.

\subsection{\texorpdfstring{Neoadjuvant Chemotherapy \(\rightarrow\)
ChemoRT}{Neoadjuvant Chemotherapy \textbackslash rightarrow ChemoRT}}\label{neoadjuvant-chemotherapy-rightarrow-chemort}

TOPGEAR study (Leong et al. 2017) randomized patients with gastric or GE
junction adenocarcinoma to neoadjuvant chemo (ECF or FLOT) vs
neoadjuvant chemo followed by chemoRT (4500cGy). Higher pCR rate with
chemo \(\rightarrow\) chemoRT (17\% vs 8\%) but no difference in
survival at 67mo followup. Higher R0 resection rate with chemo
\(\rightarrow\) chemoRT (92\% vs 88\%) ECF used for 2/3 of patients and
FLOT for 1/3.

German trial (Stahl et al. 2009) randomized patients to preoperative
chemotherapy (A) vs preoperative chemotherapy followed by preoperative
chemoradiation (B). Higher pcR rate in arm B (15.6\% vs 2\%) and ypN0
resection (64.4\% vs 37.7\%).

\section{Postoperative
chemoradiation}\label{postoperative-chemoradiation}

Intergroup-0116 (Macdonald et al. 2001) (Smalley et al. 2012) treated
556 patients with adenocarcinoma of the stomach or GE junction with
surgery along vs surgery followed by postoperative chemoradiation. After
a median followup of over 5 years, median overall survival iin the
surgery alone group was 27 months vs 36 months in the postoperative
chemoradiation group (p=0.005) Decrease in local failure as the first
site of failure in the chemoradiation group (19\% versus 29\%).

Chemoradiation after resection of GE junction tumors (Kofoed et al.
2012) among a group of 211 patients with GE junction adenocarcinoma with
positive lymph nodes with improved 3-year disease-free survival (37\% vs
24\%).

\section{Restaging}\label{restaging-1}

PET is routinely done for restaging esophageal cancer after
chemoradiation. Interval metastasis is found in 9\% of cases (Noordman
et al. 2018). Similar findings in a European multicenteral study (Zijden
et al. 2024)

\section{Active Surveillance}\label{active-surveillance}

Patients treated with neoadjuvant therapy who have a robust clinical
response (as evidenced by PET and EGD) can be treated with intensive
surveillance.

EGD is poor predictor of pCR due to low sensitivity(Sarkaria et al.
2009)

\section{GE Junction}\label{ge-junction}

(Siewert, Stein, and Feith 2006)

\chapter{Metastatic EsoCa}\label{eso_metastatic}

\section{Palliative radiation}\label{palliative-radiation}

Palliative radiation vs chemoradiation (Penniment et al. 2018)

Radiation along favored over chemoradiation in the palliaitve setting
(Penniment et al. 2018)

\section{Chemoradiation vs chemotherapy in Stage
IV}\label{chemoradiation-vs-chemotherapy-in-stage-iv}

(Guttmann et al. 2017)

\section{Stents for malignant
disease}\label{stents-for-malignant-disease}

(Vakil et al. 2001)

Review of guidelines 2010 Am Society GI (Sharma, Kozarek, and Practice
Parameters Committee of American College of Gastroenterology 2010)

\section{Perc Esophagostomy}\label{perc-esophagostomy}

Percutaneous transesophageal gastrostomy (PTEG) can be used as an
alternative to nasogastric decompression, particularly in malignant
obstruction.

Systematic review of 14 studies(Zhu et al. 2022) 3 cases of PTEG for
malignant obstruction, all placed under sedation. (Singal et al. 2010)
10 patients treated with PTEG for malignant bowel obstruction. Time from
placement to death median 15 days. Unlike venting gastrostomy, all
patients required suction to maintain resolution of malignent bowel
obstruction symptoms. (Selby et al. 2019) 38 patients treated with perc
transesophageal gastrostomy (PTEG) which was successful in 35/38. Mean
catheter duration 61 days with 5/35 requiring tube exchanges.
(Rotellini-Coltvet et al. 2023)

17 patients treated with PTEG (Udomsawaengsup et al. 2008)

\chapter{Chemoradiation}\label{eso_dcrt}

**This section addresses chemoradiation as primary (definitive) therapy
for esophageal cancer. See also \hyperref[trimodality]{Trimodality
Therapy}

\section{Phase II Studies}\label{phase-ii-studies}

Experience with patients who refuse surgery or are medically unfit:

MD Anderson report of 61 patients out of 622 tri-modality-eligible
patients who refused surgery after cCR. 5-year overall survival was
58\%. 13 developped local recurrence during surveillance, and 12 had
successful salvage esophagectomy (Taketa et al. 2012). The same group
compared those who underwent surgery vs definitive chemoradiation in a
propensity-matched fashion(Taketa et al. 2013). Irish study of 56
patients aged 70 or older treated with neoadjuvant chemoradiation
followed by selective surgery. Median survival was 28 months overall, 47
months for those with cCR, 61 months for primary resection, 46 months
for cCRs who did not undergo resection, and 29 months for those with
salvage esophagectomy(Furlong et al. 2013).

Castoro(Castoro et al. 2013)

preSANO(Chirieac et al. 2005) Clinical Response evaluation after chemoRT
for esophageal cancer with PET and EGD.

RTOG 94-05 clinical trial (Minsky et al. 2002)

\section{ChemoRT vs Trimodality
therapy}\label{chemort-vs-trimodality-therapy}

The sensitivity of squamous cell carcinoma of the esophagus to
chemoradiation has raised the question whether

Stahl Locally advanced squamous cell carcinoma randomized to induction
chemotherapy (cisplatin, etopiside, 5FU with leuocovrin) followed by
chemoradiation (4000cGy with concurrent ciplatin and etopiside) followed
by surgery compared with induction chemotherapy followed by
chemoradiation (6400cGy with concurrent cisplatin and etopiside).(Stahl
et al. 2005) progression-free survival was better in the trimodality
group (64.3\% vs 40.7\%) Treatment-related morality was substantial in
the surgery arm (13\% vs 4\%). This would be considered an excessive
rate of operative mortality by modern standards. Unsurprisingly, there
was no difference in overall survival between groups, in part because
the surgical group had an excess 9\% mortality rate from treatment.
Two-year survival in the surgery arm was 40\% vs 35\% in the definitive
chemoradiation arm.

In the French FFCD trial, 444 patients with carcinoma of the esophagus
(90\% squamous cell) were treated with two cycles of 5-FU and cisplatin
with concurrent radiation.(Bedenne et al. 2007) Patients with a partial
or complete clinical response to chemoradiation were randomized to
either surgery or a boost of radiation. Patients who did not respond to
chemoradiation were treated with surgery and were eliminated from the
study. Only 259 of the original 444 patients (59\%) went on to
randomization, with the remainder (those not responding to
chemoradiation) treated with surgery. Of the randomized group, median
survival was 17.7months in the surgery arm versus 19.3months in the
definitive chemoradiation arm. Like the Stahl study, treatment-related
mortality in the surgical arm was high (9\% versus 1\%).

\section{Active Surveillance}\label{active-surveillance-1}

EGD is poor predictor of pCR (Sarkaria et al. 2009)

\chapter{Salvage esophagectomy}\label{salvage-esophagectomy}

(Markar et al. 2014)

(Swisher et al. 2002)

SANO clinical trial (Overtoom et al. 2024) found similar outcomes for
esophagectomy \textless12mo vs \textgreater12mo after preoperative
therapy, but higher rate of respiratory complications in the delayed
surgery group.

\chapter{Nutrition}\label{eso_nutrition}

\section{Jejunostomy tubes}\label{jejunostomy-tubes-1}

Jejunostomy tubes are the conventional method used for nutritional
support in patients with esophageal cancer who may under surgery in the
future. Jejunostomy avoids potential damage to the (future) gastric
conduit most commonly used for reconstruction of the esophagus.

\section{Gastrostomy tubes}\label{gastrostomy-tubes}

PEG in esophageal cancer (Margolis et al. 2003). PEG placement planned
in 119/179 patients with new diagnosis of esophageal cancer. Successful
in 103/119. No incidence of tumor inoculation metastasis noted. 61
patients underwent surgery and none had difficulty with gastrostomy
closure. PEG patients were more likely to complete chemoRT and had
better survival at 12 months.

Case report of PEG causing injury to right gastroepiploic artery
(Ohnmacht et al. 2006)

\chapter{Esophageal Surgery}\label{esophagectomy}

Three general approaches exist for surgical therapy.

Trans-thoracic or Ivor Lewis esophagectomy(Visbal et al. 2001) removes
the intrathoracic portion of the esophagus and constructs an anastomosis
within the chest. The approach include an abdominal phase, during which
an esophageal substitute is constructed (usually from stomach). A
thoracic phase then removes the intrathoracic esophagus and constructs
an anastomosis within the chest cavity.

A McKeown esophagectomy utilizes three surgical fields: abdomen, right
chest, and neck. The right chest approach allows dissection of
peri-esophageal lymph nodes, and the cervical incision allows removal of
the total esophagus.(McKeown 1976) This approach is useful for tumors
which involve the proximal thoracic esophagus, to ensure a negative
margin. The cervical anastomosis carries a higher risk of anastomotic
leak than a thoracic anastomosis, although the morbidity of a cervical
anastomosis leak is less serious than that of a leak of a thoracic
anastomosis.

A transhiatal esohpagectomy approaches the esophagus from the abdomen
through the hiatus and from neck. By blunt dissection the esophagus is
freed up without the need for thoracotomy. An esophageal substute is
then brought from the abdomen to the neck through the mediastinum.(M. B.
Orringer and Sloan 1978) (Mark B. Orringer et al. 2007) The operation is
designed to avoid the pulmonary toxicity of the right chest approach. On
the other hand, the blunt nature of the mediastinal dissection means
that fewer lymph nodes are harvested than with a trans-thoracic
approach.

Randomized trial of transthoracic esophagectomy with extended lymph node
dissection versus transhiatal esohpagectomy showed fewer pulmonary
complications with the transhiatal approach. (Hulscher et al. 2002)
Fewer lymph nodes were havested with a transhiatal appraoch. A post-hoc
analysis showed that among patients with 1-8 positive lymph nodes,
survival with improved with the extended lymph node dissection.(Omloo et
al. 2007)

Minimally-invasive approaches to esophagectomy are now common, with
evidence for less perioperative morbidity than an open approach (Biere
et al. 2012) (Zhou et al. 2015)

Randomized trial of a hybrid MIE (with laparoscopy and thoracotomy) was
associated with lower postoperative complications than open
esophagectomy (Christophe Mariette et al. 2019)

High volume centers have lower mortality for esophageal cancer than
low-volume centers. (Birkmeyer et al. 2003) (Wouters et al. 2009)

\subsection{Trans-thoracic vs Transhiatal
Esophagectomy}\label{trans-thoracic-vs-transhiatal-esophagectomy}

Dutch randomized trial (n=262) cervical vs thoracic anastomosis.(Workum
et al. 2021) Thoracic anastomosis associated with lower leak rate (12\%
vs 34\%) and lower rate of recurrent laryngeal nerve injury (0\% vs
7.3\% ) and better quality of life (dysphagia, choking while swallowing,
and talking)

\subsection{GE Junction
Adenocarcinoma}\label{ge-junction-adenocarcinoma}

Siwert III lesions are considered gastric cancers (Rusch 2004) (Siewert,
Stein, and Feith 2006)

Laparoscopy may be helpful in Siewert III tumors (Graaf et al. 2007)

\subsection{Preoperative Evaluation}\label{preoperative-evaluation-2}

Dysphagia can be scored according to Mellow et al (Mellow and Pinkas
1985):

\begin{itemize}
\tightlist
\item
  0 No dysphagia
\item
  1 Dysphagia to normal solids
\item
  2 Dysphagia to soft solids (ground beef, poultry,fish)
\item
  3 Dysphagia to solids and liquids
\item
  4 Inability to swallow saliva
\end{itemize}

\section{Minimally-invasive
Esophagectomy}\label{minimally-invasive-esophagectomy}

Higher lymph node yield with MIE vs open approach (Kalff, Fransen, et
al. 2022)

\section{Early Recovery Pathways}\label{early-recovery-pathways}

\href{https://www.ncbi.nlm.nih.gov/pmc/articles/PMC7098920/}{ERAS
Society Guidlines}

\section{Colon Interposition}\label{colon-interposition}

Left colon technique based upon blood supply (Peters et al. 1995)

\chapter{Survivorship}\label{eso_survivorship}

\section{Nutritional consequences of
esophagectomy}\label{nutritional-consequences-of-esophagectomy}

\subsection{Vitamin D deficiency}\label{vitamin-d-deficiency}

Vitamin D deficiency is defined as serum 25(OH)D levels below 20ng/mL

\begin{itemize}
\tightlist
\item
  Replacement with 2000 IU Vitamin D3 daily (Khan and Fabian 2010)
\end{itemize}

Vitamin D insufficiency is defined as serum 25(OH)D level 20-29ng/mL

\begin{itemize}
\tightlist
\item
  Replacement with 1000-2000 IU Vitamin D3 daily
\end{itemize}

\href{https://ods.od.nih.gov/factsheets/vitamind-healthprofessional/}{NIH
ODS Vitamin D info}

(Baker et al. 2016)

Weight loss (Martin and Lagergren 2009) (Ouattara et al. 2012)

\section{Cardiac toxicity of
radiation}\label{cardiac-toxicity-of-radiation}

(Beukema et al. 2015) (Frandsen et al. 2015) (Gharzai et al. 2016)

\chapter{Surveillance}\label{eso_surveillance}

\textbf{T1a treated with endoscopic resection} EGD every 3 mo for first
year, then every 6 months for second year, then annually(Shaheen et al.
2016)

\textbf{T1b treated with endoscopic resection} EGD every 3 mon for first
year, then every 4-6 months for seond year, then annually CT
chest/abdomen every 12 months for up to 3 years (as clinically
indicated)

\textbf{T1b treated with esophagectomy} EGD every 3-6 months for first 2
years, then annually for 3 more years. CT every 6-9 months for first 2
years, then annually up to 5 years.

\textbf{Stage II or III treated with chemoradiation.}

These patients are at risk for local recurrence (Sudo et al. 2014) and
some may be candidates for salvage esophagectomy. Most relapses (95\%)
occur within 24 months. See also (Taketa et al. 2014)

\textbf{Locally-advanced treated with trimodality therapy}

Local/regional relapses are uncommon. (Dorth et al. 2014) (Oppedijk et
al. 2014) (Sudo et al. 2013) =\textgreater{} NCCN does not recommend
EGD. 90\% of relapses occur within 36 months of surgery.

CT every 6 months up to 2 years (if patient is a candidate for
additional curative-intent therapy)

\subsection{Recurrence Profile}\label{recurrence-profile}

Dutch cancer registry (Kalff, Henckens, et al. 2022)

\part{Gastric Cancer}

\chapter*{(PART*) Gastric Cancer}\label{part-gastric-cancer}
\addcontentsline{toc}{chapter}{(PART*) Gastric Cancer}

\markboth{(PART*) Gastric Cancer}{(PART*) Gastric Cancer}

\chapter*{Gastric Ca SCORE}\label{gastric-ca-score}
\addcontentsline{toc}{chapter}{Gastric Ca SCORE}

\markboth{Gastric Ca SCORE}{Gastric Ca SCORE}

\href{http://portal.surgicalcore.org/modulecontent.aspx?id=128003}{Gastric
Cancer SCORE}

\textbf{Anatomy}

\begin{itemize}
\tightlist
\item
  Lymph node stations
\item
  D1 vs D2 vs D3 nodal basins
\end{itemize}

\textbf{Epidemiology}

\begin{itemize}
\tightlist
\item
  Worldwide distribution
\item
  Risk factors
\end{itemize}

\textbf{Diagnosis and Staging}

\begin{itemize}
\tightlist
\item
  Nodal staging number vs location
\item
  Role of PET scan
\item
  Role of laparoscopy
\end{itemize}

\textbf{Operative Treatment}

\begin{itemize}
\tightlist
\item
  Reconstruction options
\end{itemize}

\textbf{Non-operative Treatments}

\begin{itemize}
\tightlist
\item
  Palliative chemotherapy
\item
  Stents
\end{itemize}

\href{http://portal.surgicalcore.org/modulecontent.aspx?id=146287}{Gastrectomy
SCORE}

\textbf{Operative Anatomy}

\textbf{Preoperative Preparation}

\textbf{Operative Details}

\textbf{Complications}

\begin{itemize}
\tightlist
\item
  Anastomotic leak
\item
  Bleeding
\item
  Obstruction
\item
  Afferent Loop Syndrome
\item
  Dumping
\item
  Alkaline reflux gastritis
\end{itemize}

\href{http://portal.surgicalcore.org/modulecontent.aspx?id=136240}{GIST
SCORE}

\textbf{GI stromal tumors}

\begin{itemize}
\tightlist
\item
  Pathology: Cell of origin/IHC
\item
  Operative treatment
\end{itemize}

\textbf{Gastric carciniods}

\begin{itemize}
\tightlist
\item
  Pathology - Cell of origin
\item
  Type I
\item
  Type II
\item
  Type III
\item
  Role for endoscopic therapy
\item
  Surgery

  \begin{itemize}
  \tightlist
  \item
    Tumor resection
  \item
    Antrectomy
  \end{itemize}
\end{itemize}

\textbf{Gastric lymphoma}

\begin{itemize}
\tightlist
\item
  Low-grade vs high-grade MALT vs non-MALT
\item
  Treatment: operative/non-op
\item
  Indications for surgery
\end{itemize}

\textbf{Gastric Polyps}

\begin{itemize}
\tightlist
\item
  Indications for resection
\item
  Non-operative treatment of hyperplastic polyps
\item
  SMAD4 mutations
\end{itemize}

\chapter{Gastrectomy}\label{gastrectomy}

\section{Proximal gastrectomy}\label{proximal-gastrectomy}

Proximal gastrectomy for small proximal tumors can be performed with a
dual-tract reconstruction to take advantage of the reservoir capacity of
the distal stomach without the risk of gastroesophageal reflux.

\includegraphics[width=0.5\linewidth,height=0.5\textheight]{index_files/mediabag/double-tract-reconst.jpg}
\#\# Total vs Subtotal Gastrectomy

Italy: Randomized trial total vs subtotal gastrectomy (Bozzetti et al.
1999). No difference in survival

\chapter{Hereditary Diffuse Gastric
Cancer}\label{hereditary-diffuse-gastric-cancer}

HDGC is a syndrome characterized by early-onset diffuse gastric cancer
and increased risk of lobular breast cancer.

Hereditary diffuse gastric cancer was first reported in a New Zealand
Maori cohort. The genetic cause was subsequently found to the a deletion
in the CDH1 gene.

\subsection{Penetrance estimates for lifetime risk of
DGC}\label{penetrance-estimates-for-lifetime-risk-of-dgc}

The original estimates of lifetime risk of gastric cancer in those with
pathogenic CDH1 mutations were made from a population of patients with
strong familiy histories of DGC(Hansford et al. 2015), in this group,
the cumulative risk of gastric cancer by 80 years of age was estimated
to be to be 70\% in men and 56\% in women.

The lifetime risk of gastric cancer in more modern cohorts(Roberts et
al. 2019) (Xicola et al. 2019), in whom a minority of patients have a
family history of gastric cancer, are much lower: 42\% in men and 33\%
in women

\section{Surveillance for CDH1
carriers}\label{surveillance-for-cdh1-carriers}

Cancer surveillance as an alternative to prophylactic total gastrectomy
in hereditary diffuse gastric cancer: a prospective cohort study Bilal
Asif\emph{, Amber Leila Sarvestani}, Lauren A Gamble, Sarah G
Samaranayake, Amber L Famiglietti, Grace-Ann Fasaye, Martha Quezado,
Markku Miettinen, Louis Korman, Christopher Koh, Theo Heller, Jeremy L
Davis Summary Background Loss of function variants in CDH1 are the most
frequent cause of hereditary diffuse gastric cancer. Endoscopy is
regarded as insufficient for early detection due to the infiltrative
phenotype of diffuse-type cancers. Microscopic foci of invasive signet
ring cells are pathognomonic of CDH1 and precede development of diffuse
gastric cancer. We aimed to assess the safety and effectiveness of
endoscopy for cancer interception in individuals with germline CDH1
variants, particularly in those who declined prophylactic total
gastrectomy. Methods In this prospective cohort study, we included
asymptomatic patients aged 2 years or older with pathogenic or likely
pathogenic germline CDH1 variants who underwent endoscopic screening and
surveillance at the National Institutes of Health (Bethesda, MD, USA) as
part of a natural history study of hereditary gastric cancers
(NCT03030404). Endoscopy was done with non-targeted biopsies and one or
more targeted biopsy and assessment of focal lesions. Endoscopy
findings, pathological data, personal and family cancer history, and
demographics were recorded. Procedural morbidity, gastric cancer
detection by endoscopy and gastrectomy, and cancer-specific events were
assessed. Screening was defined as the initial endoscopy and all
subsequent endoscopies were considered surveillance; follow-up endoscopy
was at 6 to 12 months. The primary aim was to determine effectiveness of
endoscopic surveillance for detection of gastric signet ring cell
carcinoma. Findings Between Jan 25, 2017, and Dec 12, 2021, 270 patients
(median age 46·6 years {[}IQR 36·5--59·8{]}, 173 {[}64\%{]} female
participants, 97 {[}36\%{]} male participants; 250 {[}93\%{]} were
non-Hispanic White, eight {[}3\%{]} were multiracial, four {[}2\%{]}
were non-Hispanic Black, three {[}1\%{]} were Hispanic, two {[}1\%{]}
were Asian, and one {[}\textless1\%{]} was American Indian or Alaskan
Native) with germline CDH1 variants were screened, in whom 467
endoscopies were done as of data cutoff (April 30, 2022). 213 (79\%) of
270 patients had a family history of gastric cancer, and 176 (65\%)
reported a family history of breast cancer. Median follow-up was 31·1
months (IQR 17·1--42·1). 38 803 total gastric biopsy samples were
obtained, of which 1163 (3\%) were positive for invasive signet ring
cell carcinoma. Signet ring cell carcinoma was detected in 76 (63\%) of
120 patients who had two or more surveillance endoscopies, of whom 74
had occult cancer detected; the remaining two individuals developed
focal ulcerations each corresponding to pT3N0 stage carcinoma. 98 (36\%)
of 270 patients proceeded to prophylactic total gastrectomy. Among
patients who had a prophylactic total gastrectomy after an endoscopy
with biopsy samples negative for cancer (42 {[}43\%{]} of 98),
multifocal stage IA gastric carcinoma was detected in 39 (93\%). Two
(1\%) participants died during follow-up, one due to metastatic lobular
breast cancer and the other due to underlying cerebrovascular disease,
and no participants were diagnosed with advanced stage (III or IV)
cancer during follow-up. Interpretation In our cohort, endoscopic cancer
surveillance was an acceptable alternative to surgery in individuals
with CDH1 variants who declined total gastrectomy. The low rate of
incident tumours (\textgreater T1a) suggests that\ldots{}

Advanced (T3) tumors found during endoscopic surveillance all had
mucosal abnormalities found at EGD

Consensus guidelines (2020)(Blair et al. 2020) Recommend surveillance in
whom family history of gastric cancer is weak

REcommend 28-30 random biopsies

revirew (Pilonis et al. 2021) Ann rev med 72:263 2022

guidelines (Post et al. 2015) j med genet 52:361

total gast (Chen et al. 2011)ann surg oncol 2011:18. kingham

Cambridge: Endoscopic surveillance in HDGC. 145 patients from 76
families with a family history of gastric cancer and CDH1 mutations (see
blair and vanderpost). Pathogenic variasnt of CDH1 in 92 (63\%).
Prophylatic total gastrecromy done in 36 (25\%)

During surveilance 58/145 (58\%) were found to have invasive signet ring
cell carcinoma.

Of those with CDH1 mutations, 53\% had (49/91) had invasive signet ring
cell carcinoma diagnosed

Invasive signet ring cell carcinoma detected in 6/41 (15\%) with no CDH1
pathogenic mutations

Total gastrectom in 36 of whom 32 (89\%) had signet ring cell carcinoma
on pathology.

Prophylactic total gastrectomy with CDH1 mutation in 28 patients, of
whom 26 (93\%) had signet ring cell carcinoma.

6/13 invasive signet ring cell carcinoma found on random biopsies later
developped visible lesions and underwent gastrectomy (see addpneix p8)
progressed to more advanced

\chapter{Gastric GIST}\label{gastric-gist}

\section{Genetics}\label{genetics}

\begin{itemize}
\tightlist
\item
  Mutation of KIT receptor tyrosine kinase in 80\%

  \begin{itemize}
  \tightlist
  \item
    KIT exon 9 mutations have lower response rate to imatinib and
    shorter progression-free survival than exon 11 mutations (especially
    at dose of 400mg daily)
  \end{itemize}
\item
  Mutation in PDGRFA receptor tyrosine kinase in 5-10\%. Most mutations
  respond to imatinib except D842V (which may respond to avapritinib)
\item
  No mutation of KIT or PDGFRA in 10-15\%. Most will have functional
  inactivation of SDH complex and evidenced by lack of SDHB on IHC.
  These patients should be tested for germline SDH mutations. These
  tumors tend to occur in the stomach in younger patients.
\end{itemize}

Tyrosine kinase inhibitors:

\begin{itemize}
\tightlist
\item
  imatinib - first line therapy
\item
  Sunitinib treatment is indicated for patients with imatinib-resistant
  tumors or imatinib intolerance.
\item
  Regorafenib is indicated for patients with disease progression on
  imatinib or sunitinib.
\item
  Ripretinib is indicated for patients who have received prior treatment
  with 3 or more kinase inhibitors, including imatinib.
\end{itemize}

\section{Surgery}\label{surgery}

Surgical resection ideally involved a negative histologic margin. A wide
resection margin is not required, nor is lymph node dissection unless
nodes are clinically involved. (\textbf{dematteo51?})

\section{Positive margins after
resection}\label{positive-margins-after-resection}

ACOSOG working group (\textbf{mccarter53?}) studies 819 patients
undergoing resection of GIST under Z9000 and Z9001 clinical trials. 72
patients (8.8\%) had an R1 resection, of whom 21 had intraperitoneal
rupture or bleed. Factors associated with R1 resection included tumor
size \textgreater10cm, rectum location, and tumor rupture. Disease
recurrence in patients with R1 resections was primarily due to
associated tumor rupture. There was no difference in recurrence-free
survival for patients underoing R1 vs R0 resection.

Norweigan report (Hølmebakk et al. 2019) of 401 patients wtih GIST of
whom 47had R1 resection and 52 with tumor recurrence. In patients
without rupture, there was no difference in 5-year recurrence-free
survival for R1 vs R0 resection. In multivariable analysis, tumor
rupture but not R1 resection was associates with recurrence.

\chapter{Locally-Advanced Gastric}\label{locally-advanced-gastric}

Locally-advanced gastric cancer (T3 or N\textsuperscript{+}) is
generally treated with some form of adjuvant therapy, which has been
shown to improve upon the outcomes with surgery alone.

\section{Preoperative Chemotherapy}\label{preoperative-chemotherapy}

\textbf{FLOT} chemotherapy (Al-Batran et al. 2019) with 5FU, oxaliplatin
and Taxotere superior for neoadjuvant chemotherapy compared with
epirubicin, oxaliplatin and 5FU (MAGIC regimen).

Matterhorn trial evaluated the addition of anti-PD-L1 MAb
\textbf{durvalumab} to FLOT for stage I-IVa adenocarcinoma of stomach
and GE junction. Both groups received 4 cycles of FLOT every 15 days.
Durvalumab group received 2 doses preop and 2 doses postop with FLOT in
addition to 10 monthly doses of durvalumab. Surgery was performed 4-8
weeks after the last dose of chemotherapy. Longer disease-free survival
and trend to longer overall survival with the addition of durvalumab
without additional toxicity.(Janjigian et al. 2025)

MAGIC study randomized 503 patients to perioperative `sandwich'
chemotherapy consisting of epirubicin, cisplatin, and 5-FU versus
surgery alone. In the perioperative chemotherapy group, 4 cycles were
administered prior to surgery, and 4 cycles afterwards. Tumors of the
esophagus or gastroesophageal junction comprised 26\% of the study
population. While over 90\% of patients assigned to the chemotherapy arm
completed their preoperative chemotherapy, only 66\% completed their
postoperative therapy. Survival at 5 years was 36\% in the perioperative
chemotherapy group, compared with 24\% in the surgery group
(p\textless0.001).(Cunningham et al. 2006)

ACCORD-07 clinical trial (Ychou et al. 2011): surgery +/- neoadjuvant
cisplatin+5FU. Improved R0 resection rate but no difference in survival,
perhaps due to poor accrual.

FFCD trial randomized patients to preoperative chemotherapy with 2 or 3
cycles of cisplatin and 5-FU versus surgery alone. Tumors of the lower
esophagus or gastroesophageal junction comprised 75\% of the study
population. Survival at 5 years was longer in the chemotherapy group
(38\%) versus 24\% in the surgery alone group (p=0.02).(Ychou et al.
2011)

RESONANCE-II trial of Oxaliplatin + S-1 (SOX) (Wang et al. 2021)

\section{Postoperative chemotherapy}\label{postoperative-chemotherapy}

CLASSIC clinical trial randomized 1033 patients with stage II or III
gastric cancer after D2 gastrectomy to 6 months of adjuvant chemotherapy
with capecitabine and oxaliplatinversus vs surgery alone. Three-year
survival was improved in the chemotherapy group (74\% \emph{v}
59\%).(Bang et al. 2012) (Noh et al. 2014)

\section{Postoperative
chemoradiation}\label{postoperative-chemoradiation-1}

Intergroup 0116 trial (Macdonald et al. 2001) (Smalley et al. 2012)
Surgery +/- postoperative chemoradiation. Surgical quality control was
poor, as 90\% were treated a limited (R1) lymph node dissection.
Long-term followup, however (Smalley et al. 2012) showed a persistent
benefit of postoperative chemoradiation.

ARTIST trial 450 patients treated with a D1 \(\alpha\) gastrectomy were
randomized to adjuvant capecitibine and cisplatin versus chemoradiation
consisting of two cycles of capcitabine/oxaliplatin followed by
chemoradiation followed by chemotherapy. Overall 3- year survival did
not differ between groups (78.2\% vs 74.2\% p =0.86). A post-hoc
analysis of patients with positive nodes showed a beneficial effect of
chemoradiation (77.5\% \emph{v} 72.3\% p=0.365).(Lee et al. 2012)

CRITICS trial (Cats et al. 2018) treated all patients with preoperative
chemotherapy (ECF or CapeOx) followed by surgery. Postoperative patients
were then randomized between additional chemotherapy versus
chemoradiation. No difference in survival between groups.

CRITICS-II Trial (Slagter et al. 2018) resectable gastric cancer
randomized to neoadjuvant chemo (DOC x4) vs chemo (DOC x2) followed by
ChemoRT (CROSS) vs ChemoRT (CROSS).

\section{Preoperative chemoradiation}\label{preoperative-chemoradiation}

RTOG 9904: Phase II trial of preoperative chemoradiation for gastric
adenocarcinoma(Ajani et al. 2006). Findings of pathologic complete
response (pCR) in 20\%.

TOPGEAR trial: neoajuvant chemotherapy \(\pm\) subsequent radiation
therapy: chemo+RT had better pCR rate (17\% vs 8\%). No survival
difference between groups.(Leong et al. 2017)

\section{Preop Immunotherapy}\label{preop-immunotherapy}

NEONIPIGA trial: MSI-Hi gastric and GE junction cancers treated with
neoadjuvant nivolumbab (6 cycles) and ipilumumab (2 cycles) over 12
weeks. Among 29 patients that underwent surgery, 59\% pathologic
complete response and all patients had R0 resection.(\textbf{andre288?})

MSKCC: MSI-Hi GI cancers treated with dostarlimab. Clinic response rate:
rectal 50/50 (100\%), colon 28/33 (82\%), gastric 8/14 (57\%), GE
junction 1/3 (33\%), esophageadl 1/3 (0\%) (Cercek et al. 2025)

\chapter{M1 Gastric CA}\label{m1-gastric-ca}

\section{Surgery for M1 disease}\label{surgery-for-m1-disease}

REGATTA trial: patient with a single unresectable factor (liver
metastasis, peritoneal disease, para-aortic lymphdenopathy randomized
between gastrectomy -\textgreater{} chemotherapy vs chemotherapy. No
improvement in survival with gastrectomy. Trial was closed early due to
futility.(Fujitani et al. 2016)

AIO-FLOT3 trial: evaluted FLOT followed by surgical resection in
patients with limited metastatic gastric or GEJ adenocarcinoma. Patients
were stratified: resectable (A), limited metastatic (B), and extensive
metastatic (C).~Arm B patients received at least four cycles of FLOT and
were considered for surgery if restaging indicated a chance for R0
resection of the primary and complete resection of metastatic lesions.
The primary endpoint was overall survival.~Results showed that patients
in arm B (limited metastatic) who underwent neoadjuvant FLOT and surgery
had a median overall survival of~\textbf{31.3 months}, compared
to~\textbf{15.9 months}~for those who did not proceed to surgery,
and~\textbf{10.7 months}~for arm C (extensive metastatic, palliative
approach).~The response rate in arm B was 60\%.(Al-Batran et al. 2017)

\section{Conversion Surgery}\label{conversion-surgery}

Conversion surgery is performed in patients with positive peritoneal
cytology or positive peritoneal disease who, after chemotherapy, are
found to have coversion of their findings to M0

Japan: 115 patients with advanced gastric cancer underwent second-look
laparoscopy. 56 had negative cytology and peritoneal biopsy. 26 patients
had positive cytology and negative peritoneal biopsy, of whom 12 had a
negative second-look laparoscopy and had conversion surgery (CS). 33
patients had positive cytology and peritoneal biopsy. The survival of
patients with CS was significantly longer than those without CS (median
survival time (MST); 41 vs.~11 months, respectively, P \textless{}
0.001). We observed no difference in overall survival between patients
who underwent CS and patients with CY0P0 at the first SL (P = 0.913).
All patients with P1 received chemotherapy. As peritoneal metastasis of
five patients (15.2\%) disappeared by chemotherapy, those patients
underwent the CS (R0). The survival of patients who underwent CS was
significantly longer than those who did not (MST; 31 vs.~10 months,
respectively, P = 0.034).(Nakamura et al. 2019)

\part{Colon Cancer}

\chapter{Colectomy}\label{colectomy}

\section{SCORE}\label{PColectomyObj}

\href{https://www.surgicalcore.org/modulecontent.aspx?id=129670}{Partial
Colectomy CORE}

\textbf{Indications}

\textbf{Operative Anatomy}

\textbf{Preop Prep}

\begin{itemize}
\tightlist
\item
  Antibiotics
\item
  Stoma marking
\item
  Ureteral stents
\end{itemize}

\textbf{Complications}

\begin{itemize}
\tightlist
\item
  Anastomotic leak
\item
  Ureteral injury
\end{itemize}

\section{Early Recovery Protocols}\label{early-recovery-protocols}

ASCRS and SAGES clinical guidelines for enhanced recovery after
colorectal surgery (Irani et al. 2023)

\href{https://fascrs.org/ascrs/media/files/downloads/Clinical\%20Practice\%20Guidelines/clinical_practice_guidelines_for_enhanced_recovery-3.pdf}{ASCRS
Guidelines for Early Recovery}

Same-day discharge after colectomy or stoma reversal (Paradis et al.
2024)

\section{Anesthetic Management}\label{anesthetic-management}

Restrictive vs Liberal Fluid Managment (RELIEF Trial)
(\textbf{myles378?}) 3000 patients undergoing elective major surgery
(43\% colorectal, 64\% cancer, no liver resection) randomized to
restrictive intraoperative fluid management (aim net zero fluid balance
in the OR) vs traditional fluid management. Liberal group averaged 3
liters of intraop fluid and 3 liters first 24 hours postop vs
restrictive group with 1.8 and 1.9 liters respectively. No difference in
mortality or major disability at 1 year postop. More acute kidney injury
in restrictive group. More surgical site infections (17\% vs 14\%) in
the restrictive group. Bottom line: No detectable benefit to restrictive
fluid strategy

\section{Extended Node dissection}\label{extended-node-dissection}

Short-term outcomes of complete mesocolic excision versus D2 dissection
in patients undergoing laparoscopic colectomy for right colon cancer
(RELARC): a randomised, controlled, phase 3, superiority tria

Short-term outcomes of a multicentre randomized clinical trial comparing
D2 versus D3 lymph node dissection for colonic cancer (COLD trial).
Karachun A, Panaiotti L, Chernikovskiy I, Achkasov S, Gevorkyan Y,
Savanovich N, Sharygin G, Markushin L, Sushkov O, Aleshin D, Shakhmatov
D, Nazarov I, Muratov I, Maynovskaya O, Olkina A, Lankov T, Ovchinnikova
T, Kharagezov D, Kaymakchi D, Milakin A, Petrov A. Br J Surg. 2020
Apr;107(5):499-508. doi: 10.1002/bjs.11387. Epub 2019 Dec 24. PMID:
31872869 Clinical Trial.

\part{Rectal Cancer}

\chapter{Rectal Cancer Staging}\label{rectal-cancer-staging}

\section{MRI staging}\label{mri-staging}

Rectal cancer is preferentially evaluated with MRI. At Atrium, this is
ordered as ``MRI Pelvis without Contrast Rectal Protocol''

Key findings

\begin{itemize}
\tightlist
\item
  Circumferential resection margin
\item
  Extra mesorectal vascular invasion
\item
  Mesorectal lymph nodes
\item
  Extra-mesorectal lymph nodes
\end{itemize}

MRI has supplanted endoscopic ultrasound (EUS) as the study of choice
for staging rectal cancer.

MERCURY group trial demonstrated the predictive value of MRI for rectal
cancer(Taylor et al. 2014) (MERCURY Study Group 2006)

\section{EUS}\label{eus}

Endoscopic ultrasound has now been replaced by MRI for initial staging,
in part due to interoperater variability

\chapter{Rectal Ca Surgery}\label{rectal-ca-surgery}

\section{Trans-abdominal Rectal
Surgery}\label{trans-abdominal-rectal-surgery}

Trans-abdominal procedures include

\begin{itemize}
\tightlist
\item
  Low anterior resection: Removal of the rectum with an anastomosis
  between colon and distal rectum
\item
  Hartmann resection: Removal of the rectum with end colostomy. The
  rectal stump is stapled. The anal sphincters are left in situ
\item
  Abdominoperineal resection: Removal of rectum and anus from both an
  abdominal and perineal approach. The anal sphincters are removed.
\end{itemize}

\section{Low Anterior Resection}\label{low-anterior-resection}

\subsection{Total Mesorectal Excision}\label{total-mesorectal-excision}

Importance of total mesorectal excision was championed by Bill Heald
(Heald and Ryall 1986) who emphasized sharp dissection of the mesorecum
outside of the visceral fascial envelope. In addition he advocated
dissection of the totality of the mesorectum distal to the tumor to
avoid leaving behind nodes within the mesorectum distal to the
tumor(Quirke et al. 1986) (Nagtegaal and Quirke 2008) (Paty et al.
1994).

\chapter{Rectal Adjuvant Therapy}\label{RectalAdjuvant}

\section{Trans-abdominal Rectal
Surgery}\label{trans-abdominal-rectal-surgery-1}

Trans-abdominal procedures include

\begin{itemize}
\tightlist
\item
  Low anterior resection: Removal of the rectum with an anastomosis
  between colon and distal rectum
\item
  Hartmann resection: Removal of the rectum with end colostomy. The
  rectal stump is stapled. The anal sphincters are left in situ
\item
  Abdominoperineal resection: Removal of rectum and anus from both an
  abdominal and perineal approach. The anal sphincters are removed.
\end{itemize}

\section{Rectal Adjuvant Therapy}\label{rectal-adjuvant-therapy}

Surgery as sole treatment for rectal cancer is associated with an
unacceptable rate of local recurrence. Accordingly, adjunctive
stratagies were employed to reduce risk of local recurrence.

\section{\texorpdfstring{Surgery \(\Rightarrow\)
CRT}{Surgery \textbackslash Rightarrow CRT}}\label{surgery-rightarrow-crt}

The original approach to adjuvant therapy for rectal cancer was surgery
followed by chemoradiation.

GI Tumor Study Group (GITSG) (Gastrointestinal Tumor Study Group 1985)
showed a reduction in local recurrence with postoperative chemotherapy +
radiation

NIH Consensus Conference in 1990 recommended adjuvant chemoradation
after surgery

\section{\texorpdfstring{CRT \(\Rightarrow\)
Surgery}{CRT \textbackslash Rightarrow Surgery}}\label{crt-rightarrow-surgery}

The next innovation to adjuvant therapy for rectal cancer was
chemoradiation prior to surgery, followed by adjuvant chemotherapy.

German Rectal Cancer Study (AIO-94) compared preop and postop
chemoradiation in locally-advanced rectal cancer (Sauer et al. 2001).
Among 823 patients, local recurrence was 6\% in the preoperative group
vs 13\% in the postop group (p=0.0006). Overall survival was not
different, but preop chemoradiation was less toxic. Pathologic complete
response rate in the preoperative group as 8\%.

NSABP R-03 clinical trial (Roh et al. 2009) randomized 267 patients with
T3 or T4 or node-positive rectal cancer to preop vs postop
chemoradiation. Disease-free survival was better in the preoperative
group, wiht no difference in overall survival. Pathologic complete
response rate in the preop group was 15\%. The trial did not meets its
accrual targets.

\section{TME}\label{tme}

Total Mesorectal Excision (TME) (Heald and Ryall 1986) was introduced to
reduce risk of local recurrence by removal of the mesorectum and its
lymph nodes. 115 rectal cancer patients underwent surgery over 7.5
years, of whom 69 had an anastomosis below 5cm. At median followup of
4.2 years, there were 3 pelvic recurrences and no staple-line
recurrences.

\section{TME + ChemoRT}\label{tme-chemort}

The Dutch trial (CKVO 95-04) found a reduction in local recurrence form
10.9\% to 5.6\% with preoperative radiation with no difference in
survival(Kapiteijn et al. 2001). A later update(Gijn et al. 2011)

The Medical Research Council C07 trial compared preoperative
short-course radiation with selective post-operative chemoradiation and
found no difference in local recurrence (4.4\% vs 10.6\%) and no
difference in overall survival.(Sebag-Montefiore et al. 2009)

\section{Short-course RT}\label{short-course-rt}

An alternative to preoperative chemoradiation over 6 weeks is to
administer radiation alone properative over a five-day period.

Swedish Rectal Cancer Trial compared surgery alone with preoperative
short-course therapy consisting of 5 doses of 500cGy of radiation
without chemotherapy administered in one week prior to surgery. Local
recurrence was 9\% in the therapy group vs 26\% in the control group,
with an improvement in overall survival of 38\% vs 80\%(Swedish Rectal
Cancer Trial et al. 1997). Of note, this trial was performed in the era
prior to the widespread use of total mesorectal excision.

Stockholm III trial showed that short-course radiation therapy performed
4-8 weeks prior to surgery resulted in improved rates of pathologic
complete response (12\% vs 2\%) compared with short-course radiation
therapy performed the week prior to surgery (Pettersson et al. 2015)

\section{Total Neoadjuvant}\label{total-neoadjuvant}

A more modern approach has been to administer chemotherapy in addition
to chemoradiation prior to surgery, Total Neoadjuvant Therapy (TNT). The
rationale is to increase the rates of completion of adjuvant
chemotherapy, noting that many patients do no receive postoperative
adjuvant chemotherapy due to delayed recovery from surgery.

RAPIDO trial randomized 920 patients with T4 or node-positive disease to
long-course chemoradiation followed by surgery vs short-course radiation
followed by chemotherapy and surgery. The pCR rate was significantly
higher in the short course/chemotherapy/surgery group (28\% vs 14\%) and
disease-specific surival at 3years was higher (30\% vs 24\%).(Bahadoer
et al. 2021) (Valk et al. 2020)

PRODIGE 23 was designed to reduce distant metastasis and improve
survival using TNT and randomized 461 patients with T3 or T4 rectal
cancers to long-course radiation followed by surgery (standard therapy)
vs induction chemotherapy with FOLFIRINOX followed by long-course
radiation (TNT) followed by surgery. Rectal surgery with TME was
mandatory. In both arms, adjuvant chemotherapy was administered after
surgery. Up-front induction chemotherapy followed by chemoradiation was
associated with increased 3-year survival (76\% vs 69\%) and an increase
in rate of pathologic complete response of 28\% vs 12\%. pCR rate was
28\% with TNT vs 12\% with standard therapy. (Conroy et al. 2021)

STELLAR trial (Jin et al. 2022) TNT with short-course radiation (5Gy x 5
days) + CAPEOX chemotherapy vs long-course chemoRT with capecitabine. No
differences in relapse-free survival or metastasis free survival.
Overall survival better with TNT (p=0.03). On subgroup analysis, no
group appeared to have more benefit. pCR rate 17\% with TNT vs 13.9\%
with conventional chemoradiation. Positive margins (R1) occurred in
8.5\% of TNT patients vs 12.5\% in the conventional chemoRT group.

OPRA clinical trial was designed to optimize organ preservation
(Garcia-Aguilar et al. 2022). 324 rectal cancer patients staged with MRI
randomized to INCT (induction chemo followed by chemoRT) vs CNCT
(Chemoradiation followed by chemotherapy). Patients with a response were
offered watch and wait. Patients without a response were treated with
surgery. No difference in disease-free survival or overall survival or
metastasis-free survival. 304 patients were restaged and only 26\% were
recommended to have surgery. Among 225 patients in watch and wait,
somewhat more patients with INCT had recurrences (40\% of 105 = 42
patients vs 27\% of 102 =32 patients). More organ preservation at 3
years with CNCT (60\% CNCT vs 47\% INCT).See also (Smith et al. 2015)

\section{Selective Adjuvant}\label{selective-adjuvant}

Selective adjuvant therapy approaches reserves adjuvant therapy for
high-risk patients, and treats low-risk patients with surgery alone.

MERCURY study group has examined selective approaches to adjuvant
chemoradiation in low-risk patients with rectal cancer. (Taylor et al.
2011), (Strassburg et al. 2011)

German OCUM group (Kreis et al. 2016) (Ruppert et al. 2018)

Canadian {[}Quicksilver Trial{]}
(https://jamanetwork.com/journals/jamaoncology/fullarticle/2730134)

\section{Neoadjuvant ImmunoTx}\label{neoadjuvant-immunotx}

Mismatch repair protein deficient (MSI-high) colorectal cancer (due to
either Lynch Syndrome or BRAF-1 mutation) responds pooly to chemotherapy
but responds well to immunotherapy with PD-L1 blockade

MSKCC trial of neoadjuvant PD-L1 blockade with 6 months of dostarlimab
for 12 patients with MMR-deficient (MSI-high) rectal cancer resulted in
100\% clinical response rate (Cercek et al. 2022). No patients were
subsequently treated with either chemoradiation or surgery as originally
planned.

PICC trial (Hu et al. 2022) Chinese trial of PD-L1 blockade with
neoadjuvant toripalimab + colecoxib vs toripalimab for 3 months in 34
patients with MMR-deficient locally-advanced (T3/4 or N+) colorectal
cancer. All patients were then treated with surgery. Pathologic complete
response in 88\% of dual-therapy patients and 65\% of monotherapy. Grade
3 toxicity in 1/34 patients. Previous neoadjuvant chemo in 25\%. Rectal
cancer in 18\%.

NICHE trial (Chalabi et al. 2024) 115 patients with dMMR
locally-advanced colon cancer treated with a neoadjuvant immunotherapy.
Pathologic complete response in 68\%. No recurrences with a median
followo up of 26 months.

\section{Neoadjuvant Chemotherapy}\label{neoadjuvant-chemotherapy}

The PROSPECT clinical trial (Schrag et al. 2023) randomized patients
with T2/3 recctal cancer to neoadjuvant therapy with chemoradiation vs
FOLFOX. 585 in the FOLFOX group and 543 in the chemoradiotherapy group.
FOLFOX was noninferior to chemoradiotherapy for disease-free survival
and the groups were similar with respect to overall survival. In the
FOLFOX group, 53 patients (9.1\%) received preoperative
chemoradiotherapy (if the primary tumor decreased in size by
\textless20\% or if FOLFOX was discontinued because of side effects) and
8 (1.4\%) received postoperative chemoradiotherapy.

\part{Anal Cancer}

\chapter{Anal Squamous Cell
Carcinoma}\label{anal-squamous-cell-carcinoma}

\href{https://www.nccn.org/professionals/physician_gls/pdf/anal_blocks.pdf}{NCCN
Guidelines}

Surgical Clinics Review Article: (Young et al. 2020)

\section{Chemoradiation}\label{chemoradiation}

Chemoradiation is now the standard for anal squamous cell carcinoma of
the anal canal and for perianal cancers (except for small T1 lesions).
Nigro protocol is the standard approach.(Nigro et al. 1983)

\subsection{Restaging after chemoRT}\label{restaging-after-chemort}

Based on the results of the ACT-II study, it may be appropriate to
follow patients who have not achieved a complete clinical response with
persistent anal cancer up to 6 months following completion of radiation
therapy and chemotherapy as long as there is no evidence of progressive
disease during this period of follow-up. Persistent disease may continue
to regress even at 26 weeks from the start of treatment.(James et al.
2013)

\section{Anal Intraepithelial
Neoplasia}\label{anal-intraepithelial-neoplasia}

Review of diagnosis and treatment (Siddharthan, Lanciault, and Tsikitis
2019)

HPV vaccination has been shown to reduce the incidence of AIN in at-risk
populations, particularly those \textless26 years of age (Wei et al.
2023)

\part{Sarcoma}

\chapter{Soft Tissue Neoplasms}\label{soft-tissue-neoplasms}

\section{Desmoid Tumors}\label{desmoid-tumors}

DeFi Phase 3 trial of Nirogacestat (Gounder et al. 2023)

Milan Consensus Guidelines 2023 (Kasper et al. 2024)

\section{Lipomatous Tumors}\label{lipomatous-tumors}

Use of mdm2 IHC/FISH for diagnosis of liposarcoma (Weaver et al. 2010)

\section{Retroperitoneal}\label{retroperitoneal}

\subsection{Preop Radiation}\label{preop-radiation}

STRASS trial randomized 266 patients to preop radiation followed by
surgical resection vs surgery alone. At a median followup of 43 months,
there was no difference in recurrence-free survival. Serious adverse
affects were more common in the preop radiation group (24\% vs 10\%).
One patients in the radiation group died of treatment-related toxicity
(gastrocolic fistula), compared with none in the surgery alone group
(Bonvalot et al. 2020). See commentary: (Cardona 2020)

STRASS2: Ongoing trial of neoadjuvant chemotherapy. (Lambdin et al.
2023)

\section{Peritoneal mesothelioma}\label{peritoneal-mesothelioma}

Review (Bridda et al. 2007)

Surgical Oncology Clinics review (Li and Alexander 2018)

\bookmarksetup{startatroot}

\chapter*{References}\label{references-1}
\addcontentsline{toc}{chapter}{References}

\markboth{References}{References}

\phantomsection\label{refs}
\begin{CSLReferences}{1}{0}
\bibitem[\citeproctext]{ref-ajani3953}
Ajani, Jaffer A., Kathryn Winter, Gordon S. Okawara, John H. Donohue,
Peter W. T. Pisters, Christopher H. Crane, John F. Greskovich, et al.
2006. {``Phase {II} Trial of Preoperative Chemoradiation in Patients
with Localized Gastric Adenocarcinoma ({RTOG} 9904): Quality of Combined
Modality Therapy and Pathologic Response.''} \emph{Journal of Clinical
Oncology: Official Journal of the American Society of Clinical Oncology}
24 (24): 3953--58. \url{https://doi.org/10.1200/JCO.2006.06.4840}.

\bibitem[\citeproctext]{ref-al-batran1948}
Al-Batran, Salah-Eddin, Nils Homann, Claudia Pauligk, Thorsten O.
Goetze, Johannes Meiler, Stefan Kasper, Hans-Georg Kopp, et al. 2019.
{``Perioperative Chemotherapy with Fluorouracil Plus Leucovorin,
Oxaliplatin, and Docetaxel Versus Fluorouracil or Capecitabine Plus
Cisplatin and Epirubicin for Locally Advanced, Resectable Gastric or
Gastro-Oesophageal Junction Adenocarcinoma ({FLOT4}): A Randomised,
Phase 2/3 Trial.''} \emph{Lancet (London, England)} 393 (10184):
1948--57. \url{https://doi.org/10.1016/S0140-6736(18)32557-1}.

\bibitem[\citeproctext]{ref-al-batran1237}
Al-Batran, Salah-Eddin, Nils Homann, Claudia Pauligk, Gerald Illerhaus,
Uwe M. Martens, Jan Stoehlmacher, Harald Schmalenberg, et al. 2017.
{``Effect of {Neoadjuvant} {Chemotherapy} {Followed} by {Surgical}
{Resection} on {Survival} in {Patients} {With} {Limited} {Metastatic}
{Gastric} or {Gastroesophageal} {Junction} {Cancer}: {The} {AIO}-{FLOT3}
{Trial}.''} \emph{JAMA Oncology} 3 (9): 1237--44.
\url{https://doi.org/10.1001/jamaoncol.2017.0515}.

\bibitem[\citeproctext]{ref-allum5062}
Allum, William H., Sally P. Stenning, John Bancewicz, Peter I. Clark,
and Ruth E. Langley. 2009. {``Long-Term Results of a Randomized Trial of
Surgery with or Without Preoperative Chemotherapy in Esophageal
Cancer.''} \emph{Journal of Clinical Oncology: Official Journal of the
American Society of Clinical Oncology} 27 (30): 5062--67.
\url{https://doi.org/10.1200/JCO.2009.22.2083}.

\bibitem[\citeproctext]{ref-bahadoer29}
Bahadoer, Renu R., Esmée A. Dijkstra, Boudewijn van Etten, Corrie A. M.
Marijnen, Hein Putter, Elma Meershoek-Klein Kranenbarg, Annet G. H.
Roodvoets, et al. 2021. {``Short-Course Radiotherapy Followed by
Chemotherapy Before Total Mesorectal Excision ({TME}) Versus
Preoperative Chemoradiotherapy, {TME}, and Optional Adjuvant
Chemotherapy in Locally Advanced Rectal Cancer ({RAPIDO}): A Randomised,
Open-Label, Phase 3 Trial.''} \emph{The Lancet. Oncology} 22 (1):
29--42. \url{https://doi.org/10.1016/S1470-2045(20)30555-6}.

\bibitem[\citeproctext]{ref-baker987}
Baker, M, V Halliday, RN Williams, and DJ Bowery. 2016. {``A Systematic
Review of the Nutritional Consequences of Esophagectomy.''}
\emph{Clinical Nutrition (Edinburgh, Scotland)} 35 (5): 987--97.
\url{https://doi.org/10.1016/j.clnu.2015.08.010}.

\bibitem[\citeproctext]{ref-bang315}
Bang, Yung-Jue, Young-Woo Kim, Han-Kwang Yang, Hyun Cheol Chung,
Young-Kyu Park, Kyung Hee Lee, Keun-Wook Lee, et al. 2012. {``Adjuvant
Capecitabine and Oxaliplatin for Gastric Cancer After {D2} Gastrectomy
({CLASSIC}): A Phase 3 Open-Label, Randomised Controlled Trial.''}
\emph{Lancet (London, England)} 379 (9813): 315--21.
\url{https://doi.org/10.1016/S0140-6736(11)61873-4}.

\bibitem[\citeproctext]{ref-bedenne1160}
Bedenne, Laurent, Pierre Michel, Olivier Bouché, Chantal Milan,
Christophe Mariette, Thierry Conroy, Denis Pezet, et al. 2007.
{``Chemoradiation Followed by Surgery Compared with Chemoradiation Alone
in Squamous Cancer of the Esophagus: {FFCD} 9102.''} \emph{Journal of
Clinical Oncology: Official Journal of the American Society of Clinical
Oncology} 25 (10): 1160--68.
\url{https://doi.org/10.1200/JCO.2005.04.7118}.

\bibitem[\citeproctext]{ref-bergeron765}
Bergeron, Edward J., Jules Lin, Andrew C. Chang, Mark B. Orringer, and
Rishindra M. Reddy. 2014. {``Endoscopic Ultrasound Is Inadequate to
Determine Which {T1}/{T2} Esophageal Tumors Are Candidates for
Endoluminal Therapies.''} \emph{The Journal of Thoracic and
Cardiovascular Surgery} 147 (2): 765-771: Discussion 771-773.
\url{https://doi.org/10.1016/j.jtcvs.2013.10.003}.

\bibitem[\citeproctext]{ref-beukema85}
Beukema, Jannet C., Peter van Luijk, Joachim Widder, Johannes A.
Langendijk, and Christina T. Muijs. 2015. {``Is Cardiac Toxicity a
Relevant Issue in the Radiation Treatment of Esophageal Cancer?''}
\emph{Radiotherapy and Oncology: Journal of the European Society for
Therapeutic Radiology and Oncology} 114 (1): 85--90.
\url{https://doi.org/10.1016/j.radonc.2014.11.037}.

\bibitem[\citeproctext]{ref-biere1887}
Biere, S. S., M. I. van Berge Henegouwen, K. W. Maas, L. Bonavina, C.
Rosman, J. R. Garcia, S. S. Gisbertz, et al. 2012. {``Minimally Invasive
Versus Open Oesophagectomy for Patients with Oesophageal Cancer: A
Multicentre, Open-Label, Randomised Controlled Trial.''} \emph{Lancet}
379 (9829): 1887--92.
\url{https://doi.org/10.1016/S0140-6736(12)60516-9}.

\bibitem[\citeproctext]{ref-birkmeyer2117}
Birkmeyer, John D., Therese A. Stukel, Andrea E. Siewers, Philip P.
Goodney, David E. Wennberg, and F. Lee Lucas. 2003. {``Surgeon Volume
and Operative Mortality in the {United} {States}.''} \emph{The New
England Journal of Medicine} 349 (22): 2117--27.
\url{https://doi.org/10.1056/NEJMsa035205}.

\bibitem[\citeproctext]{ref-blaire386}
Blair, Vanessa R., Maybelle McLeod, Fátima Carneiro, Daniel G. Coit,
Johanna L. D'Addario, Jolanda M. van Dieren, Kirsty L. Harris, et al.
2020. {``Hereditary Diffuse Gastric Cancer: Updated Clinical Practice
Guidelines.''} \emph{The Lancet. Oncology} 21 (8): e386--97.
\url{https://doi.org/10.1016/S1470-2045(20)30219-9}.

\bibitem[\citeproctext]{ref-bonvalot1366}
Bonvalot, Sylvie, Alessandro Gronchi, Cécile Le Péchoux, Carol J.
Swallow, Dirk Strauss, Pierre Meeus, Frits van Coevorden, et al. 2020.
{``Preoperative Radiotherapy Plus Surgery Versus Surgery Alone for
Patients with Primary Retroperitoneal Sarcoma ({EORTC}-62092: {STRASS}):
A Multicentre, Open-Label, Randomised, Phase 3 Trial.''} \emph{The
Lancet. Oncology} 21 (10): 1366--77.
\url{https://doi.org/10.1016/S1470-2045(20)30446-0}.

\bibitem[\citeproctext]{ref-bozzetti170}
Bozzetti, F., E. Marubini, G. Bonfanti, R. Miceli, C. Piano, and L.
Gennari. 1999. {``Subtotal Versus Total Gastrectomy for Gastric Cancer:
Five-Year Survival Rates in a Multicenter Randomized {Italian} Trial.
{Italian} {Gastrointestinal} {Tumor} {Study} {Group}.''} \emph{Annals of
Surgery} 230 (2): 170--78.
\url{https://doi.org/10.1097/00000658-199908000-00006}.

\bibitem[\citeproctext]{ref-bridda32}
Bridda, Alessio, Ilaria Padoan, Roberto Mencarelli, and Mauro Frego.
2007. {``Peritoneal {Mesothelioma}: {A} {Review}.''} \emph{Medscape
General Medicine} 9 (2): 32.
\url{https://www.ncbi.nlm.nih.gov/pmc/articles/PMC1994863/}.

\bibitem[\citeproctext]{ref-cao751}
Cao, Y., C. Liao, A. Tan, Y. Gao, Z. Mo, and F. Gao. 2009.
{``Meta-Analysis of Endoscopic Submucosal Dissection Versus Endoscopic
Mucosal Resection for Tumors of the Gastrointestinal Tract.''}
\emph{Endoscopy} 41 (9): 751--57.
\url{https://doi.org/10.1055/s-0029-1215053}.

\bibitem[\citeproctext]{ref-cardona1257}
Cardona, Kenneth. 2020. {``The {STRASS} Trial: An Important Step in the
Right Direction.''} \emph{The Lancet. Oncology} 21 (10): 1257--58.
\url{https://doi.org/10.1016/S1470-2045(20)30429-0}.

\bibitem[\citeproctext]{ref-castoro1375}
Castoro, Carlo, Marco Scarpa, Matteo Cagol, Rita Alfieri, Alberto Ruol,
Francesco Cavallin, Silvia Michieletto, et al. 2013. {``Complete
Clinical Response After Neoadjuvant Chemoradiotherapy for Squamous Cell
Cancer of the Thoracic Oesophagus: Is Surgery Always Necessary?''}
\emph{Journal of Gastrointestinal Surgery: Official Journal of the
Society for Surgery of the Alimentary Tract} 17 (8): 1375--81.
\url{https://doi.org/10.1007/s11605-013-2269-3}.

\bibitem[\citeproctext]{ref-cats616}
Cats, Annemieke, Edwin P. M. Jansen, Nicole C. T. van Grieken, Karolina
Sikorska, Pehr Lind, Marianne Nordsmark, Elma Meershoek-Klein
Kranenbarg, et al. 2018. {``Chemotherapy Versus Chemoradiotherapy After
Surgery and Preoperative Chemotherapy for Resectable Gastric Cancer
({CRITICS}): An International, Open-Label, Randomised Phase 3 Trial.''}
\emph{The Lancet. Oncology} 19 (5): 616--28.
\url{https://doi.org/10.1016/S1470-2045(18)30132-3}.

\bibitem[\citeproctext]{ref-cercek2297}
Cercek, Andrea, Michael B. Foote, Benoit Rousseau, J. Joshua Smith,
Jinru Shia, Jenna Sinopoli, Jill Weiss, et al. 2025. {``Nonoperative
{Management} of {Mismatch} {Repair}-{Deficient} {Tumors}.''} \emph{The
New England Journal of Medicine} 392 (23): 2297--2308.
\url{https://doi.org/10.1056/NEJMoa2404512}.

\bibitem[\citeproctext]{ref-cercek2363}
Cercek, Andrea, Melissa Lumish, Jenna Sinopoli, Jill Weiss, Jinru Shia,
Michelle Lamendola-Essel, Imane H. El Dika, et al. 2022. {``{PD}-1
{Blockade} in {Mismatch} {Repair}-{Deficient}, {Locally} {Advanced}
{Rectal} {Cancer}.''} \emph{The New England Journal of Medicine} 386
(25): 2363--76. \url{https://doi.org/10.1056/NEJMoa2201445}.

\bibitem[\citeproctext]{ref-chalabi1949}
Chalabi, Myriam, Yara L. Verschoor, Pedro Batista Tan, Sara Balduzzi,
Anja U. Van Lent, Cecile Grootscholten, Simone Dokter, et al. 2024.
{``Neoadjuvant {Immunotherapy} in {Locally} {Advanced} {Mismatch}
{Repair}-{Deficient} {Colon} {Cancer}.''} \emph{The New England Journal
of Medicine} 390 (21): 1949--58.
\url{https://doi.org/10.1056/NEJMoa2400634}.

\bibitem[\citeproctext]{ref-chatterton354}
Chatterton, B. E., I. Ho Shon, A. Baldey, N. Lenzo, A. Patrikeos, B.
Kelley, D. Wong, J. E. Ramshaw, and A. M. Scott. 2009. {``Positron
Emission Tomography Changes Management and Prognostic Stratification in
Patients with Oesophageal Cancer: Results of a Multicentre Prospective
Study.''} \emph{European Journal of Nuclear Medicine and Molecular
Imaging} 36 (3): 354--61.
\url{https://doi.org/10.1007/s00259-008-0959-y}.

\bibitem[\citeproctext]{ref-chen2594}
Chen, Yijun, Kerry Kingham, James M. Ford, James Rosing, Jacques Van
Dam, R. Brooke Jeffrey, Teri A. Longacre, Nicki Chun, Allison Kurian,
and Jeffrey A. Norton. 2011. {``A Prospective Study of Total Gastrectomy
for {CDH1}-Positive Hereditary Diffuse Gastric Cancer.''} \emph{Annals
of Surgical Oncology} 18 (9): 2594--98.
\url{https://doi.org/10.1245/s10434-011-1648-9}.

\bibitem[\citeproctext]{ref-chirieac1347}
Chirieac, Lucian R., Stephen G. Swisher, Jaffer A. Ajani, Ritsuko R.
Komaki, Arlene M. Correa, Jeffrey S. Morris, Jack A. Roth, Asif Rashid,
Stanley R. Hamilton, and Tsung-Teh Wu. 2005. {``Posttherapy Pathologic
Stage Predicts Survival in Patients with Esophageal Carcinoma Receiving
Preoperative Chemoradiation.''} \emph{Cancer} 103 (7): 1347--55.
\url{https://doi.org/10.1002/cncr.20916}.

\bibitem[\citeproctext]{ref-conroy702}
Conroy, Thierry, Jean-François Bosset, Pierre-Luc Etienne, Emmanuel Rio,
Éric François, Nathalie Mesgouez-Nebout, Véronique Vendrely, et al.
2021. {``Neoadjuvant Chemotherapy with {FOLFIRINOX} and Preoperative
Chemoradiotherapy for Patients with Locally Advanced Rectal Cancer
({UNICANCER}-{PRODIGE} 23): A Multicentre, Randomised, Open-Label, Phase
3 Trial.''} \emph{The Lancet. Oncology} 22 (5): 702--15.
\url{https://doi.org/10.1016/S1470-2045(21)00079-6}.

\bibitem[\citeproctext]{ref-cunningham11}
Cunningham, David, William H. Allum, Sally P. Stenning, Jeremy N.
Thompson, Cornelis J. H. Van de Velde, Marianne Nicolson, J. Howard
Scarffe, et al. 2006. {``Perioperative Chemotherapy Versus Surgery Alone
for Resectable Gastroesophageal Cancer.''} \emph{The New England Journal
of Medicine} 355 (1): 11--20.
\url{https://doi.org/10.1056/NEJMoa055531}.

\bibitem[\citeproctext]{ref-dorth2099}
Dorth, Jennifer A., John A. Pura, Manisha Palta, Christopher G. Willett,
Hope E. Uronis, Thomas A. D'Amico, and Brian G. Czito. 2014. {``Patterns
of Recurrence After Trimodality Therapy for Esophageal Cancer.''}
\emph{Cancer} 120 (14): 2099--2105.
\url{https://doi.org/10.1002/cncr.28703}.

\bibitem[\citeproctext]{ref-dunbar850}
Dunbar, Kerry B., and Stuart Jon Spechler. 2012. {``The Risk of
Lymph-Node Metastases in Patients with High-Grade Dysplasia or
Intramucosal Carcinoma in {Barrett}'s Esophagus: A Systematic Review.''}
\emph{The American Journal of Gastroenterology} 107 (6): 850--862; quiz
863. \url{https://doi.org/10.1038/ajg.2012.78}.

\bibitem[\citeproctext]{ref-edwardsS1550}
Edwards, Michael A., Kinga Powers, R. Wesley Vosburg, Randal Zhou,
Andrea Stroud, Nabeel R. Obeid, John Pilcher, et al. 2025. {``American
{Society} for {Metabolic} and {Bariatric} {Surgery}: Postoperative Care
Pathway Guidelines for {Roux}-En-{Y} Gastric Bypass.''} \emph{Surgery
for Obesity and Related Diseases: Official Journal of the American
Society for Bariatric Surgery}, January, S1550-7289(25)00059-0.
\url{https://doi.org/10.1016/j.soard.2025.01.005}.

\bibitem[\citeproctext]{ref-eyck1995}
Eyck, Ben M., J. Jan B. van Lanschot, Maarten C. C. M. Hulshof, Berend
J. van der Wilk, Joel Shapiro, Pieter van Hagen, Mark I. van Berge
Henegouwen, et al. 2021. {``Ten-{Year} {Outcome} of {Neoadjuvant}
{Chemoradiotherapy} {Plus} {Surgery} for {Esophageal} {Cancer}: {The}
{Randomized} {Controlled} {CROSS} {Trial}.''} \emph{Journal of Clinical
Oncology: Official Journal of the American Society of Clinical Oncology}
39 (18): 1995--2004. \url{https://doi.org/10.1200/JCO.20.03614}.

\bibitem[\citeproctext]{ref-frandsen516}
Frandsen, Jonathan, Dustin Boothe, David K. Gaffney, Brent D. Wilson,
and Shane Lloyd. 2015. {``Increased Risk of Death Due to Heart Disease
After Radiotherapy for Esophageal Cancer.''} \emph{Journal of
Gastrointestinal Oncology} 6 (5): 516--23.
\url{https://doi.org/10.3978/j.issn.2078-6891.2015.040}.

\bibitem[\citeproctext]{ref-fujitani309}
Fujitani, Kazumasa, Han-Kwang Yang, Junki Mizusawa, Young-Woo Kim,
Masanori Terashima, Sang-Uk Han, Yoshiaki Iwasaki, et al. 2016.
{``Gastrectomy Plus Chemotherapy Versus Chemotherapy Alone for Advanced
Gastric Cancer with a Single Non-Curable Factor ({REGATTA}): A Phase 3,
Randomised Controlled Trial.''} \emph{The Lancet. Oncology} 17 (3):
309--18. \url{https://doi.org/10.1016/S1470-2045(15)00553-7}.

\bibitem[\citeproctext]{ref-furlong107}
Furlong, Heidi, Gary Bass, Oscar Breathnach, Brian O'Neill, Eamonn Leen,
and Thomas N. Walsh. 2013. {``Targeting Therapy for Esophageal Cancer in
Patients Aged 70 and Over.''} \emph{Journal of Geriatric Oncology} 4
(2): 107--13. \url{https://doi.org/10.1016/j.jgo.2012.12.006}.

\bibitem[\citeproctext]{ref-ganz35}
Ganz, Robert A., Bergein F. Overholt, Virender K. Sharma, David E.
Fleischer, Nicholas J. Shaheen, Charles J. Lightdale, Stephen R.
Freeman, et al. 2008. {``Circumferential Ablation of {Barrett}'s
Esophagus That Contains High-Grade Dysplasia: A {U}.{S}. {Multicenter}
{Registry}.''} \emph{Gastrointestinal Endoscopy} 68 (1): 35--40.
\url{https://doi.org/10.1016/j.gie.2007.12.015}.

\bibitem[\citeproctext]{ref-garcia-aguilar2546}
Garcia-Aguilar, Julio, Sujata Patil, Marc J. Gollub, Jin K. Kim,
Jonathan B. Yuval, Hannah M. Thompson, Floris S. Verheij, et al. 2022.
{``Organ {Preservation} in {Patients} {With} {Rectal} {Adenocarcinoma}
{Treated} {With} {Total} {Neoadjuvant} {Therapy}.''} \emph{Journal of
Clinical Oncology: Official Journal of the American Society of Clinical
Oncology} 40 (23): 2546--56. \url{https://doi.org/10.1200/JCO.22.00032}.

\bibitem[\citeproctext]{ref-gastrointestinaltumorstudygroup1465}
Gastrointestinal Tumor Study Group. 1985. {``Prolongation of the
Disease-Free Interval in Surgically Treated Rectal Carcinoma.''}
\emph{The New England Journal of Medicine} 312 (23): 1465--72.
\url{https://doi.org/10.1056/NEJM198506063122301}.

\bibitem[\citeproctext]{ref-gharzaie0158916}
Gharzai, Laila, Vivek Verma, Kyle A. Denniston, Abhijeet R. Bhirud,
Nathan R. Bennion, and Chi Lin. 2016. {``Radiation {Therapy} and
{Cardiac} {Death} in {Long}-{Term} {Survivors} of {Esophageal} {Cancer}:
{An} {Analysis} of the {Surveillance}, {Epidemiology}, and {End}
{Result} {Database}.''} \emph{PloS One} 11 (7): e0158916.
\url{https://doi.org/10.1371/journal.pone.0158916}.

\bibitem[\citeproctext]{ref-vangijn575}
Gijn, Willem van, Corrie A. M. Marijnen, Iris D. Nagtegaal, Elma
Meershoek-Klein Kranenbarg, Hein Putter, Theo Wiggers, Harm J. T.
Rutten, et al. 2011. {``Preoperative Radiotherapy Combined with Total
Mesorectal Excision for Resectable Rectal Cancer: 12-Year Follow-up of
the Multicentre, Randomised Controlled {TME} Trial.''} \emph{The Lancet.
Oncology} 12 (6): 575--82.
\url{https://doi.org/10.1016/S1470-2045(11)70097-3}.

\bibitem[\citeproctext]{ref-gounder898}
Gounder, Mrinal, Ravin Ratan, Thierry Alcindor, Patrick Schöffski,
Winette T. van der Graaf, Breelyn A. Wilky, Richard F. Riedel, et al.
2023. {``Nirogacestat, a γ-{Secretase} {Inhibitor} for {Desmoid}
{Tumors}.''} \emph{The New England Journal of Medicine} 388 (10):
898--912. \url{https://doi.org/10.1056/NEJMoa2210140}.

\bibitem[\citeproctext]{ref-degraaf988}
Graaf, G. W. de, A. A. Ayantunde, S. L. Parsons, J. P. Duffy, and N. T.
Welch. 2007. {``The Role of Staging Laparoscopy in Oesophagogastric
Cancers.''} \emph{European Journal of Surgical Oncology: The Journal of
the European Society of Surgical Oncology and the British Association of
Surgical Oncology} 33 (8): 988--92.
\url{https://doi.org/10.1016/j.ejso.2007.01.007}.

\bibitem[\citeproctext]{ref-guttmann1131}
Guttmann, David M., Nandita Mitra, Justin Bekelman, James M. Metz, John
Plastaras, Weiwei Feng, and Samuel Swisher-McClure. 2017. {``Improved
{Overall} {Survival} with {Aggressive} {Primary} {Tumor} {Radiotherapy}
for {Patients} with {Metastatic} {Esophageal} {Cancer}.''} \emph{Journal
of Thoracic Oncology: Official Publication of the International
Association for the Study of Lung Cancer} 12 (7): 1131--42.
\url{https://doi.org/10.1016/j.jtho.2017.03.026}.

\bibitem[\citeproctext]{ref-vanhagen2074}
Hagen, P. van, M. C. C. M. Hulshof, J. J. B. van Lanschot, E. W.
Steyerberg, M. I. van Berge Henegouwen, B. P. L. Wijnhoven, D. J.
Richel, et al. 2012. {``Preoperative Chemoradiotherapy for Esophageal or
Junctional Cancer.''} \emph{The New England Journal of Medicine} 366
(22): 2074--84. \url{https://doi.org/10.1056/NEJMoa1112088}.

\bibitem[\citeproctext]{ref-haidry87}
Haidry, Rehan J., Jason M. Dunn, Mohammed A. Butt, Matthew G. Burnell,
Abhinav Gupta, Sarah Green, Haroon Miah, et al. 2013. {``Radiofrequency
Ablation and Endoscopic Mucosal Resection for Dysplastic Barrett's
Esophagus and Early Esophageal Adenocarcinoma: Outcomes of the {UK}
{National} {Halo} {RFA} {Registry}.''} \emph{Gastroenterology} 145 (1):
87--95. \url{https://doi.org/10.1053/j.gastro.2013.03.045}.

\bibitem[\citeproctext]{ref-hansford23}
Hansford, Samantha, Pardeep Kaurah, Hector Li-Chang, Michelle Woo,
Janine Senz, Hugo Pinheiro, Kasmintan A. Schrader, et al. 2015.
{``Hereditary {Diffuse} {Gastric} {Cancer} {Syndrome}: {CDH1}
{Mutations} and {Beyond}.''} \emph{JAMA Oncology} 1 (1): 23--32.
\url{https://doi.org/10.1001/jamaoncol.2014.168}.

\bibitem[\citeproctext]{ref-hardacker3739}
Hardacker, Thomas J., DuyKhanh Ceppa, Ikenna Okereke, Karen M. Rieger,
Shadia I. Jalal, Julia K. LeBlanc, John M. DeWitt, Kenneth A. Kesler,
and Thomas J. Birdas. 2014. {``Treatment of Clinical {T2N0M0} Esophageal
Cancer.''} \emph{Annals of Surgical Oncology} 21 (12): 3739--43.
\url{https://doi.org/10.1245/s10434-014-3929-6}.

\bibitem[\citeproctext]{ref-heald1479}
Heald, R. J., and R. D. Ryall. 1986. {``Recurrence and Survival After
Total Mesorectal Excision for Rectal Cancer.''} \emph{Lancet (London,
England)} 1 (8496): 1479--82.
\url{https://doi.org/10.1016/s0140-6736(86)91510-2}.

\bibitem[\citeproctext]{ref-hollis1114}
Hollis, Alexander C., Lauren M. Quinn, James Hodson, Emily Evans, James
Plowright, Ruksana Begum, Harriet Mitchell, Mike T. Hallissey, John L.
Whiting, and Ewen A. Griffiths. 2017. {``Prognostic Significance of
Tumor Length in Patients Receiving Esophagectomy for Esophageal
Cancer.''} \emph{Journal of Surgical Oncology} 116 (8): 1114--22.
\url{https://doi.org/10.1002/jso.24789}.

\bibitem[\citeproctext]{ref-holmebakk419}
Hølmebakk, T., B. Bjerkehagen, I. Hompland, S. Stoldt, and K. Boye.
2019. {``Relationship Between {R1} Resection, Tumour Rupture and
Recurrence in Resected Gastrointestinal Stromal Tumour.''} \emph{The
British Journal of Surgery} 106 (4): 419--26.
\url{https://doi.org/10.1002/bjs.11027}.

\bibitem[\citeproctext]{ref-hu38}
Hu, Huabin, Liang Kang, Jianwei Zhang, Zehua Wu, Hui Wang, Meijin Huang,
Ping Lan, et al. 2022. {``Neoadjuvant {PD}-1 Blockade with Toripalimab,
with or Without Celecoxib, in Mismatch Repair-Deficient or
Microsatellite Instability-High, Locally Advanced, Colorectal Cancer
({PICC}): A Single-Centre, Parallel-Group, Non-Comparative, Randomised,
Phase 2 Trial.''} \emph{The Lancet. Gastroenterology \& Hepatology} 7
(1): 38--48. \url{https://doi.org/10.1016/S2468-1253(21)00348-4}.

\bibitem[\citeproctext]{ref-hulscher1662}
Hulscher, Jan B. F., Johanna W. van Sandick, Angela G. E. M. de Boer,
Bas P. L. Wijnhoven, Jan G. P. Tijssen, Paul Fockens, Peep F. M.
Stalmeier, et al. 2002. {``Extended Transthoracic Resection Compared
with Limited Transhiatal Resection for Adenocarcinoma of the
Esophagus.''} \emph{The New England Journal of Medicine} 347 (21):
1662--69. \url{https://doi.org/10.1056/NEJMoa022343}.

\bibitem[\citeproctext]{ref-irani5}
Irani, Jennifer L., Traci L. Hedrick, Timothy E. Miller, Lawrence Lee,
Emily Steinhagen, Benjamin D. Shogan, Joel E. Goldberg, Daniel L.
Feingold, Amy L. Lightner, and Ian M. Paquette. 2023. {``Clinical
Practice Guidelines for Enhanced Recovery After Colon and Rectal Surgery
from the {American} {Society} of {Colon} and {Rectal} {Surgeons} and the
{Society} of {American} {Gastrointestinal} and {Endoscopic}
{Surgeons}.''} \emph{Surgical Endoscopy} 37 (1): 5--30.
\url{https://doi.org/10.1007/s00464-022-09758-x}.

\bibitem[\citeproctext]{ref-james516}
James, Roger D., Robert Glynne-Jones, Helen M. Meadows, David
Cunningham, Arthur Sun Myint, Mark P. Saunders, Timothy Maughan, et al.
2013. {``Mitomycin or Cisplatin Chemoradiation with or Without
Maintenance Chemotherapy for Treatment of Squamous-Cell Carcinoma of the
Anus ({ACT} {II}): A Randomised, Phase 3, Open-Label, 2 × 2 Factorial
Trial.''} \emph{The Lancet. Oncology} 14 (6): 516--24.
\url{https://doi.org/10.1016/S1470-2045(13)70086-X}.

\bibitem[\citeproctext]{ref-janjigian217}
Janjigian, Yelena Y., Salah-Eddin Al-Batran, Zev A. Wainberg, Kei Muro,
Daniela Molena, Eric Van Cutsem, Woo Jin Hyung, et al. 2025.
{``Perioperative {Durvalumab} in {Gastric} and {Gastroesophageal}
{Junction} {Cancer}.''} \emph{The New England Journal of Medicine} 393
(3): 217--30. \url{https://doi.org/10.1056/NEJMoa2503701}.

\bibitem[\citeproctext]{ref-jin1681}
Jin, Jing, Yuan Tang, Chen Hu, Li-Ming Jiang, Jun Jiang, Ning Li,
Wen-Yang Liu, et al. 2022. {``Multicenter, {Randomized}, {Phase} {III}
{Trial} of {Short}-{Term} {Radiotherapy} {Plus} {Chemotherapy} {Versus}
{Long}-{Term} {Chemoradiotherapy} in {Locally} {Advanced} {Rectal}
{Cancer} ({STELLAR}).''} \emph{Journal of Clinical Oncology: Official
Journal of the American Society of Clinical Oncology} 40 (15): 1681--92.
\url{https://doi.org/10.1200/JCO.21.01667}.

\bibitem[\citeproctext]{ref-kalffe749}
Kalff, Marianne C., Laura F. C. Fransen, Eline M. de Groot, Suzanne S.
Gisbertz, Grard A. P. Nieuwenhuijzen, Jelle P. Ruurda, Rob H. A.
Verhoeven, et al. 2022. {``Long-Term {Survival} {After} {Minimally}
{Invasive} {Versus} {Open} {Esophagectomy} for {Esophageal} {Cancer}:
{A} {Nationwide} {Propensity}-Score {Matched} {Analysis}.''}
\emph{Annals of Surgery} 276 (6): e749--57.
\url{https://doi.org/10.1097/SLA.0000000000004708}.

\bibitem[\citeproctext]{ref-kalff806}
Kalff, Marianne C., Sofie P. G. Henckens, Daan M. Voeten, David J.
Heineman, Maarten C. C. M. Hulshof, Hanneke W. M. van Laarhoven, Wietse
J. Eshuis, et al. 2022. {``Recurrent {Disease} {After} {Esophageal}
{Cancer} {Surgery}: {A} {Substudy} of {The} {Dutch} {Nationwide} {Ivory}
{Study}.''} \emph{Annals of Surgery} 276 (5): 806--13.
\url{https://doi.org/10.1097/SLA.0000000000005638}.

\bibitem[\citeproctext]{ref-kapiteijn638}
Kapiteijn, E., C. A. Marijnen, I. D. Nagtegaal, H. Putter, W. H. Steup,
T. Wiggers, H. J. Rutten, et al. 2001. {``Preoperative Radiotherapy
Combined with Total Mesorectal Excision for Resectable Rectal Cancer.''}
\emph{The New England Journal of Medicine} 345 (9): 638--46.
\url{https://doi.org/10.1056/NEJMoa010580}.

\bibitem[\citeproctext]{ref-kasper1121}
Kasper, Bernd, Elizabeth H. Baldini, Sylvie Bonvalot, Dario Callegaro,
Kenneth Cardona, Chiara Colombo, Nadège Corradini, et al. 2024.
{``Current {Management} of {Desmoid} {Tumors}: {A} {Review}.''}
\emph{JAMA Oncology} 10 (8): 1121--28.
\url{https://doi.org/10.1001/jamaoncol.2024.1805}.

\bibitem[\citeproctext]{ref-kato921}
Kato, Hiroyuki, Hiroyuki Kuwano, Masanobu Nakajima, Tatsuya Miyazaki,
Minako Yoshikawa, Hitoshi Ojima, Katsuhiko Tsukada, Noboru Oriuchi,
Tomio Inoue, and Keigo Endo. 2002.
{``\href{https://www.ncbi.nlm.nih.gov/pubmed/11920459}{Comparison
Between Positron Emission Tomography and Computed Tomography in the Use
of the Assessment of Esophageal Carcinoma}.''} \emph{Cancer} 94 (4):
921--28.

\bibitem[\citeproctext]{ref-kelly1191}
Kelly, Ronan J., Jaffer A. Ajani, Jaroslaw Kuzdzal, Thomas Zander, Eric
Van Cutsem, Guillaume Piessen, Guillermo Mendez, et al. 2021.
{``Adjuvant {Nivolumab} in {Resected} {Esophageal} or {Gastroesophageal}
{Junction} {Cancer}.''} \emph{The New England Journal of Medicine} 384
(13): 1191--1203. \url{https://doi.org/10.1056/NEJMoa2032125}.

\bibitem[\citeproctext]{ref-khan97}
Khan, Qamar J., and Carol J. Fabian. 2010. {``How {I} {Treat} {Vitamin}
{D} {Deficiency}.''} \emph{Journal of Oncology Practice} 6 (2): 97--101.
\url{https://doi.org/10.1200/JOP.091087}.

\bibitem[\citeproctext]{ref-kim403}
Kim, Tae Jung, Hyae Young Kim, Kyung Won Lee, and Moon Soo Kim. 2009.
{``Multimodality Assessment of Esophageal Cancer: Preoperative Staging
and Monitoring of Response to Therapy.''} \emph{Radiographics: A Review
Publication of the Radiological Society of North America, Inc} 29 (2):
403--21. \url{https://doi.org/10.1148/rg.292085106}.

\bibitem[\citeproctext]{ref-kofoed26}
Kofoed, S. C., A. Muhic, L. Baeksgaard, M. Jendresen, J. Gustafsen, J.
Holm, L. Bardram, B. Brandt, J. Brenø, and L. B. Svendsen. 2012.
{``Survival After Adjuvant Chemoradiotherapy or Surgery Alone in
Resectable Adenocarcinoma at the Gastro-Esophageal Junction.''}
\emph{Scandinavian Journal of Surgery: SJS: Official Organ for the
Finnish Surgical Society and the Scandinavian Surgical Society} 101 (1):
26--31. \url{https://doi.org/10.1177/145749691210100106}.

\bibitem[\citeproctext]{ref-kreis25}
Kreis, Martin E., R. Ruppert, H. Ptok, J. Strassburg, P. Brosi, A.
Lewin, M. R. Schön, et al. 2016. {``Use of {Preoperative} {Magnetic}
{Resonance} {Imaging} to {Select} {Patients} with {Rectal} {Cancer} for
{Neoadjuvant} {Chemoradiation}--{Interim} {Analysis} of the {German}
{OCUM} {Trial} ({NCT01325649}).''} \emph{Journal of Gastrointestinal
Surgery: Official Journal of the Society for Surgery of the Alimentary
Tract} 20 (1): 25-32; discussion 32-33.
\url{https://doi.org/10.1007/s11605-015-3011-0}.

\bibitem[\citeproctext]{ref-lambdin4573}
Lambdin, Jacob, Carrie Ryan, Stephanie Gregory, Kenneth Cardona,
Jonathan M. Hernandez, Winan J. van Houdt, and Alessandro Gronchi. 2023.
{``A {Randomized} {Phase} {III} {Study} of {Neoadjuvant} {Chemotherapy}
{Followed} by {Surgery} {Versus} {Surgery} {Alone} for {Patients} with
{High}-{Risk} {Retroperitoneal} {Sarcoma} ({STRASS2}).''} \emph{Annals
of Surgical Oncology} 30 (8): 4573--75.
\url{https://doi.org/10.1245/s10434-023-13500-9}.

\bibitem[\citeproctext]{ref-lee268}
Lee, Jeeyun, Do Hoon Lim, Sung Kim, Se Hoon Park, Joon Oh Park, Young
Suk Park, Ho Yeong Lim, et al. 2012. {``Phase {III} Trial Comparing
Capecitabine Plus Cisplatin Versus Capecitabine Plus Cisplatin with
Concurrent Capecitabine Radiotherapy in Completely Resected Gastric
Cancer with {D2} Lymph Node Dissection: The {ARTIST} Trial.''}
\emph{Journal of Clinical Oncology: Official Journal of the American
Society of Clinical Oncology} 30 (3): 268--73.
\url{https://doi.org/10.1200/JCO.2011.39.1953}.

\bibitem[\citeproctext]{ref-leichman4555}
Leichman, Lawrence P., Bryan H. Goldman, Pierre O. Bohanes, Heinz J.
Lenz, Charles R. Thomas, Kevin G. Billingsley, Christopher L. Corless,
et al. 2011. {``S0356: A Phase {II} Clinical and Prospective Molecular
Trial with Oxaliplatin, Fluorouracil, and External-Beam Radiation
Therapy Before Surgery for Patients with Esophageal Adenocarcinoma.''}
\emph{Journal of Clinical Oncology: Official Journal of the American
Society of Clinical Oncology} 29 (34): 4555--60.
\url{https://doi.org/10.1200/JCO.2011.36.7490}.

\bibitem[\citeproctext]{ref-leong2252}
Leong, Trevor, B. Mark Smithers, Karin Haustermans, Michael Michael, Val
Gebski, Danielle Miller, John Zalcberg, et al. 2017. {``{TOPGEAR}: {A}
{Randomized}, {Phase} {III} {Trial} of {Perioperative} {ECF}
{Chemotherapy} with or {Without} {Preoperative} {Chemoradiation} for
{Resectable} {Gastric} {Cancer}: {Interim} {Results} from an
{International}, {Intergroup} {Trial} of the {AGITG}, {TROG}, {EORTC}
and {CCTG}.''} \emph{Annals of Surgical Oncology} 24 (8): 2252--58.
\url{https://doi.org/10.1245/s10434-017-5830-6}.

\bibitem[\citeproctext]{ref-li539}
Li, Claire Yue, and H. Richard Alexander. 2018. {``Peritoneal
{Metastases} from {Malignant} {Mesothelioma}.''} \emph{Surgical Oncology
Clinics of North America} 27 (3): 539--49.
\url{https://doi.org/10.1016/j.soc.2018.02.010}.

\bibitem[\citeproctext]{ref-macdonald725}
Macdonald, J. S., S. R. Smalley, J. Benedetti, S. A. Hundahl, N. C.
Estes, G. N. Stemmermann, D. G. Haller, et al. 2001.
{``Chemoradiotherapy After Surgery Compared with Surgery Alone for
Adenocarcinoma of the Stomach or Gastroesophageal Junction.''} \emph{The
New England Journal of Medicine} 345 (10): 725--30.
\url{https://doi.org/10.1056/NEJMoa010187}.

\bibitem[\citeproctext]{ref-maish1777}
Maish, Mary S., and Steven R. DeMeester. 2004. {``Endoscopic Mucosal
Resection as a Staging Technique to Determine the Depth of Invasion of
Esophageal Adenocarcinoma.''} \emph{The Annals of Thoracic Surgery} 78
(5): 1777--82. \url{https://doi.org/10.1016/j.athoracsur.2004.04.064}.

\bibitem[\citeproctext]{ref-malik1}
Malik, Vinod, Ciaran Johnston, Dermot O'Toole, Julie Lucey, Naoimh
O'Farrell, Zieta Claxton, and John V. Reynolds. 2017. {``Metabolic Tumor
Volume Provides Complementary Prognostic Information to {EUS} Staging in
Esophageal and Junctional Cancer.''} \emph{Diseases of the Esophagus:
Official Journal of the International Society for Diseases of the
Esophagus} 30 (3): 1--8. \url{https://doi.org/10.1111/dote.12505}.

\bibitem[\citeproctext]{ref-mantziari7}
Mantziari, Styliani, Anastasia Pomoni, John O. Prior, Michael Winiker,
Pierre Allemann, Nicolas Demartines, and Markus Schäfer. 2020. {``{18F}-
{FDG} {PET}/{CT}-Derived Parameters Predict Clinical Stage and Prognosis
of Esophageal Cancer.''} \emph{BMC Medical Imaging} 20 (1): 7.
\url{https://doi.org/10.1186/s12880-019-0401-x}.

\bibitem[\citeproctext]{ref-margolis1694}
Margolis, Marc, Pendleton Alexander, Gregory D Trachiotis, Farid
Gharagozloo, and Timothy Lipman. 2003. {``Percutaneous Endoscopic
Gastrostomy Before Multimodality Therapy in Patients with Esophageal
Cancer.''} \emph{The Annals of Thoracic Surgery} 76 (5): 1694--98.
\url{https://doi.org/10.1016/S0003-4975(02)04890-7}.

\bibitem[\citeproctext]{ref-mariette2416}
Mariette, C., L. Dahan, F. Mornex, E. Maillard, P. A. Thomas, B.
Meunier, V. Boige, et al. 2014. {``Surgery Alone Versus
Chemoradiotherapy Followed by Surgery for Stage {I} and {II} Esophageal
Cancer: Final Analysis of Randomized Controlled Phase {III} Trial {FFCD}
9901.''} \emph{J Clin Oncol} 32 (23): 2416--22.
\url{https://doi.org/10.1200/JCO.2013.53.6532}.

\bibitem[\citeproctext]{ref-mariette152}
Mariette, Christophe, Sheraz R. Markar, Tienhan S. Dabakuyo-Yonli,
Bernard Meunier, Denis Pezet, Denis Collet, Xavier B. D'Journo, et al.
2019. {``Hybrid {Minimally} {Invasive} {Esophagectomy} for {Esophageal}
{Cancer}.''} \emph{New England Journal of Medicine} 380 (2): 152--62.
\url{https://doi.org/10.1056/NEJMoa1805101}.

\bibitem[\citeproctext]{ref-markar59}
Markar, Sheraz R., Caroline Gronnier, Arnaud Pasquer, Alain Duhamel,
Hélène Beal, Jérémie Théreaux, Johan Gagnière, et al. 2016. {``Role of
Neoadjuvant Treatment in Clinical {T2N0M0} Oesophageal Cancer: Results
from a Retrospective Multi-Center {European} Study.''} \emph{European
Journal of Cancer (Oxford, England: 1990)} 56 (March): 59--68.
\url{https://doi.org/10.1016/j.ejca.2015.11.024}.

\bibitem[\citeproctext]{ref-markar922}
Markar, Sheraz R., Alan Karthikesalingam, Marta Penna, and Donald E.
Low. 2014. {``Assessment of Short-Term Clinical Outcomes Following
Salvage Esophagectomy for the Treatment of Esophageal Malignancy:
Systematic Review and Pooled Analysis.''} \emph{Annals of Surgical
Oncology} 21 (3): 922--31.
\url{https://doi.org/10.1245/s10434-013-3364-0}.

\bibitem[\citeproctext]{ref-martin1308}
Martin, L., and P. Lagergren. 2009. {``Long-Term Weight Change After
Oesophageal Cancer Surgery.''} \emph{The British Journal of Surgery} 96
(11): 1308--14. \url{https://doi.org/10.1002/bjs.6723}.

\bibitem[\citeproctext]{ref-mckeown259}
McKeown, K. C. 1976. {``Total Three-Stage Oesophagectomy for Cancer of
the Oesophagus.''} \emph{The British Journal of Surgery} 63 (4):
259--62. \url{https://doi.org/10.1002/bjs.1800630403}.

\bibitem[\citeproctext]{ref-mellow1443}
Mellow, M. H., and H. Pinkas. 1985.
{``\href{https://www.ncbi.nlm.nih.gov/pubmed/4026476}{Endoscopic Laser
Therapy for Malignancies Affecting the Esophagus and Gastroesophageal
Junction. {Analysis} of Technical and Functional Efficacy}.''}
\emph{Archives of Internal Medicine} 145 (8): 1443--46.

\bibitem[\citeproctext]{ref-mercurystudygroup779}
MERCURY Study Group. 2006. {``Diagnostic Accuracy of Preoperative
Magnetic Resonance Imaging in Predicting Curative Resection of Rectal
Cancer: Prospective Observational Study.''} \emph{BMJ (Clinical Research
Ed.)} 333 (7572): 779.
\url{https://doi.org/10.1136/bmj.38937.646400.55}.

\bibitem[\citeproctext]{ref-messager318}
Messager, Mathieu, Xavier Mirabel, Emmanuelle Tresch, Amaury Paumier,
Véronique Vendrely, Laetitia Dahan, Olivier Glehen, et al. 2016.
{``Preoperative Chemoradiation with Paclitaxel-Carboplatin or with
Fluorouracil-Oxaliplatin-Folinic Acid ({FOLFOX}) for Resectable
Esophageal and Junctional Cancer: The {PROTECT}-1402, Randomized Phase 2
Trial.''} \emph{BMC Cancer} 16 (May): 318.
\url{https://doi.org/10.1186/s12885-016-2335-9}.

\bibitem[\citeproctext]{ref-minsky1167}
Minsky, Bruce D., Thomas F. Pajak, Robert J. Ginsberg, Thomas M.
Pisansky, James Martenson, Ritsuko Komaki, Gordon Okawara, Seth A.
Rosenthal, and David P. Kelsen. 2002. {``{INT} 0123 ({Radiation}
{Therapy} {Oncology} {Group} 94-05) Phase {III} Trial of
Combined-Modality Therapy for Esophageal Cancer: High-Dose Versus
Standard-Dose Radiation Therapy.''} \emph{Journal of Clinical Oncology:
Official Journal of the American Society of Clinical Oncology} 20 (5):
1167--74. \url{https://doi.org/10.1200/JCO.2002.20.5.1167}.

\bibitem[\citeproctext]{ref-nagtegaal303}
Nagtegaal, Iris D., and Phil Quirke. 2008. {``What Is the Role for the
Circumferential Margin in the Modern Treatment of Rectal Cancer?''}
\emph{Journal of Clinical Oncology: Official Journal of the American
Society of Clinical Oncology} 26 (2): 303--12.
\url{https://doi.org/10.1200/JCO.2007.12.7027}.

\bibitem[\citeproctext]{ref-nakamura1758}
Nakamura, Masaki, Toshiyasu Ojima, Mikihito Nakamori, Masahiro Katsuda,
Toshiaki Tsuji, Keiji Hayata, Tomoya Kato, and Hiroki Yamaue. 2019.
{``Conversion {Surgery} for {Gastric} {Cancer} with {Peritoneal}
{Metastasis} {Based} on the {Diagnosis} of {Second}-{Look} {Staging}
{Laparoscopy}.''} \emph{Journal of Gastrointestinal Surgery: Official
Journal of the Society for Surgery of the Alimentary Tract} 23 (9):
1758--66. \url{https://doi.org/10.1007/s11605-018-3983-7}.

\bibitem[\citeproctext]{ref-nigro1826}
Nigro, Nd, Hg Seydel, B Considine, Vk Vaitkevicius, L Leichman, and Jj
Kinzie. 1983. {``Combined Preoperative Radiation and Chemotherapy for
Squamous Cell Carcinoma of the Anal Canal.''} \emph{Cancer} 51 (10):
1826.
\url{https://doi.org/10.1002/1097-0142(19830515)51:10\%3C1826::aid-cncr2820511012\%3E3.0.co;2-l}.

\bibitem[\citeproctext]{ref-nilsson1015}
Nilsson, K., F. Klevebro, B. Sunde, I. Rouvelas, M. Lindblad, E. Szabo,
I. Halldestam, et al. 2023. {``Oncological Outcomes of Standard Versus
Prolonged Time to Surgery After Neoadjuvant Chemoradiotherapy for
Oesophageal Cancer in the Multicentre, Randomised, Controlled {NeoRes}
{II} Trial.''} \emph{Annals of Oncology: Official Journal of the
European Society for Medical Oncology} 34 (11): 1015--24.
\url{https://doi.org/10.1016/j.annonc.2023.08.010}.

\bibitem[\citeproctext]{ref-noh1389}
Noh, Sung Hoon, Sook Ryun Park, Han-Kwang Yang, Hyun Cheol Chung, Ik-Joo
Chung, Sang-Woon Kim, Hyung-Ho Kim, et al. 2014. {``Adjuvant
Capecitabine Plus Oxaliplatin for Gastric Cancer After {D2} Gastrectomy
({CLASSIC}): 5-Year Follow-up of an Open-Label, Randomised Phase 3
Trial.''} \emph{The Lancet. Oncology} 15 (12): 1389--96.
\url{https://doi.org/10.1016/S1470-2045(14)70473-5}.

\bibitem[\citeproctext]{ref-noordman965}
Noordman, Bo Jan, Manon C. W. Spaander, Roelf Valkema, Bas P. L.
Wijnhoven, Mark I. van Berge Henegouwen, Joël Shapiro, Katharina
Biermann, et al. 2018. {``Detection of Residual Disease After
Neoadjuvant Chemoradiotherapy for Oesophageal Cancer ({preSANO}): A
Prospective Multicentre, Diagnostic Cohort Study.''} \emph{The Lancet.
Oncology} 19 (7): 965--74.
\url{https://doi.org/10.1016/S1470-2045(18)30201-8}.

\bibitem[\citeproctext]{ref-nurkin1090}
Nurkin, Steven J., Hector R. Nava, Sai Yendamuri, Charles M. LeVea,
Chumy E. Nwogu, Adrienne Groman, Gregory Wilding, Andrew J. Bain, Steven
N. Hochwald, and Nikhil I. Khushalani. 2014. {``Outcomes of Endoscopic
Resection for High-Grade Dysplasia and Esophageal Cancer.''}
\emph{Surgical Endoscopy} 28 (4): 1090--95.
\url{https://doi.org/10.1007/s00464-013-3270-3}.

\bibitem[\citeproctext]{ref-ohnmacht311}
Ohnmacht, G. A., M. S. Allen, S. D. Cassivi, C. Deschamps, F. C.
Nichols, and P. C. Pairolero. 2006. {``Percutaneous Endoscopic
Gastrostomy Risks Rendering the Gastric Conduit Unusable for
Esophagectomy.''} \emph{Diseases of the Esophagus: Official Journal of
the International Society for Diseases of the Esophagus} 19 (4):
311--12. \url{https://doi.org/10.1111/j.1442-2050.2006.00588.x}.

\bibitem[\citeproctext]{ref-omloo992}
Omloo, Jikke M. T., Sjoerd M. Lagarde, Jan B. F. Hulscher, Johannes B.
Reitsma, Paul Fockens, Herman van Dekken, Fiebo J. W. Ten Kate, Huug
Obertop, Hugo W. Tilanus, and J. Jan B. van Lanschot. 2007. {``Extended
Transthoracic Resection Compared with Limited Transhiatal Resection for
Adenocarcinoma of the Mid/Distal Esophagus: Five-Year Survival of a
Randomized Clinical Trial.''} \emph{Annals of Surgery} 246 (6):
992-1000; discussion 1000-1001.
\url{https://doi.org/10.1097/SLA.0b013e31815c4037}.

\bibitem[\citeproctext]{ref-oppedijk385}
Oppedijk, Vera, Ate van der Gaast, Jan J. B. van Lanschot, Pieter van
Hagen, Rob van Os, Caroline M. van Rij, Maurice J. van der Sangen, et
al. 2014. {``Patterns of Recurrence After Surgery Alone Versus
Preoperative Chemoradiotherapy and Surgery in the {CROSS} Trials.''}
\emph{Journal of Clinical Oncology: Official Journal of the American
Society of Clinical Oncology} 32 (5): 385--91.
\url{https://doi.org/10.1200/JCO.2013.51.2186}.

\bibitem[\citeproctext]{ref-orringer643}
Orringer, M. B., and H. Sloan. 1978.
{``\href{https://www.ncbi.nlm.nih.gov/pubmed/703369}{Esophagectomy
Without Thoracotomy}.''} \emph{The Journal of Thoracic and
Cardiovascular Surgery} 76 (5): 643--54.

\bibitem[\citeproctext]{ref-orringer363}
Orringer, Mark B., Becky Marshall, Andrew C. Chang, Julia Lee, Allan
Pickens, and Christine L. Lau. 2007. {``Two Thousand Transhiatal
Esophagectomies: Changing Trends, Lessons Learned.''} \emph{Annals of
Surgery} 246 (3): 363-372; discussion 372-374.
\url{https://doi.org/10.1097/SLA.0b013e31814697f2}.

\bibitem[\citeproctext]{ref-ouattara1088}
Ouattara, Moussa, Xavier Benoit D'Journo, Anderson Loundou, Delphine
Trousse, Laetitia Dahan, Christophe Doddoli, Jean Francois Seitz, and
Pascal-Alexandre Thomas. 2012. {``Body Mass Index Kinetics and Risk
Factors of Malnutrition One Year After Radical Oesophagectomy for
Cancer.''} \emph{European Journal of Cardio-Thoracic Surgery: Official
Journal of the European Association for Cardio-Thoracic Surgery} 41 (5):
1088--93. \url{https://doi.org/10.1093/ejcts/ezr182}.

\bibitem[\citeproctext]{ref-overtoom}
Overtoom, Hidde C. G., Ben M. Eyck, Berend J. van der Wilk, Bo J.
Noordman, Pieter C. van der Sluis, Bas P. L. Wijnhoven, J. Jan B. van
Lanschot, Sjoerd M. Lagarde, and SANO study group*. 2024. {``Prolonged
{Time} to {Surgery} in {Patients} with {Residual} {Disease} {After}
{Neoadjuvant} {Chemoradiotherapy} for {Esophageal} {Cancer}.''}
\emph{Annals of Surgery}, August.
\url{https://doi.org/10.1097/SLA.0000000000006488}.

\bibitem[\citeproctext]{ref-paradis558}
Paradis, Tiffany, Stephan Robitaille, Anna Wang, Camille Gervais, A.
Sender Liberman, Patrick Charlebois, Barry L. Stein, Julio F. Fiore,
Liane S. Feldman, and Lawrence Lee. 2024. {``Predictive {Factors} for
{Successful} {Same}-{Day} {Discharge} {After} {Minimally} {Invasive}
{Colectomy} and {Stoma} {Reversal}.''} \emph{Diseases of the Colon and
Rectum} 67 (4): 558--65.
\url{https://doi.org/10.1097/DCR.0000000000003149}.

\bibitem[\citeproctext]{ref-paty365}
Paty, P. B., W. E. Enker, A. M. Cohen, and G. Y. Lauwers. 1994.
{``Treatment of Rectal Cancer by Low Anterior Resection with Coloanal
Anastomosis.''} \emph{Annals of Surgery} 219 (4): 365--73.
\url{https://doi.org/10.1097/00000658-199404000-00007}.

\bibitem[\citeproctext]{ref-pech1200}
Pech, O., A. Behrens, A. May, L. Nachbar, L. Gossner, T. Rabenstein, H.
Manner, et al. 2008. {``Long-Term Results and Risk Factor Analysis for
Recurrence After Curative Endoscopic Therapy in 349 Patients with
High-Grade Intraepithelial Neoplasia and Mucosal Adenocarcinoma in
{Barrett}'s Oesophagus.''} \emph{Gut} 57 (9): 1200--1206.
\url{https://doi.org/10.1136/gut.2007.142539}.

\bibitem[\citeproctext]{ref-pech652}
Pech, Oliver, Andrea May, Hendrik Manner, Angelika Behrens, Jürgen Pohl,
Maren Weferling, Urs Hartmann, et al. 2014. {``Long-Term Efficacy and
Safety of Endoscopic Resection for Patients with Mucosal Adenocarcinoma
of the Esophagus.''} \emph{Gastroenterology} 146 (3): 652--660.e1.
\url{https://doi.org/10.1053/j.gastro.2013.11.006}.

\bibitem[\citeproctext]{ref-penniment114}
Penniment, Michael G., Paolo B. De Ieso, Jennifer A. Harvey, Sonya
Stephens, Heather-Jane Au, Christopher J. O'Callaghan, Andrew Kneebone,
et al. 2018. {``Palliative Chemoradiotherapy Versus Radiotherapy Alone
for Dysphagia in Advanced Oesophageal Cancer: A Multicentre Randomised
Controlled Trial ({TROG} 03.01).''} \emph{The Lancet. Gastroenterology
\& Hepatology} 3 (2): 114--24.
\url{https://doi.org/10.1016/S2468-1253(17)30363-1}.

\bibitem[\citeproctext]{ref-peters858}
Peters, J. H., J. W. Kronson, M. Katz, and T. R. DeMeester. 1995.
{``Arterial Anatomic Considerations in Colon Interposition for
Esophageal Replacement.''} \emph{Archives of Surgery (Chicago, Ill.:
1960)} 130 (8): 858-862; discussion 862-863.
\url{https://doi.org/10.1001/archsurg.1995.01430080060009}.

\bibitem[\citeproctext]{ref-pettersson972}
Pettersson, D., E. Lörinc, T. Holm, H. Iversen, B. Cedermark, B.
Glimelius, and A. Martling. 2015. {``Tumour Regression in the Randomized
{Stockholm} {III} {Trial} of~Radiotherapy Regimens for Rectal Cancer.''}
\emph{The British Journal of Surgery} 102 (8): 972--978; discussion 978.
\url{https://doi.org/10.1002/bjs.9811}.

\bibitem[\citeproctext]{ref-phoa1209}
Phoa, K. Nadine, Frederike G. I. van Vilsteren, Bas L. A. M. Weusten,
Raf Bisschops, Erik J. Schoon, Krish Ragunath, Grant Fullarton, et al.
2014. {``Radiofrequency Ablation Vs Endoscopic Surveillance for Patients
with {Barrett} Esophagus and Low-Grade Dysplasia: A Randomized Clinical
Trial.''} \emph{JAMA} 311 (12): 1209--17.
\url{https://doi.org/10.1001/jama.2014.2511}.

\bibitem[\citeproctext]{ref-pilonis263}
Pilonis, ND, Tischkowitz M, Fitzgerald Rc, and di Pietro M. 2021.
{``Hereditary {Diffuse} {Gastric} {Cancer}: {Approaches} to {Screening},
{Surveillance}, and {Treatment}.''} \emph{Annual Review of Medicine} 72
(January): 263--80.
\url{https://doi.org/10.1146/annurev-med-051019-103216}.

\bibitem[\citeproctext]{ref-vanderpost361}
Post, Rachel S. van der, Ingrid P. Vogelaar, Fátima Carneiro, Parry
Guilford, David Huntsman, Nicoline Hoogerbrugge, Carlos Caldas, et al.
2015. {``Hereditary Diffuse Gastric Cancer: Updated Clinical Guidelines
with an Emphasis on Germline {CDH1} Mutation Carriers.''} \emph{Journal
of Medical Genetics} 52 (6): 361--74.
\url{https://doi.org/10.1136/jmedgenet-2015-103094}.

\bibitem[\citeproctext]{ref-quirke996}
Quirke, P., P. Durdey, M. F. Dixon, and N. S. Williams. 1986. {``Local
Recurrence of Rectal Adenocarcinoma Due to Inadequate Surgical
Resection. {Histopathological} Study of Lateral Tumour Spread and
Surgical Excision.''} \emph{Lancet (London, England)} 2 (8514): 996--99.
\url{https://doi.org/10.1016/s0140-6736(86)92612-7}.

\bibitem[\citeproctext]{ref-rastogi394}
Rastogi, Amit, Srinivas Puli, Hashem B. El-Serag, Ajay Bansal, Sachin
Wani, and Prateek Sharma. 2008. {``Incidence of Esophageal
Adenocarcinoma in Patients with {Barrett}'s Esophagus and High-Grade
Dysplasia: A Meta-Analysis.''} \emph{Gastrointestinal Endoscopy} 67 (3):
394--98. \url{https://doi.org/10.1016/j.gie.2007.07.019}.

\bibitem[\citeproctext]{ref-repici715}
Repici, Alessandro, Cesare Hassan, Alessandra Carlino, Nico Pagano,
Angelo Zullo, Giacomo Rando, Giuseppe Strangio, et al. 2010.
{``Endoscopic Submucosal Dissection in Patients with Early Esophageal
Squamous Cell Carcinoma: Results from a Prospective {Western} Series.''}
\emph{Gastrointestinal Endoscopy} 71 (4): 715--21.
\url{https://doi.org/10.1016/j.gie.2009.11.020}.

\bibitem[\citeproctext]{ref-reynolds401}
Reynolds, J. V., S. R. Preston, B. O'Neill, L. Baeksgaard, S. M.
Griffin, C. Mariette, S. Cuffe, et al. 2017. {``{ICORG} 10-14:
{NEOadjuvant} Trial in {Adenocarcinoma} of the {oEsophagus} and
{oesophagoGastric} Junction {International} {Study} ({Neo}-{AEGIS}).''}
\emph{BMC Cancer} 17 (1): 401.
\url{https://doi.org/10.1186/s12885-017-3386-2}.

\bibitem[\citeproctext]{ref-rice317}
Rice, T. W., D. P. Mason, S. C. Murthy, Jr. Zuccaro G., D. J. Adelstein,
L. A. Rybicki, and E. H. Blackstone. 2007. {``{T2N0M0} Esophageal
Cancer.''} \emph{J Thorac Cardiovasc Surg} 133 (2): 317--24.
\url{https://doi.org/10.1016/j.jtcvs.2006.09.023}.

\bibitem[\citeproctext]{ref-ripley226}
Ripley, R. Taylor, Inderpal S. Sarkaria, Rachel Grosser, Camelia S.
Sima, Manjit S. Bains, David R. Jones, Prasad S. Adusumilli, et al.
2016. {``Pretreatment {Dysphagia} in {Esophageal} {Cancer} {Patients}
{May} {Eliminate} the {Need} for {Staging} by {Endoscopic}
{Ultrasonography}.''} \emph{The Annals of Thoracic Surgery} 101 (1):
226--30. \url{https://doi.org/10.1016/j.athoracsur.2015.06.062}.

\bibitem[\citeproctext]{ref-roberts1325}
Roberts, Maegan E., John Michael O. Ranola, Megan L. Marshall, Lisa R.
Susswein, Sara Graceffo, Kelsey Bohnert, Ginger Tsai, Rachel T. Klein,
Kathleen S. Hruska, and Brian H. Shirts. 2019. {``Comparison of {CDH1}
{Penetrance} {Estimates} in {Clinically} {Ascertained} {Families} Vs
{Families} {Ascertained} for {Multiple} {Gastric} {Cancers}.''}
\emph{JAMA Oncology} 5 (9): 1325--31.
\url{https://doi.org/10.1001/jamaoncol.2019.1208}.

\bibitem[\citeproctext]{ref-roh5124}
Roh, Mark S., Linda H. Colangelo, Michael J. O'Connell, Greg Yothers,
Melvin Deutsch, Carmen J. Allegra, Morton S. Kahlenberg, et al. 2009.
{``Preoperative Multimodality Therapy Improves Disease-Free Survival in
Patients with Carcinoma of the Rectum: {NSABP} {R}-03.''} \emph{Journal
of Clinical Oncology: Official Journal of the American Society of
Clinical Oncology} 27 (31): 5124--30.
\url{https://doi.org/10.1200/JCO.2009.22.0467}.

\bibitem[\citeproctext]{ref-rotellini-coltvet1901}
Rotellini-Coltvet, Lisa, Alex Wallace, Gia Saini, Sailendra Naidu,
Jefferey Scott Kriegshauser, Indravadan Patel, Grace Knuttinen, Sadeer
Alzubaidi, and Rahmi Oklu. 2023. {``Percutaneous {Transesophageal}
{Gastrostomy}: {Procedural} {Technique} and {Outcomes}.''} \emph{Journal
of Vascular and Interventional Radiology: JVIR} 34 (11): 1901--7.
\url{https://doi.org/10.1016/j.jvir.2023.06.040}.

\bibitem[\citeproctext]{ref-ruppert1519}
Ruppert, R., T. Junginger, H. Ptok, J. Strassburg, C. A. Maurer, P.
Brosi, J. Sauer, et al. 2018. {``Oncological Outcome After {MRI}-Based
Selection for Neoadjuvant Chemoradiotherapy in the {OCUM} {Rectal}
{Cancer} {Trial}.''} \emph{The British Journal of Surgery} 105 (11):
1519--29. \url{https://doi.org/10.1002/bjs.10879}.

\bibitem[\citeproctext]{ref-rusch444}
Rusch, Valerie W. 2004. {``Are Cancers of the Esophagus,
Gastroesophageal Junction, and Cardia One Disease, Two, or Several?''}
\emph{Seminars in Oncology} 31 (4): 444--49.
\url{https://doi.org/10.1053/j.seminoncol.2004.04.023}.

\bibitem[\citeproctext]{ref-sarkaria764}
Sarkaria, Inderpal S., Nabil P. Rizk, Manjit S. Bains, Laura H. Tang,
David H. Ilson, Bruce I. Minsky, and Valerie W. Rusch. 2009.
{``Post-Treatment Endoscopic Biopsy Is a Poor-Predictor of Pathologic
Response in Patients Undergoing Chemoradiation Therapy for Esophageal
Cancer.''} \emph{Annals of Surgery} 249 (5): 764--67.
\url{https://doi.org/10.1097/SLA.0b013e3181a38e9e}.

\bibitem[\citeproctext]{ref-sauer173}
Sauer, R., R. Fietkau, C. Wittekind, P. Martus, C. Rödel, W.
Hohenberger, G. Jatzko, et al. 2001. {``Adjuvant Versus Neoadjuvant
Radiochemotherapy for Locally Advanced Rectal Cancer. {A} Progress
Report of a Phase-{III} Randomized Trial (Protocol
{CAO}/{ARO}/{AIO}-94).''} \emph{Strahlentherapie Und Onkologie: Organ
Der Deutschen Rontgengesellschaft ... {[}Et Al{]}} 177 (4): 173--81.
\url{https://doi.org/10.1007/pl00002396}.

\bibitem[\citeproctext]{ref-schrag322}
Schrag, Deborah, Qian Shi, Martin R. Weiser, Marc J. Gollub, Leonard B.
Saltz, Benjamin L. Musher, Joel Goldberg, et al. 2023. {``Preoperative
{Treatment} of {Locally} {Advanced} {Rectal} {Cancer}.''} \emph{The New
England Journal of Medicine} 389 (4): 322--34.
\url{https://doi.org/10.1056/NEJMoa2303269}.

\bibitem[\citeproctext]{ref-sebag-montefiore811}
Sebag-Montefiore, David, Richard J. Stephens, Robert Steele, John
Monson, Robert Grieve, Subhash Khanna, Phil Quirke, et al. 2009.
{``Preoperative Radiotherapy Versus Selective Postoperative
Chemoradiotherapy in Patients with Rectal Cancer ({MRC} {CR07} and
{NCIC}-{CTG} {C016}): A Multicentre, Randomised Trial.''} \emph{The
Lancet} 373 (9666): 811--20.
\url{https://doi.org/10.1016/S0140-6736(09)60484-0}.

\bibitem[\citeproctext]{ref-selby306}
Selby, Debbie, Amy Nolen, Cheromi Sittambalam, Karen Johansen, and Robyn
Pugash. 2019. {``Percutaneous {Transesophageal} {Gastrostomy} ({PTEG}):
{A} {Safe} and {Well}-{Tolerated} {Procedure} for {Palliation} of
{End}-{Stage} {Malignant} {Bowel} {Obstruction}.''} \emph{Journal of
Pain and Symptom Management} 58 (2): 306--10.
\url{https://doi.org/10.1016/j.jpainsymman.2019.04.031}.

\bibitem[\citeproctext]{ref-shaheen30}
Shaheen, Nicholas J., Gary W. Falk, Prasad G. Iyer, Lauren B. Gerson,
and American College of Gastroenterology. 2016. {``{ACG} {Clinical}
{Guideline}: {Diagnosis} and {Management} of {Barrett}'s {Esophagus}.''}
\emph{The American Journal of Gastroenterology} 111 (1): 30--50; quiz
51. \url{https://doi.org/10.1038/ajg.2015.322}.

\bibitem[\citeproctext]{ref-shaheen2277}
Shaheen, Nicholas J., Prateek Sharma, Bergein F. Overholt, Herbert C.
Wolfsen, Richard E. Sampliner, Kenneth K. Wang, Joseph A. Galanko, et
al. 2009. {``Radiofrequency Ablation in {Barrett}'s Esophagus with
Dysplasia.''} \emph{The New England Journal of Medicine} 360 (22):
2277--88. \url{https://doi.org/10.1056/NEJMoa0808145}.

\bibitem[\citeproctext]{ref-sharma258}
Sharma, Prateek, Richard Kozarek, and Practice Parameters Committee of
American College of Gastroenterology. 2010. {``Role of Esophageal Stents
in Benign and Malignant Diseases.''} \emph{The American Journal of
Gastroenterology} 105 (2): 258--273; quiz 274.
\url{https://doi.org/10.1038/ajg.2009.684}.

\bibitem[\citeproctext]{ref-siddharthan257}
Siddharthan, Ragavan V., Christian Lanciault, and Vassiliki Liana
Tsikitis. 2019. {``Anal Intraepithelial Neoplasia: Diagnosis, Screening,
and Treatment.''} \emph{Annals of Gastroenterology} 32 (3): 257--63.
\url{https://doi.org/10.20524/aog.2019.0364}.

\bibitem[\citeproctext]{ref-siewert260}
Siewert, J. R., H. J. Stein, and M. Feith. 2006. {``Adenocarcinoma of
the Esophago-Gastric Junction.''} \emph{Scandinavian Journal of Surgery:
SJS: Official Organ for the Finnish Surgical Society and the
Scandinavian Surgical Society} 95 (4): 260--69.
\url{https://doi.org/10.1177/145749690609500409}.

\bibitem[\citeproctext]{ref-singal402}
Singal, Ashwani Kumar, Alexander A. Dekovich, Alda L. Tam, and Michael
J. Wallace. 2010. {``Percutaneous Transesophageal Gastrostomy Tube
Placement: An Alternative to Percutaneous Endoscopic Gastrostomy in
Patients with Intra-Abdominal Metastasis.''} \emph{Gastrointestinal
Endoscopy} 71 (2): 402--6.
\url{https://doi.org/10.1016/j.gie.2009.10.037}.

\bibitem[\citeproctext]{ref-slagter877}
Slagter, Astrid E., Edwin P. M. Jansen, Hanneke W. M. van Laarhoven,
Johanna W. van Sandick, Nicole C. T. van Grieken, Karolina Sikorska,
Annemieke Cats, et al. 2018. {``{CRITICS}-{II}: A Multicentre Randomised
Phase {II} Trial of Neo-Adjuvant Chemotherapy Followed by Surgery Versus
Neo-Adjuvant Chemotherapy and Subsequent Chemoradiotherapy Followed by
Surgery Versus Neo-Adjuvant Chemoradiotherapy Followed by Surgery in
Resectable Gastric Cancer.''} \emph{BMC Cancer} 18 (1): 877.
\url{https://doi.org/10.1186/s12885-018-4770-2}.

\bibitem[\citeproctext]{ref-smalley2327}
Smalley, Stephen R., Jacqueline K. Benedetti, Daniel G. Haller, Scott A.
Hundahl, Norman C. Estes, Jaffer A. Ajani, Leonard L. Gunderson, et al.
2012. {``Updated Analysis of {SWOG}-Directed Intergroup Study 0116: A
Phase {III} Trial of Adjuvant Radiochemotherapy Versus Observation After
Curative Gastric Cancer Resection.''} \emph{Journal of Clinical
Oncology: Official Journal of the American Society of Clinical Oncology}
30 (19): 2327--33. \url{https://doi.org/10.1200/JCO.2011.36.7136}.

\bibitem[\citeproctext]{ref-smith767}
Smith, J. Joshua, Oliver S. Chow, Marc J. Gollub, Garrett M. Nash,
Larissa K. Temple, Martin R. Weiser, José G. Guillem, et al. 2015.
{``Organ {Preservation} in {Rectal} {Adenocarcinoma}: A Phase {II}
Randomized Controlled Trial Evaluating 3-Year Disease-Free Survival in
Patients with Locally Advanced Rectal Cancer Treated with Chemoradiation
Plus Induction or Consolidation Chemotherapy, and Total Mesorectal
Excision or Nonoperative Management.''} \emph{BMC Cancer} 15 (October):
767. \url{https://doi.org/10.1186/s12885-015-1632-z}.

\bibitem[\citeproctext]{ref-soetikno4490}
Soetikno, Roy, Tonya Kaltenbach, Ronald Yeh, and Takuji Gotoda. 2005.
{``Endoscopic Mucosal Resection for Early Cancers of the Upper
Gastrointestinal Tract.''} \emph{Journal of Clinical Oncology: Official
Journal of the American Society of Clinical Oncology} 23 (20): 4490--98.
\url{https://doi.org/10.1200/JCO.2005.19.935}.

\bibitem[\citeproctext]{ref-speicher1195}
Speicher, Paul J., Asvin M. Ganapathi, Brian R. Englum, Matthew G.
Hartwig, Mark W. Onaitis, Thomas A. D'Amico, and Mark F. Berry. 2014.
{``Induction Therapy Does Not Improve Survival for Clinical Stage {T2N0}
Esophageal Cancer.''} \emph{Journal of Thoracic Oncology: Official
Publication of the International Association for the Study of Lung
Cancer} 9 (8): 1195--1201.
\url{https://doi.org/10.1097/JTO.0000000000000228}.

\bibitem[\citeproctext]{ref-stahl2310}
Stahl, Michael, Martin Stuschke, Nils Lehmann, Hans-Joachim Meyer,
Martin K. Walz, Siegfried Seeber, Bodo Klump, et al. 2005.
{``Chemoradiation with and Without Surgery in Patients with Locally
Advanced Squamous Cell Carcinoma of the Esophagus.''} \emph{Journal of
Clinical Oncology: Official Journal of the American Society of Clinical
Oncology} 23 (10): 2310--17.
\url{https://doi.org/10.1200/JCO.2005.00.034}.

\bibitem[\citeproctext]{ref-stahl851}
Stahl, Michael, Martin K. Walz, Martin Stuschke, Nils Lehmann,
Hans-Joachim Meyer, Jorge Riera-Knorrenschild, Peter Langer, et al.
2009. {``Phase {III} Comparison of Preoperative Chemotherapy Compared
with Chemoradiotherapy in Patients with Locally Advanced Adenocarcinoma
of the Esophagogastric Junction.''} \emph{Journal of Clinical Oncology:
Official Journal of the American Society of Clinical Oncology} 27 (6):
851--56. \url{https://doi.org/10.1200/JCO.2008.17.0506}.

\bibitem[\citeproctext]{ref-strassburg2790}
Strassburg, Joachim, Reinhard Ruppert, Henry Ptok, Christoph Maurer,
Theodor Junginger, Susanne Merkel, and Paul Hermanek. 2011.
{``{MRI}-Based Indications for Neoadjuvant Radiochemotherapy in Rectal
Carcinoma: Interim Results of a Prospective Multicenter Observational
Study.''} \emph{Annals of Surgical Oncology} 18 (10): 2790--99.
\url{https://doi.org/10.1245/s10434-011-1704-5}.

\bibitem[\citeproctext]{ref-sudo4306}
Sudo, Kazuki, Takashi Taketa, Arlene M. Correa, Maria-Claudia Campagna,
Roopma Wadhwa, Mariela A. Blum, Ritsuko Komaki, et al. 2013.
{``Locoregional Failure Rate After Preoperative Chemoradiation of
Esophageal Adenocarcinoma and the Outcomes of Salvage Strategies.''}
\emph{Journal of Clinical Oncology: Official Journal of the American
Society of Clinical Oncology} 31 (34): 4306--10.
\url{https://doi.org/10.1200/JCO.2013.51.7250}.

\bibitem[\citeproctext]{ref-sudo3400}
Sudo, Kazuki, Lianchun Xiao, Roopma Wadhwa, Hironori Shiozaki, Elena
Elimova, Takashi Taketa, Mariela A. Blum, et al. 2014. {``Importance of
Surveillance and Success of Salvage Strategies After Definitive
Chemoradiation in Patients with Esophageal Cancer.''} \emph{Journal of
Clinical Oncology: Official Journal of the American Society of Clinical
Oncology} 32 (30): 3400--3405.
\url{https://doi.org/10.1200/JCO.2014.56.7156}.

\bibitem[\citeproctext]{ref-swedishrectalcancertrial980}
Swedish Rectal Cancer Trial, B. Cedermark, M. Dahlberg, B. Glimelius, L.
Påhlman, L. E. Rutqvist, and N. Wilking. 1997. {``Improved Survival with
Preoperative Radiotherapy in Resectable Rectal Cancer.''} \emph{The New
England Journal of Medicine} 336 (14): 980--87.
\url{https://doi.org/10.1056/NEJM199704033361402}.

\bibitem[\citeproctext]{ref-swisher175}
Swisher, Stephen G., Paula Wynn, Joe B. Putnam, Melinda B. Mosheim,
Arlene M. Correa, Ritsuko R. Komaki, Jaffer A. Ajani, et al. 2002.
{``Salvage Esophagectomy for Recurrent Tumors After Definitive
Chemotherapy and Radiotherapy.''} \emph{The Journal of Thoracic and
Cardiovascular Surgery} 123 (1): 175--83.
\url{https://doi.org/10.1067/mtc.2002.119070}.

\bibitem[\citeproctext]{ref-taketa300}
Taketa, Takashi, Arlene M. Correa, Akihiro Suzuki, Mariela A. Blum,
Pamela Chien, Jeffrey H. Lee, James Welsh, et al. 2012. {``Outcome of
Trimodality-Eligible Esophagogastric Cancer Patients Who Declined
Surgery After Preoperative Chemoradiation.''} \emph{Oncology} 83 (5):
300--304. \url{https://doi.org/10.1159/000341353}.

\bibitem[\citeproctext]{ref-taketa1139}
Taketa, Takashi, Kazuki Sudo, Arlene M. Correa, Roopma Wadhwa, Hironori
Shiozaki, Elena Elimova, Maria-Claudia Campagna, et al. 2014.
{``Post-Chemoradiation Surgical Pathology Stage Can Customize the
Surveillance Strategy in Patients with Esophageal Adenocarcinoma.''}
\emph{Journal of the National Comprehensive Cancer Network: JNCCN} 12
(8): 1139--44. \url{https://doi.org/10.6004/jnccn.2014.0111}.

\bibitem[\citeproctext]{ref-taketa95}
Taketa, Takashi, Lianchun Xiao, Kazuki Sudo, Akihiro Suzuki, Roopma
Wadhwa, Mariela A. Blum, Jeffrey H. Lee, et al. 2013.
{``Propensity-Based Matching Between Esophagogastric Cancer Patients Who
Had Surgery and Who Declined Surgery After Preoperative
Chemoradiation.''} \emph{Oncology} 85 (2): 95--99.
\url{https://doi.org/10.1159/000351999}.

\bibitem[\citeproctext]{ref-taylor34}
Taylor, Fiona G. M., Philip Quirke, Richard J. Heald, Brendan J. Moran,
Lennart Blomqvist, Ian R. Swift, David Sebag-Montefiore, Paris Tekkis,
Gina Brown, and Magnetic Resonance Imaging in Rectal Cancer European
Equivalence Study Study Group. 2014. {``Preoperative Magnetic Resonance
Imaging Assessment of Circumferential Resection Margin Predicts
Disease-Free Survival and Local Recurrence: 5-Year Follow-up Results of
the {MERCURY} Study.''} \emph{Journal of Clinical Oncology: Official
Journal of the American Society of Clinical Oncology} 32 (1): 34--43.
\url{https://doi.org/10.1200/JCO.2012.45.3258}.

\bibitem[\citeproctext]{ref-taylor711}
Taylor, Fiona G. M., Philip Quirke, Richard J. Heald, Brendan Moran,
Lennart Blomqvist, Ian Swift, David J. Sebag-Montefiore, Paris Tekkis,
Gina Brown, and MERCURY study group. 2011. {``Preoperative
High-Resolution Magnetic Resonance Imaging Can Identify Good Prognosis
Stage {I}, {II}, and {III} Rectal Cancer Best Managed by Surgery Alone:
A Prospective, Multicenter, {European} Study.''} \emph{Annals of
Surgery} 253 (4): 711--19.
\url{https://doi.org/10.1097/SLA.0b013e31820b8d52}.

\bibitem[\citeproctext]{ref-terheggen783}
Terheggen, Grischa, Eva Maria Horn, Michael Vieth, Helmut Gabbert,
Markus Enderle, Alexander Neugebauer, Brigitte Schumacher, and Horst
Neuhaus. 2017. {``A Randomised Trial of Endoscopic Submucosal Dissection
Versus Endoscopic Mucosal Resection for Early {Barrett}'s Neoplasia.''}
\emph{Gut} 66 (5): 783--93.
\url{https://doi.org/10.1136/gutjnl-2015-310126}.

\bibitem[\citeproctext]{ref-udomsawaengsup2314}
Udomsawaengsup, S., S. Brethauer, M. Kroh, and B. Chand. 2008.
{``Percutaneous Transesophageal Gastrostomy ({PTEG}): A Safe and
Effective Technique for Gastrointestinal Decompression in Malignant
Obstruction and Massive Ascites.''} \emph{Surgical Endoscopy} 22 (10):
2314--18. \url{https://doi.org/10.1007/s00464-008-9984-y}.

\bibitem[\citeproctext]{ref-vakil1791}
Vakil, N., A. I. Morris, N. Marcon, A. Segalin, A. Peracchia, N. Bethge,
G. Zuccaro, J. J. Bosco, and W. F. Jones. 2001. {``A Prospective,
Randomized, Controlled Trial of Covered Expandable Metal Stents in the
Palliation of Malignant Esophageal Obstruction at the Gastroesophageal
Junction.''} \emph{The American Journal of Gastroenterology} 96 (6):
1791--96. \url{https://doi.org/10.1111/j.1572-0241.2001.03923.x}.

\bibitem[\citeproctext]{ref-vandervalk75}
Valk, Maxime J. M. van der, Corrie A. M. Marijnen, Boudewijn van Etten,
Esmée A. Dijkstra, Denise E. Hilling, Elma Meershoek-Klein Kranenbarg,
Hein Putter, et al. 2020. {``Compliance and Tolerability of Short-Course
Radiotherapy Followed by Preoperative Chemotherapy and Surgery for
High-Risk Rectal Cancer - {Results} of the International Randomized
{RAPIDO}-Trial.''} \emph{Radiotherapy and Oncology: Journal of the
European Society for Therapeutic Radiology and Oncology} 147 (June):
75--83. \url{https://doi.org/10.1016/j.radonc.2020.03.011}.

\bibitem[\citeproctext]{ref-visbal1803}
Visbal, A. L., M. S. Allen, D. L. Miller, C. Deschamps, V. F. Trastek,
and P. C. Pairolero. 2001. {``Ivor {Lewis} Esophagogastrectomy for
Esophageal Cancer.''} \emph{The Annals of Thoracic Surgery} 71 (6):
1803--8. \url{https://doi.org/10.1016/s0003-4975(01)02601-7}.

\bibitem[\citeproctext]{ref-vanvliet547}
Vliet, E. P. M. van, M. H. Heijenbrok-Kal, M. G. M. Hunink, E. J.
Kuipers, and P. D. Siersema. 2008. {``Staging Investigations for
Oesophageal Cancer: A Meta-Analysis.''} \emph{British Journal of Cancer}
98 (3): 547--57. \url{https://doi.org/10.1038/sj.bjc.6604200}.

\bibitem[\citeproctext]{ref-wang20}
Wang, Xinxin, Shuo Li, Yihong Sun, Kai Li, Xian Shen, Yingwei Xue, Pin
Liang, et al. 2021. {``The Protocol of a Prospective, Multicenter,
Randomized, Controlled Phase {III} Study Evaluating Different Cycles of
Oxaliplatin Combined with {S}-1 ({SOX}) as Neoadjuvant Chemotherapy for
Patients with Locally Advanced Gastric Cancer: {RESONANCE}-{II}
Trial.''} \emph{BMC Cancer} 21 (1): 20.
\url{https://doi.org/10.1186/s12885-020-07764-7}.

\bibitem[\citeproctext]{ref-weaver1301}
Weaver, Joshua, Priya Rao, John R. Goldblum, Michael J. Joyce, Sondra L.
Turner, Alexander J. F. Lazar, Dolores López-Terada, Raymond R. Tubbs,
and Brian P. Rubin. 2010. {``Can {MDM2} Analytical Tests Performed on
Core Needle Biopsy Be Relied Upon to Diagnose Well-Differentiated
Liposarcoma?''} \emph{Modern Pathology: An Official Journal of the
United States and Canadian Academy of Pathology, Inc} 23 (10): 1301--6.
\url{https://doi.org/10.1038/modpathol.2010.106}.

\bibitem[\citeproctext]{ref-wei1496}
Wei, Feixue, Catharina J. Alberts, Andreia Albuquerque, and Gary M.
Clifford. 2023. {``Impact of {Human} {Papillomavirus} {Vaccine}
{Against} {Anal} {Human} {Papillomavirus} {Infection}, {Anal}
{Intraepithelial} {Neoplasia}, and {Recurrence} of {Anal}
{Intraepithelial} {Neoplasia}: {A} {Systematic} {Review} and
{Meta}-Analysis.''} \emph{The Journal of Infectious Diseases} 228 (11):
1496--1504. \url{https://doi.org/10.1093/infdis/jiad183}.

\bibitem[\citeproctext]{ref-vanworkum601}
Workum, Frans van, Moniek H. P. Verstegen, Bastiaan R. Klarenbeek,
Stefan A. W. Bouwense, Mark I. van Berge Henegouwen, Freek Daams,
Suzanne S. Gisbertz, et al. 2021. {``Intrathoracic Vs {Cervical}
{Anastomosis} {After} {Totally} or {Hybrid} {Minimally} {Invasive}
{Esophagectomy} for {Esophageal} {Cancer}: {A} {Randomized} {Clinical}
{Trial}.''} \emph{JAMA Surgery} 156 (7): 601--10.
\url{https://doi.org/10.1001/jamasurg.2021.1555}.

\bibitem[\citeproctext]{ref-worrell484}
Worrell, Stephanie G., Evan T. Alicuben, Daniel S. Oh, Jeffrey A. Hagen,
and Steven R. DeMeester. 2018. {``Accuracy of {Clinical} {Staging} and
{Outcome} {With} {Primary} {Resection} for {Local}-{Regionally}
{Limited} {Esophageal} {Adenocarcinoma}.''} \emph{Annals of Surgery} 267
(3): 484--88. \url{https://doi.org/10.1097/SLA.0000000000002139}.

\bibitem[\citeproctext]{ref-wouters1789}
Wouters, M. W. J. M., H. E. Karim-Kos, S. le Cessie, B. P. L. Wijnhoven,
L. P. S. Stassen, W. H. Steup, H. W. Tilanus, and R. a. E. M. Tollenaar.
2009. {``Centralization of Esophageal Cancer Surgery: Does It Improve
Clinical Outcome?''} \emph{Annals of Surgical Oncology} 16 (7):
1789--98. \url{https://doi.org/10.1245/s10434-009-0458-9}.

\bibitem[\citeproctext]{ref-xicola838}
Xicola, Rosa M., Shuwei Li, Nicolette Rodriguez, Patrick Reinecke,
Rachid Karam, Virginia Speare, Mary Helen Black, Holly LaDuca, and
Xavier Llor. 2019. {``Clinical Features and Cancer Risk in Families with
Pathogenic {CDH1} Variants Irrespective of Clinical Criteria.''}
\emph{Journal of Medical Genetics} 56 (12): 838--43.
\url{https://doi.org/10.1136/jmedgenet-2019-105991}.

\bibitem[\citeproctext]{ref-ychou1715}
Ychou, Marc, Valérie Boige, Jean-Pierre Pignon, Thierry Conroy, Olivier
Bouché, Gilles Lebreton, Muriel Ducourtieux, et al. 2011.
{``Perioperative Chemotherapy Compared with Surgery Alone for Resectable
Gastroesophageal Adenocarcinoma: An {FNCLCC} and {FFCD} Multicenter
Phase {III} Trial.''} \emph{Journal of Clinical Oncology: Official
Journal of the American Society of Clinical Oncology} 29 (13): 1715--21.
\url{https://doi.org/10.1200/JCO.2010.33.0597}.

\bibitem[\citeproctext]{ref-young629}
Young, Anne N., Elizabeth Jacob, Patrick Willauer, Levi Smucker, Raul
Monzon, and Luis Oceguera. 2020. {``Anal {Cancer}.''} \emph{The Surgical
Clinics of North America} 100 (3): 629--34.
\url{https://doi.org/10.1016/j.suc.2020.02.007}.

\bibitem[\citeproctext]{ref-zhang429}
Zhang, J. Q., C. M. Hooker, M. V. Brock, J. Shin, S. Lee, R. How, N.
Franco, H. Prevas, A. Hulbert, and S. C. Yang. 2012. {``Neoadjuvant
Chemoradiation Therapy Is Beneficial for Clinical Stage {T2} {N0}
Esophageal Cancer Patients Due to Inaccurate Preoperative Staging.''}
\emph{Ann Thorac Surg} 93 (2): 429-35; discussion 436-7.
\url{https://doi.org/10.1016/j.athoracsur.2011.10.061}.

\bibitem[\citeproctext]{ref-zhoue0132889}
Zhou, Can, Li Zhang, Hua Wang, Xiaoxia Ma, Bohui Shi, Wuke Chen, Jianjun
He, Ke Wang, Peijun Liu, and Yu Ren. 2015. {``Superiority of {Minimally}
{Invasive} {Oesophagectomy} in {Reducing} {In}-{Hospital} {Mortality} of
{Patients} with {Resectable} {Oesophageal} {Cancer}: {A}
{Meta}-{Analysis}.''} \emph{PloS One} 10 (7): e0132889.
\url{https://doi.org/10.1371/journal.pone.0132889}.

\bibitem[\citeproctext]{ref-zhu262}
Zhu, Clara, Rebecca Platoff, Gaby Ghobrial, Jackson Saddemi, Taylor
Evangelisti, Emily Bucher, Benjamin Saracco, et al. 2022. {``What to Do
{When} {Decompressive} {Gastrostomies} and {Jejunostomies} Are Not
{Options}? {A} {Scoping} {Review} of {Transesophageal} {Gastrostomy}
{Tubes} for {Advanced} {Malignancies}.''} \emph{Annals of Surgical
Oncology} 29 (1): 262--71.
\url{https://doi.org/10.1245/s10434-021-10667-x}.

\bibitem[\citeproctext]{ref-vanderzijden7759}
Zijden, Charlène J. van der, Pieter C. van der Sluis, Bianca Mostert,
Joost J. M. E. Nuyttens, J. Jan B. van Lanschot, Manon C. W. Spaander,
Roelf Valkema, et al. 2024. {``Interval {Metastases} {After}
{Neoadjuvant} {Chemoradiotherapy} for {Patients} with {Locally}
{Advanced} {Esophageal} {Cancer}: {A} {Multicenter} {Observational}
{Cohort} {Study}.''} \emph{Annals of Surgical Oncology} 31 (12):
7759--66. \url{https://doi.org/10.1245/s10434-024-15890-w}.

\end{CSLReferences}


\backmatter


\end{document}
